\documentclass[a4paper, german, bibliography=totoc, oneside, BCOR=3mm]{scrreprt}
\KOMAoptions{parskip=half}

% PACKAGES
% \usepackage{mathastext} % uses text font; bad arrows -> NO
% \usepackage{euler} % for egyptienne
% \usepackage{sfmath}\usepackage{sansmathaccent} % sans serif
% \usepackage[math]{anttor} % letters look slightly crippled in comparison
% \usepackage[math]{kurier}
% \usepackage{mathpazo}
%\usepackage[charter]{mathdesign} % not working :-(

\usepackage{fontspec}
\setmainfont{Latin Modern Roman}
\setsansfont{Latin Modern Sans}
%\setmainfont{Charis SIL}
%\setsansfont{Source Sans Pro}
%\setsansfont{Fira Sans}
%\setsansfont{Bitstream Vera Sans} % too wide

\usepackage{babel}
\usepackage[autostyle=try]{csquotes}

\usepackage[style=alphabetic, backend=biber]{biblatex}
\bibliography{neron_models.bib}

\usepackage{mathtools}
\usepackage{amssymb}
\usepackage{amsfonts}
\usepackage{amsthm}
\usepackage{dsfont}
\usepackage{tikz-cd}
\usetikzlibrary{babel}
\hyphenation{Né-ron=Ab-bil-dungs-ei-gen-schaft Grup-pen-sche-ma-ta
  Treu-flach-heit wo-raus}

\usepackage{enumitem}
\usepackage{booktabs}

\usepackage{hyperref}
\usepackage{tabularx}

% SCHREIBHILFEN
\newcommand*{\optcite}[2][]{\cite[#1]{#2}}

% SONSTIGE DEFINITIONEN
\newcommand*{\NAbbEig}{Néron-Abbildungseigenschaft }

% MATHE DEFINITIONEN
\newtheorem{Satz}{Satz}[chapter]
\newtheorem{Theorem}[Satz]{Theorem}
\newtheorem{Proposition}[Satz]{Proposition}
\newtheorem{Lemma}[Satz]{Lemma}
\newtheorem{Korollar}[Satz]{Korollar}
\theoremstyle{definition}
\newtheorem{Definition}[Satz]{Definition}
\newtheorem{LemmaDefinition}[Satz]{Lemma/Definition}
\newtheorem{Beispiel}[Satz]{Beispiel}
%\theoremstyle{remark}
\newtheorem{Bemerkung}[Satz]{Bemerkung}

\renewcommand*{\epsilon}{\varepsilon} % richtiges Epsilon
\renewcommand*{\c}{\alpha} % Konstanten
\newcommand*{\N}{\mathds{N}} % natürliche Zahlen
\newcommand*{\Z}{\mathds{Z}} % ganze Zahlen
\newcommand*{\Q}{\mathds{Q}} % rationale Zahlen
\renewcommand*{\G}{\mathds{G}} % additives/multiplikatives Gruppenschema
\newcommand*{\K}{\ensuremath{K}} % Körper
\newcommand*{\algK}{\ensuremath{\overline K}} % algebraischer
\newcommand*{\Rh}{\ensuremath{R^{\text{h}}}} % Henselisierung
\renewcommand*{\Rsh}{\ensuremath{R^{\text{sh}}}} % (strikte) Henselisierung
\newcommand*{\Ksh}{\ensuremath{\K^{\text{sh}}}} % (strikte) Henselisierung
\newcommand*{\p}{\ensuremath{\mathfrak{p}}} % Primideal
\renewcommand*{\P}{\ensuremath{\mathfrak{P}}} % Primideal
% Abschluss
\renewcommand*{\Im}{\text{im}} % Bild
\newcommand*{\proj}{\text{proj}} % Projektion
\newcommand*{\dra}{\dashrightarrow} % rationale Abbildung
\DeclareMathOperator{\Dom}{Dom} % Definitionsbereich einer rat. Abb.
\DeclareMathOperator{\Quot}{Quot} % Quotientenkörperoperator
\DeclareMathOperator{\Nil}{Nil} % Nilradikal
\DeclareMathOperator{\Spec}{Spec} % Spektrum
\newcommand*{\A}{\ensuremath{\mathds{A}}} % affiner Raum
\renewcommand*{\P}{\ensuremath{\mathds{P}}} % projektiver Raum
\renewcommand*{\m}{\mathfrak{m}} % Maximalideal
\renewcommand*{\k}{\kappa} % Restklassenkörperoperator
\newcommand*{\n}{\eta} % generischer Punkt
\renewcommand*{\O}{\mathcal{O}} % Garbe
\newcommand*{\I}{\mathcal{I}} % Idealgarbe
\newcommand*{\Om}[1]{\ensuremath{\Omega_{#1}^1}} % relative 1-Formen
\renewcommand*{\d}{\ensuremath{\operatorname{d}}} % (universelles) Differential
\DeclareMathOperator{\rk}{rk} % Rang (einer Matrix)
\newcommand*{\longto}{\longrightarrow}
\newcommand*{\longfrom}{\longleftarrow}
%\newcommand*{\Xn}{\underline{X}} % Variablenmenge
\newcommand*{\degs}{\operatorname{\deg}_s} % Separabilitätsgrad
\newcommand*{\degi}{\operatorname{\deg}_i} % Inseparabilitätsgrad
\DeclareMathOperator{\ord}{ord} % diskrete Bewertung
\DeclareMathOperator{\Div}{Div} % Divisorengruppe
\DeclareMathOperator{\Pic}{Pic} % Picard Gruppe
\renewcommand{\div}{\operatorname{div}}
\DeclareMathOperator{\Char}{char} % Charakteristik
\DeclareMathOperator{\Hom}{Hom} % Homomorphismenmenge
\DeclareMathOperator{\End}{End} % Endomorphismenmenge
\DeclareMathOperator{\Aut}{Aut} % Automorphismenmenge
\newcommand{\Id}{\mathrm{id}} % Identitätsabbildung
\DeclareMathOperator{\im}{im} % image
%\newcommand{\tlambda}{\tilde\lambda} % alternative Abbildung \lambda
%\newcommand{\F}{F} % endliche Untergruppe
\DeclareMathOperator{\genus}{genus} % Kurvengeschlecht
\newcommand*{\Groups}{\text{(Groups)}} % Kategorie der Gruppen
\newcommand*{\Sch}[1][S]{\text{(Sch/#1)}} % Kategorie der #1-Schemata


%%% Local Variables:
%%% mode: latex
%%% TeX-master: "neron_models"
%%% End:


% TITEL
\hypersetup{
  pdfauthor={Gesina Schwalbe},
  pdftitle={Neron Modelle}
}
\newlength{\urlogowidth}\setlength{\urlogowidth}{11.5cm}
\titlehead{%
  \parbox{2\urlogowidth}{%
    \includegraphics[width=.98\urlogowidth]{UniLogo_Text.jpg}\\%
    \hspace*{0.3333\urlogowidth}%
    {\sffamily\bfseries\scshape%
      \fontsize{0.041667\urlogowidth}{0.041667\urlogowidth}\selectfont%
      Fakultät für Mathematik}\\[2em]
    % \hspace*{0.3333\urlogowidth}{\Large Betreuer: Prof.\,Dr.~Moritz Kerz}
  }%
}

\subject{Bachelorarbeit}
\title{Néron-Modelle elliptischer Kurven}
\author{Gesina Schwalbe}
\publishers{Betreuer: Prof.\,Dr.~Moritz Kerz}
\newcommand{\abgabedatum}{28.~September~2016}
\date{\abgabedatum}

% VORVERSION:
%\includeonly{}

\begin{document}

\maketitle
\tableofcontents


% Motivation
\chapter{Motivation}

Sei $(R,\m)$ diskreter Bewertungsring, $\K=\Quot(R)$ und $E_\K$
elliptische Kurve über $\K$. Eine Weierstraßgleichung zu $E_\K$
\begin{gather*}
  0 = Y^2 Z + \c_1 XYZ + \c_3 YZ^2 - X^3 - \c_2X^2 Z - \c_4 XZ^2 - \c_6 Z^3
  \eqqcolon F\in\K[X,Y,Z]  
\end{gather*}
mit Koeffizienten in $R$ definiert ein projektives $R$-Schema
$W=V_+(F)\subset\P_R^n$ der Dimension~2. Dessen generische Faser
ist wieder $E_\K$ und die spezielle Faser ist die Kurve, die durch das
reduzierte Polynom $(F\mod\m)$ über $\k(\m)$ erzeugt wird.
$W$ enthält also Informationen zur Reduktion unserer elliptischen
Kurve.

Da elliptische Kurven besonders durch ihre Gruppenstruktur
interessant sind, stellt sich die Frage, ob sich die
Gruppenstruktur $E_\K\times_\K E_\K\to E_\K$ auf ein hinreichend
gutartiges (hier: glattes, separiertes) $R$-Schema wie $W$ erweitern
lässt.
Eine Forderung an ein solches $R$-Schema ist dann:
\begin{quote}
  \itshape
  Die Gruppenverknüpfung auf der generischen Faser erweitert sich zu
  einer Gruppenverknüpfung $E\times_R E\to E$ auf ganz $E$.
\end{quote}
Insbesondere sind dann für $R$-Schemata $Y$ alle Schnitte $E(Y)$
Gruppen. Per Konstruktion ist die kanonische Abbildung
$E(Y)\to E_\K(Y_\K)$, die durch den Basiswechsel induziert ist, ein
Gruppenmorphismus, der allerdings im Allgemeinen weder surjektiv noch
injektiv ist. Für eine \enquote{schöne} Erweiterung der
Gruppenstruktur fordern also die stärkere Eigenschaft:
\begin{quote}
  \itshape
  Für alle \enquote{gutartigen} (hier: glatten) $R$-Schemata $Y$ mit
  generischer Faser $Y_\K$ soll die kanonische Abbildung
  $E(Y)\to E_\K(Y_\K)$ ein (Gruppen-)Isomorphismus sein.
\end{quote}
Dies motiviert die Definition des Néron-Modells, welches
ein $R$-Modell von $E_\K$ ist, das ebendiese Erweiterungseigenschaften
erfüllt.

Das Hauptziel dieser Arbeit wird sein, die Existenz von Néron-Modellen
für elliptische Kurven zu zeigen und die Konstruktion zu skizzieren.
Wir beschränken uns auf Dedekindringe $R$ mit perfekten
Restklassenkörpern, für allgemeine Dedekindringe siehe auch
\cite{nonperfect}.

In \autoref{chap:gruppenstrukturenaufschemata} werden die nötigen
Grundlagen zu Gruppenschemata eingeführt, wonach dann in
\autoref{chap:neronmodelle} Néron-Modelle allgemein definiert und
einige ihrer Eigenschaften behandelt werden.
Hier wird auch ein Blick darauf geworfen, wann bzw. welcher Teil eines
Néron-Modells einer abelschen Gruppenstruktur wieder abelsch ist.

In \autoref{chap:weierstraßmodelle} wird die Konstruktion von $W$ von
oben genauer erläutert und die Eigenschaften einer solchen Kurve über
$R$ untersucht. Insbesondere werden wir feststellen, dass der glatte
Teil von $W$ tatsächlich die Gruppenverknüpfung erweitert und im
glatten (bzw. später auch regulären) Fall sogar das Néron-Modell von
$E_\K$ ist.

Ist $W$ nicht regulär, müssen erst weitere Modifikationen an $W$
vorgenommen werden, um eine hinreichend gutartige, insb. reguläre,
Kurve über $R$ zu erhalten. Solche gutartigen Kurven über
einem Dedekindring werden arithmetische Fläche genannt und in
\autoref{chap:arithmetischeflächen} näher auf ihre Eigenschaften und
die Konstruktion für elliptische Kurven untersucht.

Zuletzt wird sich in \autoref{chap:exneronmodelle} zeigen, dass der
glatte Teil einer minimal gewählten arithmetischen Fläche zu $E_\K$
bereits das Néron-Modell von $E_\K$ ist, womit die Konstruktion und
insbesondere der Existenzbeweis abgeschlossen ist.

\autoref{chap:ausblick} enthält dann noch ein paar Beispiele
bzw. Anwendungen und den Ausblick auf eine Klassifizierung
elliptischer Kurven nach der Form ihres minimalen Modells.

\section*{Notation}
Es werden Kenntnisse zur Garben-~und Schematheorie benötigt,
insbesondere zu den Begriffen glatt, flach, eigentlich, separiert.
Die Notation richtet sich nach \cite{wedhorn}.

% Mit affiner Varietät ist stets ein affines, integeres $\K$-Schema
% endlichen Typs gemeint.
Eine Kurve über einem Körper $\K$ sei immer eine projektive Varietät
der Dimension~1,
und eine elliptische Kurve über $\K$ ist eine glatte Kurve $E_\K$
vom Genus~1 zusammen mit einem $\K$-Schnitt $O\in E_\K(\K)$.
Elliptische Kurven sind $\K$"=Gruppenschemata mit Nullelement $O$
\cite{silverman}.

Es werden nur elliptische Kurven über Körpern der Charakteristik
ungleich~2, 3 betrachten und die Restklassenkörper der Dedekindringe
seien perfekt.

Im Folgenden noch einige weitere verwendete Begriffe:

% DEDEKINDRINGE
\begin{Definition}[Dedekindring]\label{def:dedekindring}
  \optcite[S.\,40]{hartshorne}\optcite[Definition B.84]{wedhorn}
  Ein Integritätsring $R$ heißt Dedekindring, wenn er eine der
  äquivalenten Bedingungen erfüllt
  \begin{enumerate}[label=(\roman*)]
  \item $R$ ist regulär und $\dim(R)\leq 1$.
  \item $R$ ist noethersch, normal und $\dim(R)\leq 1$.
  \item $R$ ist Körper oder
    jede Lokalisierung nach einem Maximalideal ist ein diskreter
    Bewertungsring
    (d.\,h. lokal, noethersch, normal von Dimension~1).
  \end{enumerate}
\end{Definition}
Für Schemata über Dedekindringen werden die Fasern über dem
generischen Punkt als generische Fasern, die über den geschlossenen
Punkten als spezielle Fasern bezeichnet.
\begin{Definition}[Dedekindschema]
  Ein Dedekindschema ist ein noethersches, integeres Schema $X$, das
  eine der äquivalenten Eigenschaften erfüllt
  \begin{enumerate}[label=(\roman*)]
  \item Die offenen affinen Unterschemata sind Spektren von Dedekindringen.
  \item $\dim(X)\leq 1$ und $X$ ist regulär
    (in Dimension 1 äquivalent zu normal nach
    \cite[Corollary 6.39, Proposition 6.40]{wedhorn}).
  \end{enumerate}
\end{Definition}
Sei $S$ allgemein ein irreduzibles Schema. Ein $S$-Schema wird auch als
$S$-Modell seiner generischen Faser bezeichnet.

Folgende Äquivalenz und Aussage zu Glattheit werden häufig verwendet:
\begin{Bemerkung}\label{thm:eigglatt}
  Sei $f\colon X\to Y$ Morphismus von Schemata.
  Für $f$ endlich präsentiert ist $f\colon X\to Y $ glatt in $x$
  äquivalent zu $f$ flach in $x$ und $X_{f(x)}$ glatt über $\k(f(x))$ nach
  \cite[8.5, Proposition 17]{bosch}.
  
  Ist $X$ ein $\K$-Schema lokal endlichen Typs sind glatt und
  geometrisch regulär äquivalent nach \cite[Corollary 6.32]{wedhorn}.
  Ist $\K$ perfekter Körper (z.\,B. $\Char(\K)=0$), so ist dies sogar
  äquivalent zu regulär nach \cite[Remark 6.33]{wedhorn}.
    
  Glatt impliziert reduziert, lokal von endlichem Typ und flach
  \cite[für flach s.][Theorem 14.22]{wedhorn}.
\end{Bemerkung}

%%% Local Variables:
%%% mode: latex
%%% TeX-master: "neron_models"
%%% End:


% Gruppenstrukturen auf Schemata
\chapter{Gruppenstrukturen auf Schemata}
\label{chap:gruppenstrukturenaufschemata}
\section{Gruppenschemata}
Im Folgenden sei $S$ ein Schema.
Zuerst werden wir in diesem Kapitel Gruppenschemata einführen, also
$S$-Schemata, deren Schnittfunktor Gruppen liefert. Als Spezialfälle
werden abelsche Schemata und abelsche Varietäten aufgeführt. Ein
wichtiges Beispiel für letzteres sind gerade elliptische Kurven.

Danach werden wir sehen, dass sich eigentliche Gruppenschemata
gutartig bzgl. rationalen Abbildungen verhalten,
d.\,h. bzgl. Morphismen, die nur auf offenen dichten Teilmengen
definiert sind.
\begin{Definition}[$S$"=Gruppenschema]\label{def:gruppenschema}
  Ein \emph{Gruppenschema über $S$} oder \emph{$S$-Gruppenschema} ist ein
  $S$-Schema ${\pi\colon G\to S}$ zusammen mit $S$-Morphismen
  \begin{align*}
    \mu\colon G\times_S G &\longto G
    &&\text{(Gruppenverknüpfung)}\\
    \epsilon\colon S&\longto G 
    &&\text{(Neutrales Element)}\\
    i\colon G &\longto G    
    &&\text{(Inverses)}
  \end{align*}
  für die folgende Diagramme kommutieren

  \begin{tabular}{>{\itshape(}l<{)}l}
    Assoziativität
    &\begin{tikzcd}
      G\times_S G\times_S G 
      \arrow[r, "{(\mu,\Id)}"] \arrow[d, "{(\Id,\mu)}"]
      & G\times_S G \arrow[d, "{\mu}"] \\
      G\times_S G \arrow[r, "\mu"]
      & G
    \end{tikzcd}\\
    Neutrales Element
    &\begin{tikzcd}
      G \arrow[r, "{(\pi,\Id)}"] \arrow[d, "{\Id}"]
      & S\times_S G \arrow[d, "{\epsilon\times\Id}"] \\
      G \arrow[from=r, "\mu"]
      & G\times_S G
    \end{tikzcd}
    und
    \begin{tikzcd}
      G \arrow[r, "{(\Id,\pi)}"] \arrow[d, "\Id"]
      & G\times_S S \arrow[d, "{\epsilon\times\Id}"] \\
      G \arrow[from=r, "\mu"]
      & G\times_S G
    \end{tikzcd}\\
    Inverses
    &\begin{tikzcd}
      G \arrow[r, "{(i,\Id)}"] \arrow[d, "{\pi}"]
      & G\times_S G \arrow[d, "{\mu}"] \\
      S \arrow[r, "\epsilon"]
      & G
    \end{tikzcd}
    und
    \begin{tikzcd}
      G \arrow[r, "{(\Id,i)}"] \arrow[d, "{\pi}"]
      & G\times_S G \arrow[d, "{\mu}"] \\
      S \arrow[r, "\epsilon"]
      & G
    \end{tikzcd}
  \end{tabular}
\end{Definition}
Die Folgende Äquivalenz zeigt, wo der Gruppenbegriff eines
Gruppenschemas zum tragen kommt:
\begin{Bemerkung}
  Die obige Definition ist äquivalent dazu, dass die $T$-wertigen Punkte
  $G(T)=\Hom_S(T,G)$ eine Gruppenstruktur tragen, die funktoriell in
  $T$ ist für $S$-Schemata $T$.\optcite[Chapter (4.15)]{wedhorn}
  In anderen Worten der Schnittfunktor faktorisiert über den
  Vergissfunktor
  \begin{align*}
    \Sch &\longto \Groups \longto \text{(Sets)}\\
    T &\longmapsto G(T)
  \end{align*}
  Die Gruppenstruktur auf $G(T)$ ist durch $\mu$ gegeben als
  \begin{align*}
    G(T)\times G(T) &\longto G(T)\\
    (\phi, \psi) &\longmapsto 
                   \left(
                   T\xrightarrow{(\phi,\psi)} G\times_S G
                   \overset{\mu}\longto G
                   \right)
  \end{align*}
  wobei der Schnitt $T\to S\xrightarrow{\epsilon} G$ die
  Identität ist und man Inverse durch Verknüpfung mit $i$ erhält.
  \optcite[Proposition IV.3.2]{silverman2}
\end{Bemerkung}

Anhand der Gruppenstruktur auf den Schnitten lassen sich
Morphismen zwischen Gruppenschemata und Untergruppenschemata
definieren.
\begin{Definition}[Morphismen von Gruppenschemata]
  \optcite[Definition 4.42]{wedhorn}
  Seien $G$ und $H$ $S$"=Gruppenschemata.
  Ein Morphismus von $S$"=Gruppenschemata ist ein Morphismus $f\colon
  G\to H$ von $S$-Schemata, für den die induzierten Abbildungen
  \begin{gather*}
    (\bullet\circ f)\colon G(T)\to H(T)
  \end{gather*}auf den $T$-wertigen Punkten
  Gruppenhomomorphismen sind für alle $S$-Schemata $T$.
  In anderen Worten
  \begin{gather*}
    f\circ\mu_G = \mu_H\circ(f\times_S f)
  \end{gather*}
\end{Definition}

% IDENTITÄTSKOMPONENTE
\begin{Definition}[Untergruppenschema]
  \optcite[Definition 4.45]{wedhorn}
  Sei $G$ ein $S$"=Gruppenschema.
  Ein $S$"=Untergruppenschema von $G$ ist ein geschlossenes
  $S$-Unterschema von $G$, dessen geschlossene Immersion ein
  Homomorphismus von Gruppenschemata ist.
\end{Definition}

Hier gleich ein wichtiges Beispiel eines Untergruppenschemas, das
durch die Topologie gegeben ist.
\begin{Definition}[Identitätskomponente]
  \optcite[Remark IV.6.1.2]{silverman2}
  Sei $G$ ein $S$"=Gruppenschema.
  Für $s\in S$ ist die Identitätskomponente einer Faser $G_s$ von $G$
  die Zusammenhangskomponente von $G_s$, die $\epsilon(s)$ enthält.
  Die Identitäts-~oder Zusammenhangskomponente von $G$ ist
  die Vereinigung aller Identitätskomponenten der Fasern von $G$.
  Die Identitätskomponente ist ein Untergruppenschema.
\end{Definition}

% ABELSCHE SCHEMATA
Wir werden nun einige wichtige Spezialfälle von Gruppenschemata
betrachten.
\begin{Bemerkung}\label{gruppenschemaaequivalenzen}
  Für ein $\K$"=Gruppenschema endlichen Typs über einem Körper $\K$
  sind glatt und geometrisch reduziert äquivalent.
  \optcite[8.5, Excercise 11]{bosch}
  Im glatten Fall sind auch irreduzibel und zusammenhängend äquivalent
  nach \cite[Exercise 3.16]{wedhorn}, \cite[Corollary 16.51]{wedhorn}.
  % , da
  % aufgrund der Glattheit alle Halme lokal und regulär, also
  % insb. integer, sind und $G$ lokal noethersch ist.
\end{Bemerkung}
\begin{Definition}[Abelsche Varietät]\label{def:abelschevarietaet}
  Eine abelsche Varietät $X/\K$ ist ein eigentliches, glattes,
  irreduzibles $\K$"=Gruppenschema.
  \optcite[9.6, Definition 1]{bosch}\optcite[Definition 16.52]{wedhorn}
\end{Definition}
\begin{Bemerkung}\label{thm:abvarietaetprojektiv}
  Abelsche Varietäten sind projektiv nach
  \cite[9.6, Proposition 4]{bosch}.
\end{Bemerkung}

\begin{Definition}[Abelsches Schema]
  Ein abelsches $S$-Schema ist ein eigentliches, glattes $S$"=Gruppenschema,
  dessen Fasern abelsche Varietäten sind 
  (also zusätzlich irreduzibel bzw. zusammenhängend,
  s. \ref{def:abelschevarietaet} und \ref{gruppenschemaaequivalenzen},
  nachdem eigentlich und glatt stabil unter Basiswechsel sind).
\end{Definition}

% S-RATIONALE ABBILDUNG
In manchen Fällen ist für ($S$-)Schemata $X$, $Y$ kein
($S$-)Morphismus $f\colon X\to Y$ gegeben, sondern nur ein Morphismus
auf einer offenen dichten Teilmenge von $X$ (z.\,B. der generischen
Faser). Solche Fälle sind besonders interessant, wenn man Aussagen
darüber treffen möchte, ob und wie sich dieser Morphismus auf ganz $X$
erweitert. Ein Beispiel sind rationale Abbildungen zwischen
projektiven Varietäten über algebraisch abgeschlossenen Körpern.
Eine Verallgemeinerung davon sind die ($S$-)rationalen Abbildungen
zwischen ($S$-)Schemata.
\begin{Definition}[rationale Abbildung]\label{def:ratabb}
  \optcite[Chapter 2.5]{neron}
  Sei $S$ Schema, und seien $X$ und $Y$ glatte $S$-Schemata.
  Eine rationale Abbildung $\phi\colon X\dra Y$ von
  $S$-Schemata ist eine Äquivalenzklasse der Menge
  \begin{gather*}
    \left\{
      \psi\colon U\to Y \text{ $S$-Morphismus}
      \;\middle|\;
      U\subset X \text{ dicht, offen in $X$}
    \right\}
  \end{gather*}
  wobei $\psi\colon U\to Y$ äquivalent zu $\psi'\colon U'\to Y$ ist,
  falls sie auf einer in $X$ dichten, offenen Teilmenge von $U\cap U'$
  übereinstimmen.
  
  Ein offenes Unterschema $U\subset X$ heißt $S$-dicht in $X$,
  falls für alle $s\in S$ die Faser $U_s$ %=U\times_S\k(s)$
  topologisch dicht in der Faser $X_s$ %=X\times_S\k(s)$
  liegt.
  Eine $S$-rationale Abbildung $\phi\colon X\dra Y$ ist eine
  rationale Abbildung, wobei alle dichten offenen Mengen aus obiger
  Definition $S$-dicht sind.
  % Eine \emph{$S$-rationale Abbildung} $\phi\colon X\dra Y$ ist eine
  % Äquivalenzklasse der Menge
  % \begin{gather*}
  %   \left\{
  %     \psi\colon U\to Y \text{ $S$-Morphismus}
  %     \;\middle|\;
  %     U\subset X \text{ $S$-dicht, offen in $X$}
  %   \right\}
  % \end{gather*}
  % wobei $\psi\colon U\to Y$ äquivalent zu $\psi'\colon U'\to Y$ ist,
  % falls sie auf einer $S$-dichten, offenen Teilmenge von $U\cap U'$
  % übereinstimmen.

  Eine ($S$-)rationale Abbildung $\phi\colon X\dra Y$ heißt definiert
  in $x\in X$, falls $\phi$ von einem Morphismus $\psi\colon U\to Y$
  präsentiert werden kann mit $x\in U$.
  Der Definitionsbereich $\Dom(\phi)$ von $\phi$ sind die
  Punkte, in denen $\phi$ definiert ist. Er ist offene, ($S$-)dichte
  Teilmenge von $X$. Ist $X/R$ separiert gibt es einen Morphismus
  $\Dom(\phi)\to Y$, der $\phi$ präsentiert.
  
  $\phi\colon X\dra Y$ heißt ($S$-)birational, falls es von einem
  $S$-Morphismus $\psi\colon U\to Y$ präsentiert werden kann, der
  die offene Immersion einer ($S$-)dichten, offenen Teilmenge von $Y$
  ist.
\end{Definition}

Nachdem unser Hauptaugenmerk auf Gruppenschemata liegt, ist für uns
besonders interessant, wie sich $S$-Gruppenschemata
bzgl. ($S$-)rationalen Abbildungen verhalten.
Tatsächlich verhalten sich eigentliche Gruppenschemata in folgender
Hinsicht gutartig:
\begin{Lemma}[{\cite[Proposition IV.6.2]{silverman2}}]
  \label{thm:rationalzumorphismus}
  Sei $G$ ein $R$"=Gruppenschema, $Y$ ein glattes $R$-Schema und
  $\phi\colon Y\dra G$ eine $R$-rationale Abbildung.
  Dann gilt
  \begin{enumerate}[label=(\roman*)]
  \item $Y\setminus \Dom(\phi)$ hat (reine) Kodimension 1.
  \item Ist $G$ eigentlich, so ist $\phi$ Morphismus.
  \end{enumerate}
\end{Lemma}

% -----

\section{Normale Komposition auf Schemata}
Sei $R$ im folgenden ein Dedekindring.
In manchen Fällen kann der Morphismus einer Gruppenstruktur nicht auf
dem gesamten Schema sondern nur als $R$-rationale Abbildung angegeben
werden. So kommt man zur Definition einer normalen Komposition, welche
eine $R$-rationale Gruppenstruktur liefert.
\begin{Definition}[normale Komposition]\label{def:normalekomposition}
  Sei $V$ glattes $R$-Schema mit nichtleeren Fasern.
  Eine normale Komposition auf $V$ ist eine $R$-rationale
  Abbildung
  \begin{gather*}
    \mu\colon V\times_R V\dra V
  \end{gather*}
  die folgende Bedingungen erfüllt
  \begin{description}[font=\normalfont\itshape]
  \item[(Assoziativität)] Es gilt, wo es definiert ist,
    \begin{gather*}
      \mu(\mu(x,y),z)=\mu(x,\mu(y,z))
    \end{gather*}
  \item[(Inverses)]
    Die $R$-rationalen Abbildungen
    \begin{align*}
      \phi\colon V\times_R V
      &\dra V\times_R V
      &\psi\colon V\times_R V
      & \dra V\times_R V\\
      (x,y)
      &\longmapsto (x,\mu(x,y))
      &(x,y)
      &\longmapsto (y,\mu(x,y))       
    \end{align*}
    haben dichtes Bild in jeder Faser und die Einschränkung auf jede
    Faser ist birational.
  \end{description}
\end{Definition}

Ein großer Vorteil einer normalen Komposition ist, dass zumindest
ein großer Teil, nämlich ein offenes dichtes Unterschema, offen und
dicht in einem richtigen Gruppenschema mit derselben Gruppenstruktur
liegt.
\begin{Satz}[{Weil,
  \cite[Theorem VIII.1.12]{artin},
  \cite[Theorem IV.6.9]{silverman2}}]\label{thm:weil}
  Sei $V$ glattes $R$-Schema mit nichtleeren Fasern
  und normaler Komposition $\mu\colon V\times_R V\dra V$.
  $V$ sei zusätzlich endlichen Typs und $V(R)$ sei dicht in jeder
  Faser. 
  % TODO!! Was heißt V(R) dicht in jeder Faser?
  Dann gibt es ein $R$-Gruppenschema $(G,\mu_G)$ endlichen Typs und
  ein gemeinsames, offenes, $R$-dichtes Unterschemata $V\hookleftarrow
  U\hookrightarrow G$ von $V$ und $G$, so dass $\mu|_U=\mu_G|_U$.
  % \begin{proof}
  %   \begin{enumerate}
  %   \item Fall zusammenhängende Fasern:
  %     \begin{enumerate}
  %     \item $X=X_\phi\cap X_\psi\cap Y_\phi\cap Y_\psi$:
  %       $V$ hat normale Komposition $\mu$ mit zugehörigen rationalen
  %       Abbildungen $\phi$ und $\psi$.
  %       Wähle $R$-dichte, offene Teilmengen
  %       $X_\phi,X_\psi,Y_\phi,Y_\psi\subset V\times_R V$,
  %       so dass $\phi\colon X_\phi\overset\sim\to Y_\phi$,
  %       $\psi\colon X_\psi\overset\sim\to Y_\psi$ Isomorphismen sind.
  %       Definiere $X=X_\phi\cap X_\psi\cap Y_\phi\cap Y_\psi\subset
  %       V^2$.
  %       $X$ ist wieder $R$-dicht und $\phi|_X,\psi|_X\colon X\cong X$
  %       sind per Definition Isomorphismen auf $X$.
  %     \item Es gibt eine offene, $R$-dichte Teilmenge $U\subset V$, so
  %       dass die Einschränkung der Projektionen
  %       $\proj_i\colon V\times_R V\to V$ auf $X$ surjektiv sind.
  %     \item Für $V=U$ gilt:
  %       Sei $a\in V$, $x\in V$ in dichter, offener Teilmenge der
  %       Faser, die $a$ enthält.
  %       Dann sind $\phi^{\pm1}, \psi^{\pm1}$ in $(x,a)\in V\times_R V$
  %       definiert.
  %     \item Sei $\Gamma_\mu\subset V^3$ der Graf von $\mu$
  %       ($R$-Schema), $\Gamma\coloneqq \overline{\Gamma_\mu}\subset
  %       V^3$.
  %       Dann sind die Einschränkungen der Projektionen
  %       $\proj_{1,2},\proj_{2,3},\proj_{1,3}\colon \Gamma\to V^2$
  %       offene Immersionen und dicht in jeder Faser von $V^2$.
  %     \item Verkleben entlang von $W_s=V\times_R s\times_R V$ für
  %       $s\in V$:
  %       …
  %     \item …
  %     \end{enumerate}
  %   \item Stein-Faktorisierung + „dense set of local sections“?:
  %     Erhalte offenes Unterschema 
  %   \end{enumerate}
  % \end{proof}
\end{Satz}

%%% Local Variables:
%%% mode: latex
%%% TeX-master: "neron_models"
%%% End:


% Néron-Modelle
\chapter{Néron-Modelle}\label{chap:neronmodelle}
Sei $R$ im gesamten Kapitel ein Dedekindring mit Quotientenkörper
$\K$.

Nachdem jetzt die nötigen Grundlagen für Gruppenschemata vorhanden
sind, lassen sich Néron-Modelle wie zuvor motiviert definieren.
Dazu wird erst die formale Definition angegeben und für
Gruppenschemata spezialisiert. Danach werden wichtige allgemeine
Eigenschaften untersucht, die direkt aus der universellen Eigenschaft
folgen, unter anderem die Eindeutigkeit und Lokalität.
Zuletzt wird noch einen kurzen Blick auf Néron-Modelle
allgemeiner abelscher Varietäten geworfen.

\section{Definition}
Das Néron-Modell eines $\K$"=Gruppenschemas ist – falls es existiert –
eine besonders gutartige (genauer: glatte, separierte) Erweiterung der
Gruppenstruktur auf ein $R$"=Gruppenschema, das
minimal mit dieser Eigenschaft gewählt ist.

In diesem Abschnitt wird vorerst die allgemeine Definition eines
Néron-Modells nach \cite[Definition~1.2/1]{neron} anhand der
Erweiterungsforderung eingeführt. Daraus leiten sich bereits einige
Eigenschaften ab, die zu einer alternativen Definition für
$\K$"~Grup\-pen\-sche\-ma\-ta führen, welche die Gruppenstruktur wieder mit
einbezieht.

\begin{Definition}[Néron-Modell]
  \optcite[1.2, Definition 1]{neron}
  Sei $R$ ein Dedekindring, $\K=\Quot(R)$ und
  $X_\K$ ein glattes, separiertes $\K$"=Schema endlichen Typs.
  Ein Néron-Modell von $X_\K$ über $R$ ist ein glattes, separiertes
  $R$"=Modell $X$ von $X_\K$ mit der Néron-Abbildungseigenschaft:
  \begin{quote}
    Für jedes weitere glatte $R$"=Schema $Y$ wird jeder
    $\K$"=Morphismus $Y_\K\to X_\K$ auf den generischen Fasern durch
    einen eindeutigen $R$"=Morphismus $Y\to X$ erweitert.
  \end{quote}
  In anderen Worten die kanonische Abbildung $X(Y)\to X_\K(Y_\K)$, die
  durch die Einschränkung induziert wird, ist ein Isomorphismus.
\end{Definition}

% ERWEITERUNG DER GRUPPENSCHEMASTRUKTUR
Die Motivation als minimale Erweiterung von Gruppenschemata ist wie
folgt gerechtfertigt:
\begin{Lemma}\label{thm:gruppenschemaerweiterung}
  \optcite[1.2, Proposition 6]{neron}
  Sei $X$ ein Néron-Modell seiner generischen Faser $X_\K$ über $R$. Ist
  $X_\K$ ein $\K$"~Grup\-pen\-sche\-ma, so wird dessen $\K$"=Gruppenstruktur
  durch die Néron-Abbildungseigenschaft eindeutig auf eine $S$"=Gruppenschemastruktur auf
  $X$ erweitert.
  \begin{proof}[Beweisskizze]
    Die Gruppenverknüpfung $\mu_\K\colon X_\K\times_\K X_\K\to X_\K$
    erweitert sich wegen der Néron-Abbildungseigenschaft zu einem Morphismus
    $\mu\colon X\times_R X\to X$, da $X\times_R X$ wieder glatt ist.
    Genauso liften das Inverse, der Einheitsschnitt und, wegen der
    Eindeutigkeit, die Diagramme aus \ref{def:gruppenschema}, was
    $\mu$ zur Gruppenverknüpfung auf $X$ macht.
  \end{proof}
\end{Lemma}

% NÉRON-MODELL K-GRUPPENSCHEMA
\begin{Lemma}\label{nerongruppenschemaglatt}
  \optcite[1.2, Remark 7]{neron} 
  Ist $X_\K$ ein $\K$"=Gruppenschema und $X$ ein
  glattes $R$"=Modell und Gruppenschema von $X_\K$, das die
  Néron-Abbildungseigenschaft erfüllt, so ist $X$ ein Néron-Modell von $X_\K$.
  \begin{proof}
    Als Faser eines glatten Schemas ist $X_\K$ glatt. Es bleibt also
    zu zeigen, dass $X_\K$ und $X$ separiert sind.
    
    Ein Gruppenschema $X$ ist genau dann separiert, wenn der
    Einheitsschnitt ${\epsilon\in X(R)}$ eine abgeschlossene
    Immersion, siehe
    \cite[Lemma~38.6.1]{stacksproject} oder \cite[Lemma~7.1/2]{neron}.
    Entsprechend sind $\K$"=Gruppenschemata immer separiert, siehe
    \cite[Lemma~38.7.3]{stacksproject}. 

    Sei wegen \ref{thm:neronmodelllokal} ohne Einschränkung $R$ diskreter
    Bewertungsring mit Restklassenkörper $k$.    
    Es wird im Folgenden gezeigt, dass topologisch
    $\im(\epsilon)=\overline{\im(\epsilon_\K)}$,
    also dass $\overline{\im(\epsilon_\K)}$ aus genau dem Punkt
    $\epsilon(\Spec\K)$ in der generischen Faser und dem Punkt
    $\epsilon(\Spec k)$ in der speziellen Faser besteht.
    Dann ist $\im(\epsilon)$ topologisch abgeschlossen in $X$ und der
    Schnitt $\epsilon$ demnach eine abgeschlossene Immersion, da er
    aufgrund der Glattheit von $X$ bereits Homeomorphismus und als
    Schnitt surjektiv auf den Halmen ist.
    
    Auf der generischen Faser gilt
    $\overline{\im(\epsilon_\K)}\cap X_\K=\im(\epsilon_\K)$,
    da $\epsilon_\K$ eine abgeschlossene Immersion ist.
    Außerdem stimmen hier $\im(\epsilon_\K)=\{\epsilon_\K(\Spec\K)\}$
    und $\im(\epsilon)\cap X_\K$ wegen der Néron-Abbildungseigenschaft punktweise
    überein.

    Auf der speziellen Faser $X_k$ ist zu zeigen, dass für jeden
    Punkt $x\in\overline{\im(\epsilon_\K)}\cap X_k$ bereits
    $x\in\im(\epsilon)$ gilt. Sei also $x$ ein solcher Punkt.
    Das schematische Bild $\text{Im}(\epsilon_\K)$ von $\epsilon_\K$ ist
    nach \cite[Remark~10.32]{wedhorn} genau
    $\overline{\im(\epsilon_\K)}=\overline{\{\epsilon_\K(\Spec(\K))\}}$
    mit reduzierter Unterschemastruktur, was
    $\overline{\im(\epsilon_\K)}$ mit einer integren
    Unterschemastruktur ausstattet.
    Sei $U$ eine offene, affine Umgebung von $x$ in $X$ und
    $\Spec(A)=U\cap\overline{\im(\epsilon_\K)}$ ein offenes, affines
    Unterschema von $\overline{\im(\epsilon_\K)}$.
    Später werden die Inklusionen $R\subset A\subset\K$ gezeigt.
    Nachdem $R$ als diskreter Bewertungsring der größte Unterring
    von $\K$ ist, folgt daraus dann $R=A$ und es ergibt sich zusammen
    mit der Immersion zu $\Spec(A)$ ein Schnitt
    \begin{gather*}
      X(R)\ni e\colon
      \Spec(R)\xrightarrow{\sim}\Spec(A)\hookrightarrow X\;,
    \end{gather*}
    der $\epsilon_\K\colon\Spec(\K)\to X_\K$ erweitert.
    Dieser muss allerdings wegen der Eindeutigkeit aus der Néron-Abbildungseigenschaft
    gleich $\epsilon\in X(R)$ sein, also liegt $x$ in
    $\im(e)=\im(\epsilon)$.
    Nachdem $x\in \overline{\im(\epsilon_\K)}\cap X_k$ beliebig
    gewählt war, kann $\overline{\im(\epsilon_\K)}\cap X_k$ nur
    den Punkt $\epsilon(\Spec k)$ enthalten.

    Betrachte also $A$ wie oben, um $R\subset A\subset\K$ zu
    zeigen. Es ist $e_\K\coloneqq\epsilon_\K(\Spec(\K))$ der
    generische Punkt von $\overline{\im(\epsilon_\K)}$ und
    entsprechend auch vom offenen Unterschema $\Spec(A)$. Da
    $\overline{\im(\epsilon_\K)}$ und $\Spec(A)$ integer sind, ist
    $\O_{\overline{\Im(\epsilon_\K)}, e_\K}=\O_{A,e_\K}$ Körper.
    Es folgt eine Sequenz von Körperhomomorphismen auf den Halmen
    \begin{gather*}
      \Id_\K\colon
      \K \xrightarrow{\pi_\K^\#}
      \O_{\Im(\epsilon_\K), e_\K} = \O_{A,e_\K} = \Quot(A)
      \xrightarrow{\epsilon_\K^\#} \K
    \end{gather*}
    weshalb $\Quot(A)\cong\K$ und daher $A\hookrightarrow\K$.
    Für den zu $\pi|_{\Spec(A)}$ assoziierten
    Ringhomomorphismus $\pi^\#=\pi_\K^\#|_R$ kommutiert außerdem wegen
    $\pi\circ\epsilon_\K=\Id_{\Spec(\K)}$ das folgende Diagramm
    \begin{center}
      \begin{tikzcd}
        A\arrow[r, hook]\arrow[from=d, "{\pi^\#}"]
        & \Quot(A)\arrow[d, "{\wr}"{swap}, "{\epsilon_\K^\#}"]\\
        R\arrow[r, hook]
        & \K
      \end{tikzcd}
    \end{center}
    weshalb $\pi^\#\colon R\hookrightarrow A$ gelten muss.
    
    Die Mengen $\overline{\im(\epsilon_\K)}=\im(\epsilon_\K)$ stimmen
    also überein, insbesondere ist $\im(\epsilon_\K)$ eine
    abgeschlossene Immersion.
  \end{proof}
\end{Lemma}

% Elliptische Kurven sind abelsche Varietäten und nach
% \ref{thm:exneronmodellabvarietaet} hat somit jede elliptische
% Kurve ein Néron-Modell. Wir werden uns damit befassen, wie sich diese
% Modelle konkret angeben lassen mithilfe der Weierstraßgleichung der
% Kurve.

\begin{SatzDefinition}[Néron-Modelle elliptischer Kurven]
  \label{def:neronmodellellkurve}
  Sei $E_\K$ eine elliptische Kurve, das heißt ein
  glattes, projektives, insbesondere separiertes, integres
  $\K$"=Schema endlichen Typs vom Geschlecht~1.
  Dann ist ein Néron-Modell von $E_\K$ nach
  \ref{thm:gruppenschemaerweiterung} und \ref{nerongruppenschemaglatt}
  ein glattes $R$"=Modell und Gruppenschema $E/R$ mit der
  Néron-Abbildungseigenschaft:
  \begin{quote}
    Für jedes weitere glatte $R$"=Schema $Y$ wird jeder
    $\K$"=Morphismus ${\phi_\K\colon Y_\K\to E_\K}$ auf den generischen
    Fasern durch einen eindeutigen $R$"=Morphismus ${Y\to E}$ erweitert.
  \end{quote}
  \optcite[Chapter IV.5]{silverman2}
\end{SatzDefinition}
\begin{Bemerkung}\label{thm:dichtefaser}
  Das $\phi_\K$ aus \ref{def:neronmodellellkurve} kann auch als
  $\K$"=rationale Abbildung gewählt werden. Da $E_\K$ projektiv, also
  eigentlich, ist, erweitert sich $\phi_\K$ nach
  \ref{thm:rationalzumorphismus} zu einem Morphismus.

  Im Fall eines diskreten Bewertungsrings $R$ ist der generische Punkt
  $(0)$ von $\Spec(R)$ offen, da das Komplement
  $\Spec(R)\setminus\{(0)\}$ aus nur endlich vielen abgeschlossenen
  Punkten besteht. Entsprechend ist die generische Faser $Y_\K$ jedes
  $R$"=Schemas $Y$ offen.
  Ist $Y$ glatt – also insbesondere flach und reduziert – über $R$,
  ist $Y_\K$ dicht in $Y$ nach \cite[Proposition~III.9.7]{hartshorne}.
  Also präsentiert der Morphismus ${\phi_\K\colon X_\K\to E_\K}$ eine
  rationale Abbildung wie in \ref{def:ratabb}
  \begin{gather*}
    \phi\colon X\dra E
  \end{gather*}
  Die Néron-Abbildungseigenschaft bedeutet, dass diese sich Abbildung eindeutig zu einem
  Morphismus $X\to E$ erweitert.
\end{Bemerkung}


\section{Eigenschaften}
\subsection{Allgemeine Eigenschaften}
Wie oben sei $R$ immer noch ein Dedekindring mit $\K=\Quot(R)$.

% WARUM SEPARIERTHEIT
Ein Grund, weshalb für ein Néron-Modell Separiertheit gefordert wird,
ist die Eindeutigkeit von Morphismen für separierte Schemata im Sinne
des folgenden Satzes. Für den Beweis siehe
\cite[Corollary~9.9]{wedhorn}.
\begin{Satz}\label{thm:erweindeutig}
  Seien $Y$ und $X$ zwei $R$"=Schemata, sei $X$ separiert und $Y$
  reduziert.
  Dann ist ein Morphismus $\phi\colon Y\to X$ durch jede
  Einschränkung auf eine dichte Teilmenge von $Y$ eindeutig bestimmt.
  Insbesondere sind, falls $Y$ flach über $X$ ist, Erweiterungen von
  Morphismen der generischen Fasern eindeutig.
  % \begin{proof}
  %   Sei $x\in\Dom(\phi)$.
  %   Wir werden induktiv nach der Dimension von $x$ vorgehen. Für
  %   $\dim_Y(x)=0$ ist $x$ generischer Punkt und $\phi$ ist hier nach
  %   Voraussetzung festgelegt.
  %   Sei $\phi$ eindeutig definiert in $y\in\Dom(\phi)$ für
  %   $\dim_Y(y)=(n-1)$ und sei $\dim_Y(x)=n$.
  %   Wir können nun eine Generalisierung $y\in\Spec(\O_{Y,x})\subset Y$
  %   von $x$ mit $\dim_Y(y)=(n-1)$ betrachten.
  %   Im geschlossenen, irreduziblen, reduzierten Unterschema $V_y$ mit
  %   topologischem Raum $\overline{\{y\}}\subset Y$ ist
  %   $\dim_{V_y}(x)=1$.
  %   $V_y$ ist integer, $\Spec(\O_{V_y,x})=\{x,y\}$ und daher 
  %   $\O_{V_y,x}$ ein Bewertungsring mit
  %   $y=(0)\in\Spec(\O_{V_y,x})$ und $\Quot(\O_{V_y,x})=\k(y)$.
  %   Nach dem Bewertungskriterium für Separiertheit von $W$ ist
  %   $\phi|_{\Spec(\O_{V_y,x})}$ die eindeutige Erweiterung von
  %   $\phi|_{\Spec(\k(y))}$. Da $\phi$ auf $y$ bereits eindeutig war
  %   und $Y$ reduziert ist, ist $\phi$ auch eindeutig auf $x$.

  %   Ist $Y$ flach, so liegen alle generischen Punkte von $Y$ in der
  %   generischen Faser nach \cite[Proposition III.9.7]{hartshorne}.
  % \end{proof}
\end{Satz}
Somit ist die Eindeutigkeit der Erweiterung eines Schnitts bereits
durch die Separiertheit des Néron-Modells gegeben. Außerdem ist nun
bekannt, dass die kanonische Einschränkung $X(Y)\to X_\K(Y_\K)$ der
$R$"=Schnitte eines Néron-Modells $X$ nicht nur für glatte, sondern
auch für reduzierte Schemata $Y$ injektiv ist.

% EINDEUTIGKEIT, BASISWECHSEL
Nun folgen einige weitere grundlegende Eigenschaften von
Néron-Modellen. Hierbei ist zu beachten, dass sich Néron-Modelle nicht
so gutartig verhalten, wie es die universelle Eigenschaft vermuten lässt.
Unter anderem sind sie zwar stabil unter Produktbildung (hier
nicht benötigt und nicht gezeigt), allerdings nicht unter
beliebigem Basiswechsel.
\begin{Satz}\label{thm:eigneronmodell}
  \optcite[1.2, Proposition 2]{neron}
  \optcite[Proposition IV.5.2]{silverman2}
  Sei $X$ ein Néron-Modell seiner generischen Faser $X_\K$ über
  $R$. Dann gilt
  \begin{enumerate}[label=(\alph*)]
  \item $X$ ist eindeutig bis auf kanonische Isomorphie.
  \item Für einen étalen Basiswechsel $R\to R'$ mit $K'$
    Funktionenkörper von $R'$ ist $X_{R'}=X\times_R R'$ ein Néron-Modell
    seiner generischen Faser $X_{\K'}=X_\K\times_{\K}\K'$ über $R'$,
    das heißt étaler Basiswechsel kommutiert mit der Bildung von
    Néron-Modellen.
  \end{enumerate}
  \begin{proof}
    Die erste Eignschaft folgt direkt aus der universellen
    Eigenschaft, für die zweite sind weitere Argumente zur Stabilität
    unter Basiswechsel nötig. 
    \begin{enumerate}[resume*,start=1]
    \item Sei $X'$ ein weiteres Néron-Modell von $X_\K$.
      Dann liefert die Néron-Abbildungseigenschaft zur Identität
      ${\Id_{X_\K}\colon X_\K\to X_\K}$ eindeutige Morphismen
      ${\phi\colon X'\to X}$ und ${\psi\colon X\to X'}$.
      Die Verknüpfungen ${\phi\circ\psi\colon X\to X}$ und
      ${\psi\circ\phi\colon X'\to X'}$ sind jeweils die Identität auf
      der generischen Faser $X_\K$, welche nach der Eindeutigkeit in der
      Néron-Abbildungseigenschaft eindeutig durch die Identitäten auf
      $X$ und auf $X'$ erweitert wird. Also sind $\phi,\psi$
      zueinander inverse Isomorphismen.
    \item
      $X_{R'}$ ist glatt, separiert und von endlichem Typ, das heißt
      lokal von endlichem Typ und quasikompakt, da diese Eigenschaften
      stabil unter Basiswechsel sind. Es ist somit nur die
      Néron-Abbildungseigenschaft zu prüfen.
      Sei also $Y'$ glattes $R'$"~Sche\-ma mit generischer Faser
      $Y_{\K'}'$ und ${\phi_{\K'}\colon Y_{\K'}'\to X_{\K'}}$ ein
      $\K'$"=Morphismus. 
      \begin{description}
      \item[Existenz:] 
        $Y'$ ist ein glattes $R$"=Schema durch die Komposition
        \begin{gather*}
          Y'\xrightarrow{\text{glatt}}\Spec(R')
          \xrightarrow{\text{glatt}}\Spec(R)
        \end{gather*}
        glatter Abbildungen.
        Außerdem lässt sich der $\K$"=Morphismus
        \begin{gather*}
          \phi_{\K} \coloneqq \proj_{X_\K}\circ\phi_{\K'}
          \colon
          Y_{\K'}' \xrightarrow{\phi_{\K'}} X_{\K'}=X\times_\K\K'
          \xrightarrow{\proj_{X_\K}} X_{\K}
        \end{gather*}
        mit der Néron-Abbildungseigenschaft von $X$ zu einem eindeutigen
        $R$"=Morphismus ${\phi_R\colon Y'\longto X}$ erweitern.
        So ergibt sich ein kanonischer, eindeutiger $R$"~Mor\-phis\-mus
        (beziehungsweise ein $R'$"=Morphismus)
        $\phi_{R'} = (\phi_R\;,\;\pi_{R'})$ auf dem Faserprodukt 
        \begin{center}
          \begin{tikzcd}
            Y'
            \arrow[drr, bend left, "\pi_{R'}"]
            \arrow[ddr, bend right, "{\phi_R}"]
            \arrow[dr, dashed, "{\exists!\phi_{R'}}"]\\
            & X\times_{R} R'
            \arrow[dr, phantom, "\ulcorner", very near start]
            \arrow[r, "\proj_{R'}"]\arrow[d, "\proj_X"]
            & \Spec(R')\arrow[d]\\
            & X \arrow[r]
            & \Spec(R)
          \end{tikzcd}
        \end{center}
        $\phi_{R'}$ erweitert $\phi_{\K'}$, denn
        \begin{align*}
          \phi_{R'}|_{Y_{\K'}'}
          &= \left(
            \phi_R|_{Y_{\K'}'}\;,\;
            \pi_{R'}|_{Y_{\K'}'}\colon
            Y_{\K'}'\hookrightarrow Y'\to\Spec(R')
            \right)\\
            % \pi|_{Y_{\K'}'})\colon Y'\times_{R'}\K'\to Y'\to\Spec(R')
            % \overset{\text{Faserprod.}}{=}
            % \pi_{\K'}\colon Y'\times_{R'}\K'\to \Spec(K')\to\Spec(R')
          &\overset{\mathllap{\text{Faserprod.}}}{=} \left(
            \phi_R|_{Y_{\K'}'}\;,\;
            \pi_{\K'}\colon
            Y_{\K'}'\to \Spec(K')\hookrightarrow\Spec(R')
            \right)\\
          &\overset{\mathllap{\text{Def. }\phi_R}}{=} \left(
            \phi_{\K}\;,\; \pi_{\K'}
            \right)\\
          &\overset{\mathllap{\text{Def. }\phi_\K}}{=} \left(
            \proj_{X_\K}\circ\phi_{\K'}\;,\;
            \proj_{\K'}\circ\phi_{\K'}
            \right)\\
          &= \phi_{\K'}
        \end{align*}
      \item[Eindeutigkeit:]
        Für jeden weiteren Lift $\psi_{R'}$ von $\phi_{\K'}$ ist
        $\proj_{X}\circ\psi_{R'}$ ein Lift von $\phi_{K}$ und es gilt
        aufgrund der Eindeutigkeit in der Néron-Abbildungseigenschaft
        ${\proj_{X}\circ\psi_{R'}=\phi_{R}}$.
        Damit lässt $\psi$ ebenfalls das obige Diagramm kommutieren
        und ist nach der universellen Eigenschaft des Faserprodukts
        bereits gleich $\phi_{R'}$.
        \qedhere
      \end{description}
    \end{enumerate}
  \end{proof}
\end{Satz}

% LOKALITÄT
Eine sehr wichtige und praktische Eigenschaft von Néron-Modellen ist
die Lokalität, also die Möglichkeit, lokal zu untersuchen, ob es sich
um ein Néron-Modell handelt.
\begin{Satz}\label{thm:neronmodelllokal}
  Für ein $R$"=Schema $X$ endlichen Typs ist äquivalent:
  \begin{enumerate}[label=(\roman*)]
  \item $X$ ist das Néron-Modell seiner generischen Faser.
  \item $X\times_R\Spec(\O_{\Spec(R),s})$ ist das Néron-Modell
    seiner generischen Faser über $\O_{R,s}=R_s$ für alle
    abgeschlossenen Punkte $s\in\Spec(R)$.
  \end{enumerate}
  Das heißt, es kann ohne Einschränkung immer davon ausgegangen
  werden, dass $R$ lokal ist, also ein diskreter Bewertungsring.
\end{Satz}

Im Folgenden wird nur die Implikation
\emph{(ii)}$\Rightarrow$\emph{(i)} benötigt, daher sei für die
andere Richtung auf \cite[Proposition~1.2/4]{neron} verwiesen.
Erst ist dazu das folgende Lemma und Korollar zur Erweiterung von
Morphismen auf offene Umgebungen nötig.

% ERWEITERUNG V. MORPHISMEN AUF OFFENE TEILMENGEN
\begin{Lemma}\label{thm:morphismuserweiterung}
  % http://math.stackexchange.com/questions/1259876/extending-a-morphism-of-schemes
  \optcite[Exercise 3.2.4]{liu}
  Sei $S$ ein lokal noethersches Schema,
  seien $X$ und $Y$ zwei $S$"=Schemata und
  sei $X$ endlichen Typs.
  Seien weiterhin $y\in Y$ und 
  ein $S$"=Morphismus ${f_y\colon \Spec\O_{Y,y}\to X}$ gegeben.
  Dann gibt es eine offene Umgebung ${U\subset Y}$ von $y$, auf der
  $f_y$ von einem Morphismus ${f\colon U\to X}$ erweitert wird.
  \begin{proof}
    Es reicht, sich auf den affinen Fall $X=\Spec(B)$ zu beschränken.
    Denn $\O_{Y,y}$ ist lokal mit Maximalideal $y$, also ist $y$ in
    jeder Zariski-abgeschlossenen, nichtleeren Menge von
    $\Spec\O_{Y,y}$ enthalten. Anders gesagt die einzige offene
    Umgebung von $y$ in $\Spec\O_{Y,y}$ ist $\Spec\O_{Y,y}$ selbst.
    Somit ist das offene Urbild jeder offenen, affinen Umgebung
    ${\Spec(B)\subset X}$ von $f_y(y)$ ganz $\Spec\O_{Y,y}$.

    $X$ ist nach Voraussetzung von endlichem Typ und $S$ lokal
    noethersch.
    Es kann also eine genügend kleine, offene, affine Umgebung
    ${\Spec(B)\subset X}$ von $f_y(y)$ und eine offene Teilmenge
    ${\Spec(R)\subset S}$ ausgewählt werden,
    so dass $R$ noethersch und $B$ eine $R$"=Algebra endlichen Typs
    – wegen noethersch auch endlich präsentiert – ist. Das heißt $B$
    hat die Form
    \begin{gather*}
      B=R[X_0,\dotsc, X_n]/(g_0,\dotsc, g_r)\qquad n,r\in\N
    \end{gather*}
    Wähle weiterhin eine genügend kleine, offene, affine Umgebung
    ${\Spec(A)\subset Y}$ von $y$ aus, dass $A$ ebenfalls $R$"=Algebra ist.
    Hier gilt ${y\in\Spec\O_{Y,y}=\Spec(A_y)\subset\Spec(A)}$.
    
    Der Morphismus ${f_y\colon\Spec(A_y)\to\Spec(B)\subset X}$
    korrespondiert also zu einem Ringhomomorphismus
    ${\Gamma(f_y)\colon B\to A_y}$.
    Insgesamt ergibt sich
    \begin{center}
      \begin{tikzcd}[row sep=0pt]
        \phi\colon R\left[X_0,\dotsc, X_n\right] \arrow[r,two heads]
        &\frac{R\left[X_0,\dotsc, X_n\right]}{\left(g_0,\dotsc, g_r\right)}
        \arrow[r,"{\Gamma(f_y)}"]
        &A_y\\
        X_i \arrow[rr,mapsto]&& \frac{\alpha_i}{\beta_i} &i=0,\dotsc, n\\
        g_j \arrow[rr,mapsto]&& \frac{\alpha_j'}{\beta_j'} &j=0,\dotsc, r
      \end{tikzcd}
    \end{center}
    Die Bedingung $\frac{\alpha_j'}{\beta_j'}=\Gamma(f_y)(g_j)=0\in A_y$
    heißt, dass es $\gamma_j$ gibt, so dass
    $\alpha_j'\gamma_j=0\in A$.
    Definiere
    $\gamma\coloneqq \prod_{i=0}^n \beta_i \prod_{j=0}^r \gamma_j$.
    Dann gilt für den Morphismus
    \begin{align*}
      \phi_\gamma\colon
      R\left[X_0,\dotsc, X_n\right]
      &\longto
        A_\gamma\\
      X_i
      &\longmapsto
        \dfrac{
        \alpha_i\cdot
        \prod_{k\neq i}^n \beta_k
        \prod_{j=0}^r \gamma_j
        }{f}
        \qquad i=0,\dotsc,n\;,
    \end{align*}
    dass $(g_0, \dotsc, g_r)\subset\ker(\phi_f)$ und das folgende
    Diagramm kommutiert
    \begin{center}
      \begin{tikzcd}
        R\left[X_0,\dotsc, X_n\right] \arrow[r,"\phi_\gamma"]
        \arrow[d, two heads]
        & A_\gamma\arrow[d]\\
        B \arrow[r, "\Gamma(f_y)"{swap}]\arrow[ur,"\overline\phi_\gamma"]
        & A_y
      \end{tikzcd}\,.
    \end{center}
    Insbesondere faktorisiert ${\Gamma(f_y)\colon B\to A_\gamma\to A_y}$
    über $A_\gamma$, was heißt, dass $f_y$ über einen Morphismus
    ${f\colon D(\gamma)\to\Spec(B)}$ faktorisiert:
    \begin{gather*}
      f_y\colon
      \Spec\O_{Y,y}=\Spec(A_y)
      \hookrightarrow \Spec(A_\gamma)=D(\gamma)
      \overset{f}\longto \Spec(B)
    \end{gather*}
    $D(\gamma)$ ist eine offene Umgebung von $y$, womit die Existenz
    gezeigt ist.
  \end{proof}
\end{Lemma}

\begin{Korollar}\label{thm:allgmorphismuserweiterung}
  Sei $S$ ein Dedekindschema,
  sei $\pi\colon Y\to S$ ein integres, flaches $S$"=Schema
  und $X$ ein separiertes $S$"=Schema endlichen Typs.
  Sei für ein $s\in S$ ein $S$"=Morphismus
  \begin{gather*}
    f_s\colon Y\times_S\O_{S,s}\to X\times_S\O_{S,s}
  \end{gather*}
  gegeben.
  Dann gibt es eine offene Umgebung $U\subset Y$ von
  ${Y\times_S\O_{S,s}}$, auf der $f_s$ eindeutig von einem Morphismus
  ${f\colon U\to X}$ erweitert wird.
  Insbesondere gilt dies für glatte $Y$, da $Y$ dann die disjunkte
  Vereinigung integrer, flacher $S$"=Schemata ist.
  \begin{proof}
    Sei $y\in Y\times_S\O_{S,s}\subset Y$ und
    $s'\coloneqq\pi(y)\in\Spec\O_{S,s}$ die zugehörige Spezialisierung
    von $s$. Dann ist $y\in Y\times_S\O_{S,s'}$ und $\pi$ induziert einen
    Ringhomomorphismus ${\O_{S,s'}\to\O_{Y,y}}$ auf den Halmen. Der
    zugehörige Morphismus
    ${\Spec\O_{Y,y}\to\Spec\O_{S,s'}\subset\Spec\O_{S,s}}$
    von affi\-nen $R$"=Schemata ist die Einschränkung von $\pi$ auf
    ${\Spec\O_{Y,y}\subset Y}$, insbesondere  liegt ganz $\Spec\O_{Y,y}$
    im Urbild von $\Spec\O_{S,s}$.
    Das zeigt, dass ${\Spec\O_{Y,y}\subset Y\times_S\O_{S,s}}$ für jeden
    Punkt ${y\in Y\times_S\O_{S,s}}$, und die Einschränkung von $f_s$
    liefert für jedes solches $y$ einen $R$"=Morphismus
    \begin{gather*}
      \SwapAboveDisplaySkip
      f_{s,y}\colon \Spec\O_{Y,y}\longto X\times_S\O_{S,s}\subset X
      \qquad \text{für alle } y \in Y\times_S\O_{S,s}
    \end{gather*}
    Jeder davon erweitert sich nach \ref{thm:morphismuserweiterung}
    zu einem Morphismus auf einer offenen Umgebung ${U_y\subset Y}$ von
    $y$.
    Diese Erweiterungen sind eindeutig aufgrund der Separiertheit.
    Die Definitionsmengen $U_y$ sowie ihre Schnitte ${U_y\cap
    U_{y'}\subset Y}$, ${y,y'\in Y}$, sind offene, integre
    Unterschemata mit demselben generischen Punkt $\eta$ wie $Y$, der
    entsprechend dicht in allen $U_y$ liegt.
    Da $Y$ nach Voraussetzung flach ist, liegt $\eta$ außerdem in
    ${Y\times_S\O_{S,s}}$ nach \cite[Proposition~14.14]{wedhorn}, sprich
    im Definitionsbereich von $f_s$. Da $X$ separiert ist, sind
    deswegen nach \ref{thm:erweindeutig} alle Erweiterungen von
    $f_s|_{U_y}$ auf ganz $U_y$ und von $f_s|_{U_y\cap U_{y'}}$ auf
    ${U_y\cap U_{y'}}$ eindeutig durch $f_s|_{\{\eta\}}$ bestimmt.
    Insbesondere sind die Erweiterungen $f_{s,y}$ eindeutig und
    stimmen auf den Schnitten ihrer Definitionsmengen überein.
    Daher verkleben sie zu einem eindeutigen Morphismus
    \begin{gather*}
      f\colon U=\bigcup_{\mathclap{y\in Y\times_S\O_{S,s}}} U_y\longto X
    \end{gather*}
    der $f_s$ erweitert.
  \end{proof}
\end{Korollar}
% \begin{Lemma}\label{thm:limesoffenemengen}
%   \cite[1.2, Lemma 5]{neron}
%   Für ein Basisschema $S$, einen Punkt $s\in S$ und endlich
%   präsentierte $S$"=Schemata $X$, $Y$ ist die kanonische
%   Abbildung
%   \begin{gather*}
%     \varinjlim \Hom_{S'}(X\times_S S', Y\times_S S')
%     \longto
%     \Hom_{\O_{R,s}}(X\times_S \O_{R,s}, Y\times_S \O_{R,s})
%   \end{gather*}
%   bijektiv, wobei der direkte Limes über die offenen Umgebungen
%   $S'\subset S$ von $s$ läuft.
% \end{Lemma}

\begin{proof}[Beweis von \ref{thm:neronmodelllokal}]
  Vorerst ist zu zeigen, dass $X$ glatt und separiert ist.
  Letzteres gilt, da das Bewertungskriterium für Separiertheit
  erfüllt ist: Sei $(A,\m)$ ein Bewertungsring und über
  ${\pi\colon\Spec(A)\to\Spec(R)}$ eine $R$"=Algebra mit ${\pi(\m)=s}$.
  Dann faktorisiert $\pi$ über
  ${\Spec(R_s)}$ und für jeden Morphismus ${\Spec(\Quot(A))\to X}$
  faktorisiert eine Erweiterung ${\Spec(A)\to X}$ über
  ${X\times_R\O_{R,s}}$:
  \begin{gather*}
    \Spec(A)\to X_s\hookrightarrow X\times_R\O_{R,s}\hookrightarrow X
  \end{gather*}
  Hier ist sie eindeutig wegen Separiertheit von ${X\times_R\O_{R,s}}$.

  Für Glattheit auf $X$ ist nur noch nötig, dass $X$ ein flaches
  $R$"=Schema ist, da $R$ noethersch und $X$ nach Voraussetzung von
  endlichem Typ, also endlich präsentiert, ist.
  Für ein ${s\in\Spec(R)}$ und ein ${x\in X_s\subset X}$ ist der Halm
  ${\O_{X,x}=\O_{(X\times_R\O_{R,s}),x}}$ wegen der Glattheit von
  ${X\times_R\O_{R,s}}$ bereits flach über $R_s$.

  Es bleibt die Néron-Abbildungseigenschaft zu zeigen.
  Sei ${\K=\Quot(R)}$ und betrachte ein glattes $R$"=Schema $Y$ sowie
  einen $\K$"=Morphismus ${f_\K\colon Y_\K\to X_\K}$ der
  generischen Fasern.
  Als glattes $R$"=Schema ist $Y$ flach, reduziert und insbesondere
  regulär, weshalb die irreduziblen Komponenten alle offen und
  disjunkt sind. Es ist also ohne Einschränkung anzunehmen, dass
  $Y$ integer ist.
  Nach Voraussetzung erweitert sich $f_\K$ für einen abgeschlossenen
  Punkt $s\in\Spec(R)$ eindeutig zu einem Morphismus
  \begin{gather*}
    f_s\colon Y\times_R\O_{R,s}\to X\times_R\O_{R,s}\;,
  \end{gather*}
  da $X\times_R\O_{R,s}$ Néron-Modell seiner generischen Faser $X_\K$
  ist.
  % Nach \ref{thm:limesoffenemengen} ist $u_s$ die Einschränkung
  % einer eindeutigen Äquivalenzklasse im direkten Limes
  % $\varinjlim \Hom_{S'}(X\times_S S', Y\times_S S')$ über die
  % offenen Umgebungen $S'\subset S$ von $s$.
  % Daher gibt es eine offene Umgebung $S'(s)$ von $s$, über der $u_s$
  % eindeutig auf einen $R$"=Morphismus $u(s)\colon S'(s)\to X$
  % erweitert wird.
  % % ??? Warum kann ich eine offene Menge wählen?
  Nach \ref{thm:allgmorphismuserweiterung} gibt es eine offene
  Umgebung $U_s$ von $Y\times_R\O_{R,s}$, auf der
  $f_s|_{Y\times_R\O_{R,s}}$ wiederum
  eindeutig auf einen Morphismus $u_s\colon U_s\to X$ erweitert
  wird.
  Die $u_s$ stimmen wieder wegen Separiertheit von $X$
  nach \ref{thm:erweindeutig} auf dem Schnitt ihrer
  Definitionsbereiche überein. Zudem überdecken sie ganz $Y$,
  da $\Spec(R)$ genau aus seinen abgeschlossen Punkten und dem
  generischen Punkt besteht, und letzterer in jeder offenen Teilmenge
  enthalten ist. 
  Daher liefert Verkleben der eindeutigen $u_s$ einen
  eindeutig bestimmten Morphismus $f\colon Y\to X$, der
  $f_\K$ erweitert.
  Das heißt die Néron-Abbildungseigenschaft ist erfüllt.
\end{proof}


\subsection{Néron-Modelle abelscher Varietäten}
Nachdem das Hauptaugenmerk der Arbeit explizit auf Néron-Modellen
abelscher Varietäten liegt, seien hier einige Eigenschaften aufgeführt.
Wichtig ist unter anderem, dass die Existenz in diesem Fall allgemein
gesichert ist, was für elliptische Kurven später explizit nachgewiesen
wird.
$R$ sei in diesem Abschnitt wieder ein Dedekindring und
$\K\coloneqq\Quot(R)$.
% ALLG. EXISTENZ
% \begin{Satz}[Existenz von Néron-Modellen]\label{thm:existenz}
%   \cite[1.3, Theorem 1]{neron}
%   Sei $\Rsh$ die strikte Henselisierung von $R$ und $\Ksh$ der
%   Quotientenkörper. Ein glattes $\K$"=Gruppenschema $X_\K$ endlichen
%   Typs hat genau dann ein Néron-Modell über $R$, wenn $X_\K(\Ksh)$
%   beschränkt in $X_\K$ ist.
%   %   eigentlich => beschränkt
% \end{Satz}
Zur Existenz von Néron-Modellen im Fall abelscher Varietäten besagt
\cite[Corollary~1.3/2]{neron} vorerst ohne explizite Konstruktion:
\begin{Satz}\label{thm:exneronmodellabvarietaet}
  Abelsche Varietäten über $\K$ haben ein Néron-Modell über $R$.
  % \begin{proof}
  %   Abelsche Varietäten über $\K$ sind insbesondere glatte und
  %   eigentliche $\K$"=Gruppenschemata, siehe
  %   \ref{def:abelschevarietaet}.
  %   Nachdem eigentlich beschränkt impliziert, ist die Bedingung aus
  %   \ref{thm:existenz} erfüllt.
  % \end{proof}
\end{Satz}

% Für die folgenden Aussagen ist beim Übergang zum allgemeinen Fall
% eines allgemeinen Dedekindschemas die Eigenschaft zusammenhängend
% nötig.

Abelsche Varietäten unterscheiden sich von generellen glatten,
separierten $\K$"=Gruppenschemata dadurch, dass die Gruppen ihrer
Schnitte abelsch sind. Da diese Eigenschaft für Gruppenschemata
sehr wünschenswert ist, stellt sich die Frage, inwieweit das
Néron-Modell einer abelschen Varietät wieder abelsch ist. 
\begin{Satz}\label{thm:abelscheneronmodelle}
  Sei $A_\K$ eine abelsche Varietät über $\K$ mit guter Reduktion,
  das heißt es gibt ein glattes, eigentliches $R$"=Modell $A$ von $A_\K$.
  Dann ist $A$ das Néron-Modell von $A_\K$ und die induzierte
  Gruppenschemastruktur macht $A$ zu einem abelschen Schema.
  \begin{proof}[Beweisskizze] Für den vollständigen Beweis siehe
    \cite[Proposition~1.4/2]{neron}.
    Nach \ref{thm:exneronmodellabvarietaet} existiert ein
    Néron-Modell $X$ von $A_\K$, das nach
    \ref{thm:gruppenschemaerweiterung} eine Gruppenschemastruktur
    trägt.
    Es bleibt zu zeigen, dass $A$ isomorph zu $X$ ist
    mit zusammenhängenden Fasern, denn dann ist $A$ ist abelsches
    Schema, da es nach Voraussetzung bereits glatt und eigentlich ist.
    \begin{description}[font=\normalfont\itshape]
    \item[$A$ ist isomorph zu $X$:]
      Die Néron-Abbildungseigenschaft von $X$ liefert zur Identität
      $\Id_{A_\K}$ einen $R$"=Morphismus $f\colon A\to X$, der bereits
      offene Immersion, siehe dazu \cite[4.3/1~(ii) oder 4.4/1]{neron}.
      % ?? Wieso ist die Erweiterung von id hier offene Immersion?
      Es ist zu zeigen, dass das Bild von $f$ ganz $X$ ist, woraus
      wegen Reduziertheit von $X$ bereits folgt, dass $f$ ein
      Isomorphismus ist.
      $X$ ist glatt, also nach
      \ref{thm:eigglatt} flaches und lokal endlich präsentiertes
      $R$"=Schema.
      Als solches ist es (universell) offen nach
      \cite[Theorem~14.33]{wedhorn} und somit zusammenhängend nach
      \cite[Proposition~3.24]{wedhorn},
      da die generische Faser $A_\K$ zusammenhängend und $\Spec(R)$
      ein integres Schema ist.
      Es ist also nur zu zeigen, dass das Bild von $f$ abgeschlossen
      ist, da es dann als offene und abgeschlossene Teilmenge des
      zusammenhängenden $X$ schon ganz $X$ sein muss.
      $A$ ist nach Voraussetzung eigentlich, weswegen ${A\to R}$
      universell abgeschlossen und insbesondere der Basiswechsel
      ${A\times_R X\to X}$ abgeschlossen ist.
      Nachdem $X$ separiert ist, ist der Graph
      ${\Gamma_f\colon A\to A\times_R X}$ und daher das Bild von
      \begin{gather*}
        f\colon A\xrightarrow{\Gamma_f} A\times_R X
        \xrightarrow{f\times_R\Id} X
      \end{gather*}
      abgeschlossen in $X$.
    \item[$A$ hat zusammenhängende Fasern:]
      Ohne Einschränkung sei $R$ ein diskreter Bewertungsring nach
      \ref{thm:neronmodelllokal}.
      Dann ist nur noch die spezielle Faser über dem abgeschlossenen
      Punkt zu überprüfen, die nach \cite[{}5.5.1]{EGAIII-1} zusammenhängend
      ist.
      % ??? –> Zariski's Connectedness Thm/Stein-Faktoriesierung?
      % \cite[Exercise 12.29]{wedhorn}
      \qedhere
    \end{description}
  \end{proof}
\end{Satz}

Mit diesem Ergebnis kann für Néron-Modelle abelscher Varietäten
eine Aussage dazu getroffen werden, wo die induzierte
Gruppenschemastruktur abelsch ist.
\begin{Satz}\label{thm:teilwabelscheneronmodelle}
  Sei $A_\K$ eine abelsche Varietät über $\K$
  und $X$ ihr Néron-Modell über $R$.
  Sei weiterhin ${S\subset\Spec(R)}$ die Menge des generischen Punkts
  zusammen mit allen abgeschlossenen Punkten ${s\in\Spec(R)}$, über
  denen $A_\K$ gute Reduktion hat, also für die $A_\K$ ein glattes,
  eigentliches $R_s$"~Mo\-dell hat.
  $S$ ist entsprechend offen und dicht in $\Spec(R)$.
  Dann ist $X\times_R S$ ein abelsches Schema über $S$.
  \begin{proof}[Beweisskizze] Siehe \cite[Theorem~1.4/3]{neron} für
    den vollständigen Beweis.
    Es gilt zu zeigen, dass $A_\K$ in einer Umgebung $U_s$ jedes
    Punktes $s\in S$ gute Reduktion hat, also die generische Faser eines
    eigentlichen, glatten Schemas über $U_s$ ist.
    Diese Schemata sind dann nach \ref{thm:abelscheneronmodelle}
    bereits jeweils über $U_s$ Néron-Modelle und abelsch.
    Wegen der Néron-Abbildungseigenschaft und der resultierenden Eindeutigkeit von
    Néron-Modellen über den offenen Teilmengen von $\Spec(R)$
    lassen sie sich zu einem abelschen $S$"=Schema verkleben.

    In den abgeschlossenen Punkten von $S$ hat $A_\K$ nach Definition
    gute Reduktion, welche sich auf eine offene Umgebung erweitern
    lässt.
    % Gute Reduktion in s\in S <=> X\times_S\O_{S,s} eigentlich
    % (\ref{thm:neronmodell}, \ref{thm:abelscheneronmodelle}
    % Erweiterung von Morphismen auf Umgebung:
    % \ref{thm:allgmorphismuserweiterung}
    % ?? Wieso erhält Morphismuserweiterung Eigentlichkeit?
    Für den generischen Punkt kann der schematische
    Abschluss $A$ des nach \ref{thm:abvarietaetprojektiv} projektiven
    $A_\K$ in $\P_R^n$ betrachtet werden, das heißt das schematische
    Bild von
    \begin{gather*}
      \SwapAboveDisplaySkip
      A_\K\hookrightarrow\P_\K^n\to\P_R^n
    \end{gather*}
    Hier ist es explizit der Abschluss des Bildes mit reduzierter
    Unterschemastruktur nach \cite[Remark~10.32]{wedhorn}, da $A_\K$
    reduziert ist.
    Dieser ist glatt über dem generischen Punkt von $\Spec(R)$
    und somit auch über einer offenen Umgebung $S'$ dessen.
    Insgesamt ist es ein glattes, projektives, also eigentliches,
    $S'$"=Modell von $A_\K$ um den generischen Punkt, was gesucht war.
  \end{proof}
\end{Satz}

% \section{Beispiele}

% \begin{Satz}
%   \cite[1.2, Proposition 8]{neron}
%   Ein abelsches Schema über $\Spec(R)$ ist Néron-Modell seiner
%   generischen Faser.
% \end{Satz}

% \begin{Satz}
%   \cite[1.2, Criterion 9]{neron}
%   Ein glattes, separiertes $\Spec(R)$"=Schema endlichen Typs ist genau
%   dann ein Néron-Modell seiner generischen Faser, wenn es die
%   Erweiterungseigenschaft für étale Punkte erfüllt.
% \end{Satz}

% ----------


%%% Local Variables:
%%% mode: latex
%%% TeX-master: "neron_models"
%%% End:


% Weierstraßmodelle
\chapter{Weierstraßmodelle}\label{chap:weierstraßmodelle}
In diesem Kapitel wird näher auf die Konstruktion eines
projektiven $R$"~Schemas $W$ aus einer Weierstraßgleichung über dem
Quotientenkörper eines Dedekindrings eingegangen.
Ein solches Schema ist zwar eigentlich – erfüllt also die \NAbbEig – ist
aber nicht notwendigerweise glatt und die Gruppenverknüpfung
erweitert sich nicht immer zu einem Morphismus auf ganz $W$.
Immerhin letzteres ist auf dem Unterschema $W^0$ der glatten Punkte
von $W$ behoben. Denn hier erweitert sich die Gruppenverknüpfung und
macht $W^0$ zum Gruppenschema, wie in \autoref{chap:weierstraßgruppe}
zu sehen sein wird.

Vorweg sei schon gesagt, dass $W^0$ nur groß genug für die
\NAbbEig ist, wenn $W$ regulär ist. Allerdings wird sich aus
der späteren Konstruktion ergeben, dass $W^0$ immer ein großer Teil
des Néron-Modells ist.

Die behandelte Konstruktion des Néron-Modells einer elliptischen Kurve
beginnt in jedem Fall mit einem Weierstraßmodell wie unten beschrieben.

\section{Konstruktion}
\subsection{Weierstraßgleichungen}
Jede elliptische Kurve $E_\K$ über einem beliebigen Körper $\K$ kann
durch eine homogene Gleichung der Form
\begin{gather*}
  0 = Y^2 Z + \c_1 XYZ + \c_3 YZ^2 - X^3 - \c_2X^2 Z - \c_4 XZ^2 - \c_6 Z^3
  \eqqcolon F\in\K[X,Y,Z]
\end{gather*}
mit Koeffizienten $\c_1,\dotsc,\c_6\in\K$ und Nullpunkt $O=[0,1,0]$
dargestellt werden. Das heißt sie ist isomorph zur projektiven Kurve
$V_+(F)\subset\P_\K^2$.
Umgekehrt ist jede glatte Kurve, die durch eine
solche Gleichung definiert wird, eine elliptische Kurve mit genanntem
Nullpunkt nach \cite[Proposition~III.3.1]{silverman}.
Gleichungen dieser Form werden Weierstraßgleichungen genannt.

Für $\Char(\K)\neq 2,3$ lässt sich eine Weierstraßgleichung durch
homogene Koordinatentransformation vereinfachen zur Form
\begin{gather*}
  Y^2 Z = X^3 + \beta XZ^2 + \gamma Z^3
  \quad\text{oder dehomogenisiert}\quad
  y^2 = x^3 + \beta x + \gamma\;,
\end{gather*}
wobei $\beta,\gamma\in\K$.
Für Berechnungen wird im Folgenden von dieser vereinfachten Form
ausgegangen. Für $\Char(\K)=2,3$ siehe
\cite[Appendix: Elliptic Curves in Characteristics 2~and 3]{silverman}.


% DISKRIMINANTE
Zu einer Weierstraßgleichung ist die Diskriminante und die
$j$"~Invariante assoziiert, welche für eine Gleichung in obiger
Form definiert sind als 
\begin{align*}
  \Delta &= -16\left(4\beta^3 + 27\gamma^2\right)
  &j &= (-16)\cdot 2^8\cdot 3^3\cdot\beta^3
\end{align*}
Eine Weierstraßgleichung ist genau dann regulär, definiert also eine
elliptische Kurve, wenn $\Delta\neq 0$ gilt,
und im singulären Fall gibt es genau einen singulären Punkt
nach \cite[Proposition~III.1.4]{silverman}.


\subsection{Weierstraßmodelle}
% (MINIMALE) WEIERSTRASSGLEICHUNG ÜBER $R$
Sei im folgenden Abschnitt $R$ Dedekindring,
$\K=\Quot(R)$ mit $\Char(\K)\neq2,3$ und $E_\K$ elliptische Kurve über
$\K$.

Es ist nun interessant, ob und wie sich eine Weierstraßgleichung mit
Koeffizienten in $R$, also ein projektives $R$"~Modell von
$E_\K$, finden lässt.
Weiterhin unter der Annahme ${\Char(\K)\neq2,3}$ sind die
Isomorphismen zwischen Weierstraßkurven zu $E_\K$ genau durch
Koordinatenwechsel der Form
\begin{gather*}
  [x,y,z] \longmapsto [a^{-2n}x, a^{-3n}y, z] \qquad a\in\K, n\in\Z
\end{gather*}
mit dem Koeffizientenwechsel
\begin{align*}
  \beta&\mapsto a^{4n}\beta
  &\gamma&\mapsto a^{6n}\gamma
  &\Delta&\mapsto a^{12n}\Delta
\end{align*}
gegeben, siehe \cite[Lemma~1.5/2]{neron}
oder \cite[Chapter~VII.1]{silverman}.
Sind ${\beta=\frac{\beta_1}{\beta_2}}$,
${\gamma=\frac{\gamma_1}{\gamma_2}}$ mit ${\beta_i,\gamma_i\in R}$,
erzeugt ein solcher Koordinatenwechsel mit ${a=\gamma_2\beta_2}$ eine
Weierstraßgleichung mit Koeffizienten in $R$.
\begin{Definition}[Weierstraßmodell]
  Im Folgenden wird ein projektives $R$"~Schema
  \begin{gather*}
    W=V_+(F)\subset\P_R^2\;,
  \end{gather*}
  das von einer homogenen
  Weierstraßgleichung $0=F\in R[X,Y,Z]$ von $E_\K$ mit Koeffizienten
  in $R\subset\K$ erzeugt wird, ein Weierstraßmodell von $E_\K$ über
  $R$ genannt.
\end{Definition}

Ist $R$ ein diskreter Bewertungsring, kann die Bewertung der
Diskriminante durch einen Koordinatenwechsel manipuliert werden. Daher
macht es Sinn, eine bezüglich der gegebenen diskreten Bewertung minimale
Form der Weierstraßgleichung in $R$ zu betrachten.
\begin{Definition}[minimale Weierstraßgleichung]
  \optcite[Chapter 1.5, S.\,22]{neron}
  \optcite[Chapter VII.1]{silverman}
  Sei $R$ ein diskreter Bewertungsring.
  Eine Weierstraßgleichung ($\Char(\K)\neq2,3$) von $E_\K$
  \begin{gather*}
    Y^2Z = X^3 + \beta XZ^2 + \gamma Z^3
  \end{gather*}
  heißt minimal, wenn für die Koeffizienten ${\beta, \gamma\in R}$ gilt
  und die Bewertung $\ord_\K(\Delta)$ der Diskriminante minimal
  bezüglich homogenem Koordinatenwechsel ist.
  Ein Weierstraßmodell von $E_\K$ über $R$ heißt minimal,
  wenn es von einer minimalen Weierstraßgleichung von $E_\K$ erzeugt
  wird.
\end{Definition}
Es stellt sich heraus, dass sich die spezielle Faser des
Weierstraßmodells einer elliptischen Kurve anhand der minimalen
Bewertung der Diskriminante und der $j$"~Invariante in eine von zehn
Kategorien klassifizieren lässt. Siehe dazu auch \autoref{chap:ausblick}. 

% -----

\section{Gruppenschemastruktur}\label{chap:weierstraßgruppe}
\subsection{Allgemeine Eigenschaften}
Im Folgenden werden ein paar nützliche Eigenschaften von
Weierstraßmodellen aufgeführt, die unter anderem für den Nachweis der
Gruppenstruktur auf dem glatten Teil eines Weierstraßmodells nötig
sind.

% INTEGER, PROJEKTIV, FLACH
\begin{Lemma}\label{thm:eigweierstrassmodelle}
  Ein Weierstraßmodell $W=V_+(F)$ zu einer Weierstraßgleichung ${0=F}$
  von $E_\K$ über einem Dedekindring $R$ mit $\K=\Quot(R)$ ist
  integer,
  projektiv, insbesondere quasikompakt und noethersch,
  endlich präsentiert, insbesondere von endlichem Typ, und
  flach.
  \begin{proof}
    Die meisten Eigenschaften sind direkt ersichtlich aus der Definition.
    \begin{description}
    \item[projektiv:] Per Konstruktion $W$ ist projektiv. Als
      abgeschlossenes Unterschema des quasikompakten, noetherschen Raums
      $\P_R^n$ ist es quasikompakt und noethersch.
    \item[integer:]
      Das Schema $W=V_+(F)$ ist genau dann irreduzibel, wenn
      ${F\in R[X,Y,Z]}$ irreduzibel ist.
      Da nach Voraussetzung $E_\K\cong V_+(F)\subset\P_\K^2$
      elliptische Kurve ist, vor allem irreduzibel, ist $F$ irreduzibel
      als Polynom in $\K[X,Y,Z]$. Nach Algebra impliziert das bereits
      die Irreduzibilität von $F\in R[X,Y,Z]$.
      % Dass $W$ irreduzibel ist, kann auf der affinen, offenen
      % Standardüberdeckung des $\P_R^2$ durch drei Kopien des $\A_R^2$
      % überprüft werden, da diese nicht-leeren Schnitt haben. Hier
      % erhalten wir die Kurven im $\A_R^2$, die durch die
      % dehomogenisierten Formen der Weierstraßgleichung $f(X,Y,Z)=0\in
      % R[X,Y,Z]$, d.\,h. $f_Z\in R[\frac{X}{Z},\frac{Y}{Z}]=R[x,y]$,
      % $f_Y\in R[x,z]$, $f_X\in R[y,z]$, erzeugt werden.
      % Diese stimmen nach Konstruktion mit den Erzeugern für die affinen
      % Kurven zu $E_\K$ überein und sind daher irreduzibel über $\K$,
      % also auch irreduzibel über $R$.
      % Entsprechend ist die gewählte affine, offene Überdeckung von $W$
      % irreduzibel und reduziert und somit $W$ integer.
    \item[endlich präsentiert:]
      $W$ ist offen überdeckt von seinen drei dehomogenisierten,
      affinen Teilmengen $V(F_i)=\Spec(R[x,y]/(F_i))$, wobei $F_X,F_Y,F_Z$
      Dehomogenisierungen von $F$ nach $X,Y,Z$ seien. Diese sind jeweils
      endlich präsentiert, also ist $W$ endlich präsentiert.
      Entsprechend folgt, dass $W$ von endlichem Typ ist.
    \item[flach:]
      Es ist zu zeigen, dass aus integer und nichtleerer generischer
      Faser $E_\K$ folgt, dass der eindeutige generische Punkt $\eta$
      von $W$ auf den generischen Punkt von $R$ geschickt wird.
      Das zeigt, dass $W$ flach ist
      nach \cite[Proposition~III.9.7]{hartshorne}, denn
      $\Spec(R)$ ist integer, regulär und von Dimension~1 und
      $W$ ist reduziert, womit alle Voraussetzungen des Satzes erfüllt
      sind.
      
      Nachdem $W$ integer ist, hat es einen eindeutigen generischen
      Punkt $\eta$.
      Angenommen, dieser wird durch die Abbildung, die $W$ zum
      $R$"~Schema macht, auf einen abgeschlossenen Punkt $\m$ von $\Spec(R)$
      geschickt.
      Dann liegt $\eta$ im abgeschlossenen Urbild von $\m$ –
      aber der Abschluss $\overline{\{\eta\}}=X$ von $\eta$ ist darin
      enthalten. Entsprechend wäre die generische Faser $E_\K$
      leer, was ein Widerspruch ist.
      Also muss der generische Punkt auf den generischen Punkt geschickt
      werden.
      \qedhere
    \end{description}
  \end{proof}
\end{Lemma}
% Für Regularität siehe auch \cite[Lemma IV.9.5]{silverman2}

% ÄQ GLATTHEIT
% \begin{Lemma}\label{thm:glattüberdvr}
%   Sei $R$ diskreter Bewertungsring,
%   $\K=\Quot(R)$,
%   $X$ irreduzibles,
%   lokal endlich präsentiertes
%   (für lokal noethersch äquivalent:
%   lokal endlichen Typs \cite[Remark 10.36]{wedhorn})
%   $R$"~Schema mit nichtleerer, generischer Faser $X_\K$.
%   Dann ist $X$ genau dann glattes $R$"~Schema, wenn die Fasern glatt
%   sind.
%   \begin{proof}
%     %     siehe auch \cite[Proposition IV.2.9]{silverman2}
%     Ist $X$ glatt, so auch jede Faser, da Glattheit stabil unter
%     Basiswechsel ist.

%     Seien also die Fasern von $X$ glatt, insb. reduziert.
%     Dann ist $X$ ebenfalls reduziert, da Reduziertheit auf den Halmen
%     geprüft werden kann.
%     Außerdem ist $X$ nach Voraussetzung irreduzibel, daher integer,
%     und (lokal) endlich präsentiert.
%     Wir wissen, dass $X$ glatt ist, wenn es lokal endlich präsentiert
%     und flach über $R$ ist und seine Fasern glatt sind
%     \cite[8.5, Proposition 17]{bosch}.
%     Es ist also Flachheit zu zeigen.

%     Wir werden zeigen, dass aus integer und nichtleerer generischer
%     Faser $X_\K$ folgt, dass der eindeutige generische Punkt $\eta$
%     von $X$ auf den generischen Punkt von $R$ geschickt wird.
%     Das zeigt, dass $X$ flach ist
%     \cite[Proposition III.9.7]{hartshorne}
%     (hierbei wird zusätzlich benutzt, dass $\Spec(R)$ integer,
%     regulär und von Dimension 1 ist).

%     Nachdem $X$ integer ist, hat es einen eindeutigen generischen
%     Punkt $\eta$.
%     Angenommen, dieser wird durch die Abbildung, die $X$ zum
%     $R$"~Schema macht, auf einen geschlossenen Punkt $\m$ von $\Spec(R)$
%     geschickt.
%     Dann liegt $\eta$ im geschlossenen Urbild von $\m$ –
%     aber der Abschluss $\overline{\{\eta\}}=X$ von $\eta$ ist darin
%     enthalten. Entsprechend wäre die generische Faser $X_\K$
%     leer, was ein Widerspruch ist.
%     Also muss der generische Punkt auf den generischen Punkt geschickt
%     werden.
%   \end{proof}
% \end{Lemma}

\begin{Satz}\label{thm:weierstrassglatt}
  \optcite[Remark IV.5.4.1]{silverman2}
  \optcite[Chapter VII.5 für die Definition von guter Reduktion]{silverman}
  % Gute Reduktion auf elliptischen Kurven heißt, dass die spezielle
  % Faser von $W$ regulär (also glatt also ebenfalls elliptische
  % Kurve) ist
  % \optcite[Chapter VII.5]{silverman}.
  Sei $(R,\m)$ ein diskreter Bewertungsring,
  ${\K=\Quot(R)}$,
  und sei der Restklassenkörper ${k=R/\m}$ perfekt.
  Sei weiter $E_\K$ eine elliptische Kurve über $\K$ und
  $W$ ein Weierstraßmodell von $E_\K$ über $R$
  mit spezieller Faser $W_\m$ über $k$.
  Dann gilt:
  \begin{enumerate}[label=(\alph*)]
  \item $W$ ist genau dann glatt, wenn die spezielle Faser von $W$ glatt
    ist.
  \item Ist $W$ nicht glatt, so ist der eindeutige, abgeschlossene,
    singuläre Punkt $\gamma$ der speziellen Faser der eindeutige
    Punkt, in dem $W$ nicht glatt ist, das heißt
    \begin{gather*}
      W^0\coloneqq W\setminus \{\gamma\}
    \end{gather*}
    ist das offene Unterschema der glatten Punkte von $W$.
    $W^0$ ist ein glattes $R$"~Modell von $E_\K$.
  \end{enumerate}
  \begin{proof}
    Hier sind die Ergebnisse aus \ref{thm:eigweierstrassmodelle}
    nötig, dass $W$ flaches, lokal endlich präsentiertes
    $R$"~Schema ist.
    \begin{enumerate}[resume*,start=1]
    \item 
      Nach \cite[Proposition~8.5/17]{bosch} ist $W$ genau dann
      glatt, wenn seine Fasern glatt sind.
      Nachdem $W_\K=E_\K$ nach Voraussetzung glatt ist, ist $W$ genau
      dann glatt, wenn $W_\m$ glatt ist.
    \item Für Schemata lokal endlichen Typs über einem Körper sind
      glatt und geometrisch regulär äquivalent
      \cite[Corollary~6.32]{wedhorn}.
      Da $k$ perfekt ist, ist dies für $k$"~Schemata sogar
      äquivalent zu regulär
      \cite[Remark~6.33]{wedhorn}.
      Schemata lokal endlichen Typs über einem Körper sind genau dann
      regulär, wenn sie in allen abgeschlossenen Punkten regulär sind
      \cite[Remark~6.25~(3)]{wedhorn}. Insbesondere sind alle
      singulären Punkte abgeschlossen.

      Als Weierstraßkurve ist $W_\m$ über $k$ ein $k$"~Schema
      lokal endlichen Typs, also treffen obige Aussagen zu.
      Nach \cite[Proposition~III.1.4]{silverman} ist unter den
      abgeschlossenen Punkten von $W_\m$ maximal ein singulärer Punkt
      $\gamma$. Demnach kann $W_\m$ maximal in einem abgeschlossenen
      Punkt $\gamma\in W_\m$ singulär, also nicht glatt, sein.

      Angenommen, $W_\m$ habe den singulären Punkt $\gamma$.
      $W_\m=W\setminus E_\K$ ist als Urbild des abgeschlossenen Punktes
      $\m$ von $\Spec(R)$ ein abgeschlossenes Unterschema von $W$.
      Also ist der Punkt $\gamma$, der abgeschlossen in $W_\m$ ist, auch
      abgeschlossener Punkt von $W$.
      
      Das Unterschema $W^0=W\setminus\gamma$ ist dementsprechend offen.
      Außerdem ist $W^0$ irreduzibel, da der Abschluss
      $\overline{W^0}=W$ in $W$ irreduzibel ist nach
      \ref{thm:eigweierstrassmodelle}.
      Es ist als Unterschema von $W$ ebenfalls reduziert, endlich
      präsentiert und flach, siehe \ref{thm:eigweierstrassmodelle}.
      Es hat wie $W$ die nichtleere, glatte, generische Faser $E_\K$ und
      per Konstruktion die spezielle Faser ${W_{\m}\setminus\{\gamma\}}$,
      welche nach Voraussetzung beide regulär über perfekten Körpern,
      also glatt, sind.
      Damit ist $W^0$ wieder nach \cite[Proposition~8.5/17]{bosch}
      glatt über $R$, heißt ein glattes $R$"~Modell von $E_\K$ und das
      offene Unterschema der glatten Punkte von $W$.
      \qedhere
    \end{enumerate}
  \end{proof}
\end{Satz}

\begin{Bemerkung}\label{thm:weierstraßmodellglatt}
  Für den Fall eines diskreten Bewertungsringes $(R,\m=(\pi))$ liefert
  \ref{thm:weierstrassglatt} ein sehr einfaches Kriterium für die
  Glattheit eines Weierstraßmodell $W$ von $E_\K$.
  $W$ ist genau dann glatt, wenn seine spezielle Faser glatt
  beziehungsweise regulär ist.
  Dies lässt sich recht leicht an der Weierstraßgleichung ablesen.
  
  Die spezielle Faser $W_{\m}=W\times_R (R/\m)$ ist definiert
  durch die Reduktion der Weierstraßgleichung von $W$ modulo dem
  Maximalideal $\m$ von $R$, also ebenfalls eine Weierstraßkurve.
  Entsprechend ist $W_{\m}$, und damit $W$, genau dann glatt, wenn
  eine der äquivalenten Bedingungen erfüllt ist:
  \begin{enumerate}[label=(\roman*)]
  \item $0\neq\Delta_{\m}=\bar\Delta\in R/\m$,
  \item $\Delta\not\in\m$ bzw. $\pi\nmid\Delta$,
  \item $\ord_R(\Delta)=0$ bzw. $\Delta\in R^\times$,
  \end{enumerate}
  wobei $\Delta_{\m}$ Diskriminante von $W_{\m}$ und $\Delta$
  Diskriminante von $W$ ist,
  \cite[vgl.][Proposition~VII.5.1~(a)]{silverman}.
  Die letzte Bedingung ist einfach nachzuprüfen.
\end{Bemerkung}

\subsection{Gruppenstruktur auf Weierstraßmodellen}
Das Folgende Lemma weist die Gruppenstruktur auf dem glatten Teil
eines Weierstraßmodells nach. Es wird nur eine kurze Beweisskizze
geliefert, da die Aussage nur ein Spezialfall von
\ref{thm:exneronmodelle} ist.
\begin{Lemma}\label{thm:erweiterunggruppenstruktur}
  Sei $(R,\m)$ wieder ein diskreter Bewertungsring
  mit perfektem Restklassenkörper ${k=R/\m}$ und
  ${\K=\Quot(R)}$.
  Sei $W$ ein Weierstraßmodell von $E_\K$ über $R$ und
  $W^0$ das offene $R$"~Unterschema der glatten Punkte von $W$,
  welches ebenfalls ein $R$"~Modell von $E_\K$ ist nach
  \ref{thm:weierstrassglatt}. 
  $W^0$ ist ein glattes $R$"~Gruppenschema
  mit Addition, Inversenbildung und neutralem Schnitt
  \begin{align*}
    \mu_0\colon W^0\times_R W^0 &\longto W^0
    &i_0\colon W^0&\longto W^0
    &\epsilon_0\colon R&\longto W^0
  \end{align*}
  welche die Gruppenstruktur (Addition $\mu$, Inversenbildung $i$ und
  neutraler Schnitt $\epsilon$)
  der abelschen Varietät $E_\K$ erweitern.
  \optcite[Theorem IV.5.3 (c)]{silverman2}

  % % ??? Referenzen
  % $W^0(R)\to W(R) \cong E_\K(\K)$ ist ein Gruppenisomorphismus,
  % falls $W$ regulär ist.
  % \optcite[Theorem IV.5.3 (b)]{silverman2}
  
  \begin{proof}[Beweisskizze] Sei wieder $\Char(\K)\neq 2,3$.
    \ref{thm:erweiterunggruppenstruktur} folgt als Spezialfall von
    \ref{thm:exneronmodelle}.
    Für das Nachweisen der Gruppenstruktur ist zu zeigen, dass die von
    $E_\K$ induzierten, rationalen Abbildungen $\mu\colon W\times_R
    W\dra W$ und $i\colon W\dra W$ auf $W^0$ definiert sind.

    $i$ ist in projektiven Koordinaten (also als graduierter
    Ringmorphismus der zugrundeliegenden graduierten Ringe)
    dargestellt als
    \begin{gather*}
      i\colon [X,Y,Z]\mapsto [X,-Y,Z]
    \end{gather*}
    Dies definiert einen Isomorphismus auf ganz $\P_R^2$, ist also
    insbesondere auf $W^0$ definiert.

    Für $\mu$ kann die zugehörige Darstellung als Morphismus
    graduierter Polynomringe über $\K$ auf einen Morphismus von
    Polynomringen über $R$ erweitert werden \cite[vgl. Formeln für
    Addition in][Group Law Algorithm~III.2.3]{silverman}.
    Berechnungen wie in \cite[Proposition~2.5]{silverman} zeigen, dass
    $\mu$ auf $W\times_R W$ höchstens im singulären Punkt nicht
    definiert ist, falls dieser existiert.
    Dies kann auf der affinen Überdeckung von $W$ nachgerechnet werden.    
    % \begin{enumerate}
    % \item $W^0=W\setminus\{\gamma\}$: \ref{thm:weierstrassglatt}
    %     %   \cite[Remark IV.5.4.2]{silverman2}
    %     %   Die nicht-glatten Punkte von $W$ sind die singulären Punkte
    %     %   der Fasern.
    %     %   Die generische Faser ist regulär (ell. Kurve), die spezielle
    %     %   Faser ist durch eine Weierstraßgleichung über $k$
    %     %   definiert. Sie hat also maximal einen singulären Punkt $\gamma$.
    %     %   \cite[Proposition III.1.4 (a)]{silverman}
    % \item $W^0$ ist $R$"~Modell von $E_\K$: \ref{thm:weierstrassglatt}
    %     %   $W$ ist $R$"~Modell von $E_\K$ per Definition,
    %     %   $W^0=W\setminus\{\gamma\}$ hat dieselbe generische Faser
    % \item Es treffen alle Voraussetzungen für
    %   \ref{thm:ratpkteregulaeremodelle} auf $W$ zu
    %   ($W$ ist projektiv, also eigentlich, von Dimension 2 mit
    %   projektiven Kurven als Fasern, und nach
    %   \ref{thm:weierstrassglatt} genau dann glatt in einem Punkt,
    %   wenn die Fasern dort regulär sind).
    %   Also gilt $E_\K(\K)=W(R)=W^0(R)$.
    %     %   \cite[Corollary IV.4.4]{silverman2}
    %     %   $E_\K(\K)=W(R)=W^0(R)$
    %     %   (s. \ref{thm:ratpkteregulaeremodelle}):
    %     %   \begin{enumerate}
    %     %   \item $E_\K(\K)=W(R)$: $W$ eigentlich
    %     %   \item $W(R)\subset W^0(R)$:
    %     %     nach \cite[Proposition IV.4.3 (b)]{silverman2}
    %     %   \item $W(R)=W^0(R)$: Alle Bilder von Schnitten $W(R)$ sind
    %     %     glatt
    %     %   \end{enumerate}
    % \item rationale Abbildungen
    %   $\mu\colon W\times_R W\dra W$,
    %   $i\colon W\dra W$:
    %   von generischer Faser $E_\K$ induziert, erfüllen Gruppenaxiome
    %   auf (größter, offener) Definitionsmenge (wo es Morphismus ist)
    % \item $\mu$ definiert auf $W^0$:
    %   \begin{enumerate}
    %   \item Aufteilen in affine Überdeckung
    %     ($W_Z=W\cap U_Z=\{[x,y,z]\in W| z\neq0\}$,
    %     $W_Y=W\cap U_Y=\{[x,y,z]\in W| y\neq0\}$)
    %     $W_Z\times W_Z$, $W_Z\times W_Y$, $W_Y\times W_Z$,
    %     $W_Y\times W_Y$
    %   \item $\mu$ auf $W_Z\times W_Z\setminus(\gamma,\gamma)$
    %     definiert:
    %     Schreibe rationale Abbildung von $\mu$ auf und finde
    %     Gleichungen für Unterschema, an denen der Nullpunkt getroffen
    %     würde
    %   \end{enumerate}
    % \item $i$ definiert auf $W^0$:
    %   $i\colon [X,Y,Z]\mapsto [X,-Y,Z]$ ist Morphismus in $\P_R^2$,
    %   also definiert auf ganz $W$
    % \end{enumerate}
  \end{proof}
\end{Lemma}

% IDENTITÄTSKOMPONENTE
\begin{Bemerkung}
  Für ein minimales Weierstraßmodell $W$ von $E_\K$ über $R$
  ist $W^0$ isomorph zur Identitätskomponente des Néron-Modells von
  $E_\K$ über $R$. Siehe dazu \cite[Chapter~1.5, S.\,23]{neron},
  \cite[Corollary~IV.9.1]{silverman2} oder \cite[S.~46]{tate}.
  Diese Tatsache ist eine direkte Folge der Konstruktion.
\end{Bemerkung}

% Im Folgenden sei $E_\K$ immer eine elliptische Kurve über dem
% Quotientenkörper $\K$ des Dedekindringes $R$.

Im glatten Fall ist das Néron-Modell wie versprochen bereits das
Weierstraßmodell mit der induzierten Gruppenstruktur.
\begin{Satz}\label{thm:neronmausweierstrassgl}
  \optcite[Corollary IV.6.3]{silverman2}
  Sei $R$ hier ein Dedekindring mit perfekten Restklassenkörpern,
  ${\K=\Quot(R)}$,
  $E_\K$ eine elliptische Kurve über $\K$ und
  $W$ ein Weierstraßmodell von $E_\K$ über $R$.
  Ist $W$ glatt, so ist $W$ bereits das Néron-Modell von $E_\K$ über
  $R$.
  \begin{proof}
    Nach \ref{thm:neronmodelllokal} ist nur zu prüfen, ob die 
    $W\times_R R_s$ für abgeschlossene Punkte $s\in\Spec(R)$
    Néron-Modelle von $E_\K$ über $R_s$ sind.
    Ist $W$ glatt, so ist auch $W\times_R R_s$ ein glattes
    Weierstraßmodell von $E_\K$ über $R_s$, da alles stabil unter
    Basiswechsel ist.
    Sei daher $R$ ohne Einschränkung bereits lokal, also ein diskreter
    Bewertungsring.

    Für die \NAbbEig ist nun von Nutzen, dass $W$ nach
    \ref{thm:erweiterunggruppenstruktur} ein $R_s$"~Gruppenschema ist.
    Sei $Y$ ein weiteres glattes $R$"~Schema mit generischer Faser
    $Y_\K$ und ${\phi_\K\colon Y_\K\to E_\K}$ ein Morphismus.
    In \ref{thm:erwrrational} unten wird gezeigt, dass $\phi_\K$ eine
    eindeutige $R$"~rationale Abbildung ${\phi\colon Y\dra W}$ liefert.
    Nachdem $W$ ein eigentliches $R$"~Gruppenschema und $Y$ glatt ist,
    ist $\phi$ nach \ref{thm:rationalzumorphismus} ein Morphismus, der
    $\phi_\K$ erweitert, und die \NAbbEig ist erfüllt.
  \end{proof}
\end{Satz}

\begin{Lemma}\label{thm:erwrrational}
  Sei $R$ ein Dedekindring, $\K=\Quot(R)$, $W$ allgemein ein
  eigentliches $R$"~Schema, $Y$ lokal noethersches, flaches, reguläres,
  integres
  $R$"~Schema und ${\phi_\K\colon Y_\K\to W_\K}$ ein Morphismus auf den
  generischen Fasern.
  Dann erweitert sich $\phi_\K$ eindeutig zu einer
  $R$"~rationalen Abbildung ${\phi\colon Y\dra W}$.
  Dasselbe gilt für ein glattes $Y$, da es dann als reguläres Schema
  die disjunkte Vereinigung seiner irreduziblen Komponenten ist.
  \begin{proof}
    $\Spec(R)$ und $Y$ sind nach Voraussetzung lokal noethersch,
    und der Morphismus ${\pi\colon Y\to R}$ ist flach über $R$. Daher
    gilt für ein $x\in Y$ nach \cite[Theorem~4.3.12]{liu}
    \begin{gather}\label{eq:stalkdimension}
      \dim\O_{Y,x} = \dim R_{\pi(x)} + \dim\O_{Y_{\phi(x)}, x}\;.
    \end{gather}
    Für $s\in\Spec(R)$ und den generischen Punkt $\eta_{s}$ einer
    irreduziblen Komponente der Faser $Y_s$ von $Y\to R$ gilt
    $\dim\O_{Y_s, \eta_{s}}=0$, also $\dim\O_{Y,\eta_{s}}=\dim R_s=1$.
    Da $Y$ regulär über $R$ ist, ist $\O_{Y,\eta_{s}}$ ein diskreter
    Bewertungsring.
    Der Quotientenkörper von $\O_{Y,\eta_{s}}$ ist der lokale Ring
    eines generischen Punkts $\eta$ von $Y$. Da $Y$ flach
    ist, liegt $\eta$ nach \cite[Proposition~III.9.7]{hartshorne} in
    der generischen Faser.
    Es kann also das Bewertungskriterium für Eigentlichkeit für $W$
    angewendet werden:
    \begin{center}
      \begin{tikzcd}
        \eta\arrow[r, equal]
        &\Spec(\k(\eta))
        \arrow[r, "{\phi_\K|_{\eta}}"] \arrow[d]
        & E_\K \arrow[r, hook]
        & W\arrow[d]\\
        \eta\cup\eta_{s}\arrow[r, equal]
        &\Spec(\O_{Y,\eta_{s}}) \arrow[rr]
        \arrow[urr, dashed, "{\exists!}", out=20, in=215]
        && \Spec(R)
      \end{tikzcd}
    \end{center}
    Damit erweitert sich $\phi_\K$ zu einem eindeutigen Morphismus
    ${\Spec(\O_{Y,\eta_{s}})\to W}$, der sich nach
    \ref{thm:morphismuserweiterung} und \ref{thm:erweindeutig}
    eindeutig auf einen Morphismus in einer Umgebung von $\eta_{s}$
    erweitert.
    
    Die so gewonnenen Morphismen in Umgebungen der generischen Punkte
    der Fasern können wiederum wegen der Eindeutigkeit aus
    \ref{thm:erweindeutig} verklebt. Das Ergebnis ist eine eindeutige
    $R$"~rationale Abbildung $\phi\colon U\dra W$ auf einem offenen
    Unterschema $U\subset Y$.
    Letzteres enthält alle generischen Punkte der Fasern von $Y\to R$
    und ist somit $R$"~dicht.
    $\phi$ ist demnach eindeutige $R$"~rationale Erweiterung von
    $\phi_\K$.
  \end{proof}
\end{Lemma}
\begin{Bemerkung}
  Im Fall diskreter Bewertungsringe kann die Glattheit anhand der Fasern
  überprüft werden:
  Aus \ref{thm:weierstraßmodellglatt} ist bekannt,
  dass ein Weierstraßmodell $W$ von $E_\K$ über $R$ genau dann glatt
  ist, wenn $\ord_R(\Delta)=0$ für die Diskriminante der
  Weierstraßgleichung gilt.
  Somit ist beispielsweise das Weierstraßmodell $W$ über $\Z_{(5)}$
  zur (hier dehomogenisierten) Weierstraßgleichung
  \begin{gather*}
    y^2=x^3+x
  \end{gather*}
  einer elliptischen Kurve $E_\Q$ über $\Q$ glatt, da
  $\Delta=-16\cdot4\cdot1^3$ eine Einheit in $\Z_{(5)}$ ist.
  $W$ ist deshalb bereits das Néron-Modell von $E_\Q$.

  Nach \ref{thm:abelscheneronmodelle} gilt auch, dass die
  Gruppenstruktur auf glatten Weierstraßmodellen die eines abelschen
  Schemas ist.
  Dies ist direkt aus der Konstruktion klar, da die Fasern des
  Weierstraßmodells in diesem Fall ja alle elliptische Kurven sind.
\end{Bemerkung}

%%% Local Variables:
%%% mode: latex
%%% TeX-master: "neron_models"
%%% End:


% Arithmetische Flächen
\chapter{Arithmetische Flächen}\label{chap:arithmetischeflächen}
Arithmetische Flächen sind grob gesagt besonders gutartige Kurven über
einem Dedekindring $R$, d.\,h. $R$-Schemata der Dimension~2.
Wichtig istbesonders, dass ihre Fasern alle projektive Kurven sind,
was sie leicht zu beschreiben macht.

In diesem Kapitel werden wir erst arithmetische Flächen einführen und
dann zeigen, dass es für elliptische Kurven immer eine solche Fläche
als $R$-Modell gibt. Für die Konstruktion wird das Weierstraßmodell
aus \autoref{chap:weierstraßmodelle} nötig sein.

Weiterhin werden einige Eigenschaften arithmetischer Flächen
untersucht. U.\,a. wird gezeigt, dass alle $R$-Schnitte einer
arithmetischen Fläche bereits $R$-Schnitte des glatten Teils
sind. Dies wird für die Konstruktion des Néron-Modells später von
Nutzen sein, welches sich als der glatte Teil einer minimal gewählten
arithmetischen Fläche entpuppen wird.

Diese Form der Minimalität wird in \autoref{chap:minmodelle}
eingeführt und die Existenz solcher minimalen Modelle für elliptische
Kurven gezeigt.

\section{Definition}
\begin{Definition}[arithmetische Fläche]
  Sei $R$ Dedekindring mit Quotientenkörper $\K$.
  % \optcite[IV.4]{silverman2}
  % Eine arithmetische Fläche $X$ über $R$ ist ein
  % integeres, normales, exzellentes Schema, flach und von endlichem
  % Typ über $R$,
  % dessen generische Faser eine reguläre, zusammenhängende,
  % projektive Kurve $X_\K/\K$ ist
  % und dessen spezielle Fasern Vereinigungen von Kurven über den
  % entsprechenden Restklassenkörpern sind.

  % % exzellent-Bediungung:
  % Hat $R$ Quotientenkörper der Charakteristik Null, so ist jedes
  % $R$-Schema lokal endlichen Typs bereits excellent.
  % \optcite[Theorem 12.51]{wedhorn}
  
  \optcite[Definition 8.3.1, 8.3.14]{liu}
  Eine arithmetische Fläche ist ein projektives, integeres, flaches,
  reguläres $R$-Schema von Dimension 2 (also Dimension 1 über $R$).
  Die generische Faser $X_\K$ ist integere Kurve über $\K$ und jede
  spezielle Faser $X_s$ ($s\in\Spec(R)$) ist projektive Kurve über
  $\k(s)$ nach \cite[Lemma 8.3.3]{liu}.
  
  $X$ ist also ein $X_\K$-Modell.
\end{Definition}

\begin{Beispiel}\label{bsp:regweierstraßmodell}
  Ein reguläres Weierstraßmodell einer elliptischen Kurve ist
  arithmetische Fläche.
  Beachte, dass $W$ regulär sein kann, aber dennoch nicht glatt, also
  singulär in einer Faser.
\end{Beispiel}

\begin{Satz}\label{thm:exarithfl}
  Für eine elliptische Kurve $E_\K$ über $\K$ existiert eine
  arithmetische Fläche über $R$, deren generische Faser $E_\K$ ist
  (kurz: eigentliches Modell von $E_\K$ über $R$).
  \begin{proof}
    Aus \autoref{chap:weierstraßmodelle} wissen wir, dass wir zu $E_\K$
    ein Weierstraßmodell $W$ über $R$ finden, d.\,h. ein projektives,
    integeres, flaches $R$-Schema, dessen generische Faser genau
    $E_\K$ ist. Zur arithmetischen Fläche fehlt $W$ also nur
    noch die Regularität.
    \cite[Corollary 8.3.51]{liu} besagt nun, dass $W$ eine starke,
    projektive Desingularisierung besitzt. D.\,h. es gibt ein
    projektives, reguläres $W$-Schema $\pi\colon Z\to W$, das
    birational und ein Isomorphismus auf den regulären Punkten von $W$
    ist.
    Birationale Abbildungen erhalten per Definition
    Irreduzibilitätskomponenten, Reduziertheit und Dimension.
    Daher ist $Z$ insbesondere wie $W$ integer
    und von Dimension~2.

    Weiterhin schicken birationale Morphismen generische Punkte auf
    generische Punkte, weshalb das reduzierte Schema $Z$ flach über
    $R$ ist nach \cite[Proposition III.9.7]{hartshorne}, weil
    $W$ flach über $R$ ist.

    Außerdem hat $Z$ dieselbe generische Faser $E_\K$, da
    $W$ in $E_\K$ glatt, insbesondere regulär ist.
    Insgesamt ist $Z$ arithmetischen Fläche über $R$ mit
    generischer Faser $E_\K$.
  \end{proof}
\end{Satz}

\begin{Bemerkung}\label{bem:konstruktioneigmodell}
  Die Desingularisierung von $W$ kann man induktiv durch eine
  Abfolge von abwechselnden Normalisierungen und Blow-Ups konstruieren:
  \begin{gather*}
    Z=Z_{n}\to Z_{n-1}'\to Z_{n-1}\to\dotsb
    \to Z_{2}\to Z_{1}'\to Z_{1}\to Z_{0}'=W
  \end{gather*}
  Die $Z_i\to Z_{i-1}'$ sind die Normalisierung der $Z_{i-1}'$,
  d.\,h. birationale, integere Morphismen mit normalen $Z_{i}$ (welche
  wegen Birationalität auch wieder integer sind). Sie existieren nach
  \cite[Proposition 4.1.22]{liu}.
  
  Die $Z_{i}'\to Z_{i}$ sind Blow-ups entlang des singulären Teils von
  $Z_{i}$, der in unserem Fall jeweils geschlossen ist.
  Zu Blow-Ups siehe auch
  \cite[Chapter 8.1]{liu} oder \cite[Chapter 13.19]{wedhorn}.
  Wichtig hier ist, dass Blow-Ups außerhalb des aufgeblähten Bereichs
  Isomorphismen sind, also insbesondere unsere generische Faser und
  die Eigenschaft integral erhalten.

  \cite[Corollary 8.3.51]{liu} besagt, dass diese Folge in endlich
  vielen Schritten ein reguläres $Z_{n}$ liefert.
\end{Bemerkung}

% -----
\section{Eigenschaften}
Eine recht praktische Eigenschaft von arithmetischen Flächen ist, dass
sie in den Bildern all ihrer Schnitte regulär sind.
Dieser Abschnitt befasst sich mit dieser Aussage sowie einem
zugehörigen Korollar.

\begin{Lemma}\label{thm:arithflschnittbilderglatt}
  \optcite[Proposition IV.4.3]{silverman2}
  Sei $R$ Dedekindring,
  $\pi\colon X\to\Spec(R)$ eine arithmetische Fläche über $R$,
  $\p\in\Spec(R)$ und $\phi\in X(R)$.
  Dann ist $X_\p$ regulär in $\phi(\p)$.
  \begin{proof}
    Nachdem die generische Faser einer arithmetischen Fläche bereits
    glatt, also regulär ist, sei $\p\subset\Spec(R)$ ein geschlossener
    Punkt.
    
    Sei $x=\phi(\p)\in X_\p\subset X$ und $\m_{X_\p,x}\subset\O_{X_\p,x}$
    das Maximalideal des Halmes von $x$ in $X_\p$,
    analog $\m_{X,x}\subset\O_{X,x}$.
    Der Punkt $x\in X_\p(\k(\p))\subset X_\p$ ist geschlossen in $X_\p$,
    daher gilt nach \cite[Corollary 2.5.24]{liu}
    $\dim\O_{X_\p,x}=\dim X_\p=1$ ($X_p$ ist equidimensional nach
    \cite[Proposition 4.4.16]{liu} und projektive Kurve).
    Wir müssen für die Regularität von $X_\p$ in $x$ also zeigen, dass
    $\m_{X_\p,x}/\m_{X_\p,x}^2$ als $(\O_{X_p,x}/\m_{X_p,x})$-Vektorraum
    von einem Element $f_1$ erzeugt wird.
    
    $\O_{X,x}$ ist $R$-Algebra über den induzierten lokalen
    Ringhomomorphism $\pi^\#\colon R_\p\to\O_{X,x}$.
    Sei $\phi^\#\colon\O_{X,x}\to R_\p$ der induzierte lokale
    Ringhomomorphismus zu $\phi$ auf den Halmen, wobei
    $\phi^\#(\m_{X,x})=\p$ gelten muss, da $\pi^\#\circ\phi^\#(\p)=\p$.
    
    Wir werden zeigen, dass es ein Element
    $t\in\pi^\#(\p)\setminus\m_{X,x}^2$ gibt und
    daraus die gewünschte Dimension von $\m_{X_\p,x}/\m_{X_\p,x}^2$
    bzw. die Regularität von $X_\p$ in $x$ folgern.

    Angenommen, es gelte $\pi^\#(\p)\subset\m_{X,x}^2$.
    Mit $\pi\circ\phi=\Id_{\Spec(R)}$ erhalten wir die Abschätzung
    \begin{gather*}
      \p = (\pi\circ\phi)^\#(\p)
      = \phi^\#\circ\pi^\#(\p)
      \overset{\text{Ann.}}{\subset} \phi^\#(\m_{X,x}^2)
      = \phi^\#(\m_{X,x})^2
      = \p^2
    \end{gather*}
    Dies ist ein Widerspruch, da $\p\subset R_\p$ Maximalideal eines
    diskreten Bewertungsringes ist und deshalb wegen Regularität von
    $R_\p$ gelten muss $\p/\p^2\neq 0$.
    Also gilt $\pi^\#(\p)\not\subset\m_{X,x}^2$.

    Nach \ref{eq:stalkdimension} gilt
    $\dim\O_{X,x}=\dim R_\p+\dim\O_{X_\p,x}=2$ und $X$ ist nach
    Voraussetzung regulär, weshalb $\O_{X,x}$ regulär von Dimension~2
    ist. Das bedeutet, dass es zwei Elemente $f_1, f_2\in\m_{X,x}$
    gibt, so dass
    \begin{gather}\label{eq:maximalideal}
      \m_{X,x} = f_1\O_{X,x} + f_2\O_{X,x} + \m_{X,x}^2
    \end{gather}
    Wegen $\pi^\#(\p)\not\subset\m_{X,x}^2$ gibt es ein
    $t\in\pi^\#(\p)\setminus\m_{X,x}^2\subset\m_{X,x}$.
    \eqref{eq:maximalideal} bedeutet für $t$, dass es $\alpha_1,
    \alpha_2\in\O_{X,x}$ gibt, so dass
    \begin{gather*}
      t = \alpha_1 f_1 + \alpha_2 f_2 \mod \m_{X,x}^2
    \end{gather*}
    Da $t\not\in\m_{X,x}^2$, muss mindestens einer der Koeffizienten
    in $\O_{X,x}\setminus\m_{X,x}=\O_{X,x}^\times$ liegen, also einen
    Einheit sein. \OE{} $\alpha_2\in\O_{X,x}^\times$, dann können wir
    umschreiben
    \begin{gather*}
      f_2 = \frac{1}{\alpha_2}t + \frac{\alpha_1}{\alpha_2}f_1 
    \end{gather*}
    und aus \eqref{eq:maximalideal} wird
    \begin{gather}\label{eq:maximalidealneu}
      \m_{X,x} = t\O_{X,x} + f_1\O_{X,x} + \m_{X,x}^2
    \end{gather}

    Für den lokalen Ring $\O_{X_\p,x}$ von $x$ in der generischen
    Faser $X_\p = X\times_R (R/\p)$ gilt
    \begin{align*}
      \O_{X_\p,x} &= \O_{X,x}/\pi^\#(\p)\\
      \m_{X_\p,x} &= \m_{X,x}/\pi^\#(\p)
                    \overset{\eqref{eq:maximalidealneu}}{=}
                    \bar f_1\O_{X_\p,x} + \m_{X_\p,x}^2
    \end{align*}
    Wir sehen, $\m_{X_\p,x}$ ist von $\m_{X_\p,x}^2$ und einem Element
    erzeugt, was $\O_{X_\p,x}$ zu einem regulären Ring der Dimension~1
    macht.
  \end{proof}
\end{Lemma}

\begin{Korollar}\label{thm:ratpkteregulaeremodelle}
  \optcite[Corollary IV.4.4]{silverman2}
  Sei $R$ Dedekindring
  mit perfekten Restklassenkörpern,
  $\K=\Quot(R)$,
  $X$ eine arithmetische Fläche,
  $X^0$ offenes Unterschema der glatten Punkte,
  $X_\K/\K$ generische Faser.
  Dann gilt
  \begin{gather*}
    X_\K(\K)=X(R)=X^0(R)
  \end{gather*}
  % $X$ ein $R$-Schema,
  % $X^0$ offenes Unterschema der glatten Punkte,
  % $X_\K/\K$ generische Faser.
  % Dann gilt
  % \begin{enumerate}[label=(\roman*)]
  % \item Ist $X$ eigentlich über $R$, so gilt $X(R)=X_\K(\K)$.
  % \item Ist $X$ regulär über $R$ mit Dimension 2 und mit
  %   Fasern der Dimension 1, und ist $X$ genau dann glatt in einem
  %   Punkt $x\in X$, wenn die ihn enthaltene Faser regulär in $x$ ist,
  %   so gilt $X(R)=X^0(R)$.
  % \end{enumerate}

  % Es gilt also für ein reguläres, eigentliches $R$-Modell von $X_K$,
  % das genau dann glatt in $x\in X$ ist, wenn die ihn enthaltene
  % Faser in $x$ regulär ist,
  % \begin{gather*}
  %   X_\K(\K)=X^0(R)=X(R)
  % \end{gather*}
  \begin{proof}
    Die erste Gleichheit ist allein eine Anwendung des
    Bewertungskriteriums für Eigentlichkeit auf das projektive (also
    eigentliche) Schema $X$.
    % \begin{enumerate}[label=(\roman*)]
    % \item
    %   Die Einschränkung von $\Hom_R(\Spec(R),X)=X(R)$ auf 
    %   $\Spec(\K)\subset\Spec(R)$ liefert eine kanonische Abbildung
    %   \begin{center}
    %     \begin{tikzcd}
    %       \Hom_R(\Spec(R),X) \arrow[d, equal]
    %       \arrow[r, "{\text{Res.}}"]
    %       &\Hom_R(\Spec(K),X)
    %       \arrow[r, equal, "{\text{$R$-lin.}}"]
    %       &\Hom_\K(\Spec(\K),X_\K) \arrow[d, equal]\\
    %       X(R)&& X_\K(\K)
    %     \end{tikzcd}
    %   \end{center}
    %   Nach Voraussetzung ist $X$ eigentlich über $R$, d.\,h. nach
    %   dem Bewertungskriterium für $X$ gibt es für ein $\phi_\K\in
    %   X_\K(\K)$ genau einen Lift in $\Hom_R(R,X)$,
    %   bzw. ein $\phi\colon\Spec(R)\to X$, so dass das folgende
    %   Diagramm kommutiert
    %   \cite[Theorem II.4.7]{hartshorne}:
    %   \begin{center}
    %     \begin{tikzcd}
    %       \Spec(\K)\arrow[d, hook, "{\text{incl}}"]
    %       \arrow[r,"\phi_\K"]
    %       &X_\K\arrow[r, hook, "{\text{incl}}"]
    %       &X\arrow[d,"{\text{eig.}}"]
    %       \\
    %       \Spec(R)\arrow[rr,"{\Id}"]
    %       \arrow[to=urr, dashed, "{\exists!\phi}"{near start}]
    %       &&\Spec(R)
    %     \end{tikzcd}
    %   \end{center}
    %   Dadurch sind Injektivität und Surjektivität der
    %   Einschränkungsabbildung gegeben und es gilt
    %   $X(R)=\Hom_R(R,X)\cong\Hom_\K(\K,X_\K)=X_\K(\K)$.
    % \item
    Da $X^0\subset X$, gibt es eine natürliche Inklusion
    $\circ \text{incl}\colon X^0(R)\hookrightarrow X(R)$.
    Jetzt gilt es zu zeigen, dass im regulären Fall das Bild jedes
    Schnitts $\phi\in X(R)$ in $X^0$ liegt (also nur glatte Punkte
    enthält). Dann ist die Inklusion ein Isomorphismus.
    
    Nach \ref{thm:arithflschnittbilderglatt} enthält das Bild
    $\phi(\Spec(R))$ eines Schnitts $\phi\in X(R)$ nur reguläre
    Punkte der Fasern. Diese sind glatte Punkte von $X$, denn:
    $X$ ist flach und lokal endlich präsentiert, daher genau dann
    glatt, wenn die Fasern glatt sind nach
    \cite[8.5, Proposition 17]{bosch}.
    Die generische Faser ist bereits glatt und die speziellen Fasern
    sind als Schemata endlichen Typs über einem perfektem Körper
    genau dann glatt, wenn sie regulär sind.
    Damit ist die Behauptung gezeigt.
    % \end{enumerate}
  \end{proof}
\end{Korollar}

% -----

\section{Minimale Modelle}\label{chap:minmodelle}
Eigentliche Modelle können für elliptische Kurven minimal
bzgl. birationaler Morphismen gewählt werden, wie wir in diesem
Abschnitt zeigen werden. Aus der Konstruktion wird auch ersichtlich,
dass der Algorithmus aus \ref{bem:konstruktioneigmodell} bereits ein
minimales Modell liefert.

Ein solches minimales Modell ist die Grundlage für unseren
Existenzbeweis des Néron-Modells, welches genau der glatte Teil des
minimalen Modells ist.

\begin{Definition}
  Sei $R$ Dedekindring mit perfekten Restklassenkörpern,
  $\K=\Quot(R)$,
  $C_\K$ reguläre, projektive Kurve über $\K$.
  Ein minimales eigentliches Modell $X_{\text{min}}$ für $C_\K$ über
  $R$ (kurz: minimales Modell von $C_\K$ über $R$) ist ein
  eigentliches $R$-Modell von $C_\K$, das folgende
  Minimalitätseigenschaft erfüllt:
  \begin{quote}
    Für jedes andere eigentliche Modell $X$ von
    $C_\K$ über $R$ und für jeden $\K$-Isomorphismus $\phi_\K$ der
    generischen Fasern ist die nach \ref{thm:erwrrational}
    induzierte $R$-birationale Abbildung $\phi\colon X\dra
    X_{\text{min}}$ ein $R$-Morphismus.
  \end{quote}
\end{Definition}
% % exzellent oder glatt nötig!!! evtl. nur Konstruktion?
% \begin{Definition}
%   Sei $R$ Dedekindring mit perfekten Restklassenkörpern,
%   $\K=\Quot(R)$,
%   $C/\K$ reguläre, projektive Kurve vom Geschlecht $g$.
%   \begin{enumerate}[label=(\roman*)]
%   \item\cite[Proposition IV.4.5(a)]{silverman2}
%     Dann existiert ein eigentliches $R$-Modell $X$ von $C$,
%     das eine arithmetische Fläche darstellt.
%   \item\cite[Proposition IV.4.5(b)]{silverman2} Ist $g\geq1$ kann
%     ein minimales eigentliches Modell $X_{\text{min}}$ für $C$ über
%     $R$ gewählt werden (kurz: minimales Modell von $C$ über $R$).
%     D.\,h. für jedes andere eigentliche Modell $X$ von
%     $C/\K$ über $R$ und für jeden $\K$-Isomorphismus $\phi_\K$ der
%     generischen Fasern ist die nach \ref{thm:erwrrational}
%     induzierte $R$-birationale Abbildung $\phi\colon X\dra
%     X_{\text{min}}$ ein $R$-Morphismus.
%   \end{enumerate}
% \end{Definition}

\begin{Satz}\label{thm:eigminmodell}
  Seien $R$, $\K$, $C_\K$ wie oben.
  Ein minimales Modell $X_{\text{min}}$ von $C_\K$ über $R$ hat
  folgende Eigenschaften:
  \begin{enumerate}[label=(\roman*)]
  \item $X_{\text{min}}$ ist eindeutig bis auf eindeutige Isomorphie.
  \item\optcite[Proposition IV.4.6]{silverman2}
    Für ein minimales Modell $X_{\text{min}}$ von $C_\K$ über $R$
    gilt, dass jeder $\K$"=Automorphismus von $C_\K$ sich zu
    einem eindeutigen $R$-Automorphismus von $X_{\text{min}}$
    erweitert, der glatte Punkte auf glatte Punkte abbildet.
  \item\optcite[Beweis von Proposition IV.6.10]{silverman2}
    Minimale Modelle sind wie folgt stabil unter Basiswechsel:
    Sei $Y$ glattes $R$-Schema, $x\in Y$ und der Halm $R'=\O_{Y,x}$ in
    $x$  ein diskreter Bewertungsring mit $\K'\coloneqq\Quot(R')$.
    Dann ist das minimale Modell von $C_\K\times_\K\K'$ genau
    $X_{\text{min}}\times_R R'$.
    Ist $E\subset X_{\text{min}}$ das Unterschema der glatten Punkte
    von $X_{\text{min}}$, ist
    $E\times_R R'\subset X_{\text{min}}\times_R R'$ das Unterschema
    der glatten Punkte von $X_{\text{min}}\times_R R'$.
  \end{enumerate}
  \begin{proof}
    \begin{enumerate}[label=(\roman*)]
    \item Für zwei minimale Modelle $X_1$, $X_2$ von $C_\K$ über $R$
      liefert die Minimalitätseigenschaft zwei $R$-Morphismen
      $X_1\to X_2$, $X_2\to X_1$, deren Komposition jeweils die
      Identität auf $C_\K$ erweitert. Wegen der Eindeutigkeit aus
      \ref{thm:erweindeutig} müssen die Kompositionen jeweils die
      Identität sein und die Morphismen sind zueinander inverse
      $R$-Isomorphismen.
    \item Die Minimalitätseigenschaft liefert zu einem
      $\K$-Automorphismus $\tau\colon C_\K\to C_\K$ die Erweiterung auf
      einen $R$-Automorphismus
      $\tau\colon X_{\text{min}}\to X_{\text{min}}$, der wegen
      Eigentlichkeit eindeutig ist.
      Als bijektive, also offene Abbildungen, schicken $\tau$ und
      sein Inverses offene glatte Umgebungen auf offene glatte
      Umgebungen, und somit glatte Punkte auf glatte Punkte. Wir
      erhalten also einen $R$-Automorphismus auf dem Unterschema der
      glatten Punkte.
    \item Die Eigenschaften flach, projektiv und integer sind stabil
      unter Basiswechsel. Regulär ist stabil unter glattem
      Basiswechsel \cite[Chapter 2.3, Page 49]{neron}, weshalb
      $C_\K\times_R Y$ regulär ist.
      Nachdem $C_\K\times_R R'$ dieselben Halme hat, ist es ebenfalls
      regulär.
      Die Minimalitätseigenschaft von $X_{\text{min}}\times_R R'$
      folgt aus der von $X_{\text{min}}$ und der universellen
      Eigenschaft des Faserprodukts.
      
      $E\times_R R'$ ist glattes Unterschema von
      $X_{\text{min}}\times_R R'$, da glatt stabil unter Basiswechsel
      ist. Nachdem $Y$ glatt ist, ist ein Punkt in
      $X_{\text{min}}\times_R Y$ und insbesondere in
      $X_{\text{min}}\times_R R'\subset
      X_{\text{min}}\times_R Y$ genau dann glatt, wenn seine
      Projektion auf $X_{\text{min}}$ glatt ist, also wenn sie in $E$
      liegt. Damit ist $E\times_R R'$ das Unterschema aller glatten
      Punkte in $X_{\text{min}}\times_R R'$.
    \end{enumerate}
  \end{proof}
\end{Satz}

\begin{Satz}[{\cite[Theorem 9.3.21]{liu}}]\label{thm:exminmodell}
  Minimale $R$-Modelle elliptischer Kurven existieren und sind nach
  \ref{thm:eigminmodell} eindeutig bis auf eindeutige Isomorphie.
  \begin{proof}[Beweisskizze]
    Aus \ref{thm:exarithfl} wissen wir, dass wir zu einer elliptischen
    Kurve $E_\K$ über $\K$ zumindest eine arithmetische Fläche finden.
    \cite[Theorem 9.3.21]{liu} besagt, dass allgemein für
    arithmetische Flächen mit generischer Faser vom Genus größer~0 ein
    minimales Modell existiert.
    Genutzt wird hierbei, dass die Existenz eines minimalen Modells
    äquivalent dazu ist, dass folgende Bedingungen erfüllt sind:
    \begin{enumerate}[label=(\alph*)]
    \item Es existiert ein \emph{relativ minimales Modell}, d.\,h. ein
      $R$-Modell mit der abgeschwächten Minimalitätseigenschaft:
      Jeder birationale Morphismus in ein anderes $R$-Modell ist
      Isomorphismus.
    \item Alle relativ minimalen Modelle sind isomorph, es gibt also
      ein eindeutiges relativ minimales Modell.
    \end{enumerate}
    Relative Minimalität ist nach \cite[Theorem 9.2.2]{liu} äquivalent
    dazu, dass das Modell keine außerordentlichen Divisoren besitzt,
    also keine integralen, zusammenziehbaren Kurven als Teilmengen
    hat.
    Außerordentliche Divisoren kann man u.\,a. mit Castelnuovos
    Kriterium \cite[Theorem 9.3.8]{liu} identifizieren und es gibt nur
    eine endliche Anzahl an Fasern, die welche enthalten nach
    \cite[Lemma 9.3.17]{liu}.

    Nach \cite[Proposition 9.3.19]{liu} ist es nun möglich, die
    zusammenziehbaren Divisoren in endlich vielen Schritten
    nacheinander zusammenzuziehen zu einem dann relativ minimalen
    Modell.

    Die zweite Bedingung kann wie in \cite[Theorem 9.3.21]{liu}
    beschrieben durch Widerspruchsbeweis mithilfe der Annahme zum
    Genus der generischen Faser gezeigt werden.
    Wir haben also durch das Zusammenziehen bereits ein minimales
    Modell konstruiert.
  \end{proof}
\end{Satz}

\begin{Bemerkung}\label{bem:konstruktionminmodell}
  Wählt man in \ref{thm:exarithfl} das $n$ minimal, erhält man bereits
  ein relativ minimales, also minimales, Modell. Denn
  bei jedem Zusammenziehen eines Divisors würde per Konstruktion die
  Regularität des Modells verloren gehen.
\end{Bemerkung}

%%% Local Variables:
%%% mode: latex
%%% TeX-master: "neron_models"
%%% End:


% Existenz von Néron-Modellen auf elliptischen Kurven
\chapter{Existenz von Néron-Modellen auf elliptischen Kurven}
\label{chap:exneronmodelle}
% EXISTENZSATZ ELL. KURVEN
\begin{Satz}\label{thm:exneronmodelle}
  % Der diskrete Bewertungsring $R$ sei strikt henselsch mit algebraisch
  % abgeschlossenem Quotientenkörper $\K$ und
  % $W$ ein minimales reguläres, eigentliches, flaches $R$-Modell der
  % elliptischen Kurve $E_\K$.
  % Dann ist das $R$-Unterschema der glatten Punkte von $W$ ein
  % Néron-Modell von $E_\K$.
  % \cite[1.5, Proposition 1, S.\,21]{neron}
  % \begin{proof}
  %   …
  % \end{proof}

  % Fehlt: \cite[Theorem IV.6.1]{silverman2}
  % (ähnlich, ohne Bed. an den Ring, dafür mehr an das Modell)
  Sei $R$ Dedekindring mit perfekten Quotientenkörpern,
  $\K=\Quot(R)$, $E_\K$ elliptische Kurve über $\K$ und $C$ minimales
  $R$-Modell von $E_\K$,
  das heißt reguläres, integeres, projektives,
  flaches $R$-Schema mit Minimalitätseigenschaft.
  Sei $E$ das offene Unterschema der glatten Punkte von $C$.
  Dann ist $E$ Néron-Modell von $E_\K$.

  Nachdem minimale Modelle elliptischer
  Kurven nach \ref{thm:exminmodell} existieren, hat jede
  elliptische Kurve über einem Dedekindring mit perfekten
  Restklassenkörpern ein Néron-Modell.
\end{Satz}
Dieses Kapitel ist dem Beweis des vorangegangenen Satzes gewidmet.

Von der Idee her werden wir uns vorerst in
\autoref{chap:gruppenstrukturhinreichend} klar machen, dass $E$ genau 
dann ein Néron-Modell von $E_\K$ ist, wenn sich die Gruppenverknüpfung
von $E_\K$ auf ganz $E$ erweitert (das \enquote{genau dann} ist
bereits aus \ref{thm:gruppenschemaerweiterung} bekannt).

Nun bleibt es zu zeigen, dass sich die Gruppenverknüpfung tatsächlich
auf $E$ erweitert.
Wir betrachten dies in \autoref{chap:gruppestrikthenselsch} vorerst für
den Spezialfall eines strikt henselschen diskreten Bewertungsrings
$R$. Dabei sehen wir, dass man eine normale Komposition auf $E$
erhält. Nach \ref{thm:weil} wissen wir nun, dass $E$ birational zu
einem $R$-Gruppenschema mit der gewünschten Struktur ist. Für den
Abschluss des Beweises stellen wir fest, dass
diese $R$-birationale Abbildung bereits ein $R$-Isomorphismus und $E$
daher ein $R$-Gruppenschema ist.

In \autoref{chap:exneronmodellebeweis} zeigen wir zuletzt, dass wir
uns ohne Einschränkung auf strikt henselsche Bewertungsringe
beschränken können, was dann mit den vorangegangenen Ergebnissen den
Beweis vollendet.


% \begin{proof}
  %   \begin{enumerate}
  %   \item \ref{thm:rationalzumorphismus}:
  %     Eine rationale Abbildung $X\dra G$ von einem glatten
  %     $R$-Schema in ein eigentliches $R$-Gruppenschema ist Morphismus.
  %   \item \ref{def:henselscheringe}: (strikt) Henselsche Ringe
  %   \item Eigenschaften henselscher Ringe auf Schnitteinschränkung:
  %     $(R,\m)$ diskreter Bewertungsring mit Restklassenkörper $k$, $X$
  %     glattes $R$-Schema, $X_k$ spezielle Faser.
  %     Dann
  %     \begin{enumerate}[label=(\roman*)]
  %     \item Falls $R$ henselsch, ist $X(R)\to X_k(k)$ surjektiv.
  %     \item Falls $R$ strikt henselsch, ist $X(R)\to X_k(k)$ surjektiv
  %       und hat dichtes Bild in $X_k$.
  %       (Beachte, dass $X_k$ glatt, also von endlichem Typ ist und $k$
  %       nach Voraussetzung algebraisch abgeschlossen ist. Daher
  %       gilt $X_k(k)\cong\{x\in X_k\text{ geschl.}\}\subset X_k$
  %       \cite[Example 4.1]{wedhorn}.)
  %     \end{enumerate}
  %     \cite[Proposition IV.6.4]{silverman2},
  %     \cite[2.3, Proposition 5]{neron}
  %   \item Henselisierung: \ref{thm:exhenselisierung}
  %   \item Fall $R$ strikt henselscher, diskreter BWR:
  %     \begin{enumerate}
  %     \item Ist $E$ Gruppenschema, so ist es Néron-Modell
  %       \ref{thm:fallstriktehenselisierung}
  %     \item Def. Normale Komposition \ref{def:normalekomposition}
  %     \item Für $R$ strikter henselscher, diskreter BWR
  %       definiert $E_\K$ eine normale Komposition auf $E$.
  %     \item (Theorem von Weil) Ein glattes $R$-Schema mit normaler
  %       Komposition ist birational zu einem $R$-Gruppenschema endlichen
  %       Typs: \ref{thm:weil}
  %     \item Also ist $E$ birational zu einem Gruppenschema $G$.
  %       Die rationalen Abbildungen sind Morphismen.
  %       \ref{thm:egruppenschema}
  %     \end{enumerate}
  %   \item Fall diskreter BWR:
  %     \begin{enumerate}
  %     \item Minimale eigentliche, reguläre Modelle sind stabil unter
  %       flachem, unverzweigtem Basiswechsel.
  %     \item $E^\text{sh}=E\times_R\Rsh$ ist glatter Teil von
  %       $C^\text{sh}$, also Néron-Modell von $E_\K^\text{sh}$.
  %     \item $E$ ist Néron-Modell wegen „faithfully flat descent“
  %     \end{enumerate}
  %   \item Allg. Dedekindring: lokal (\ref{thm:neronmodelllokal})
  %   \end{enumerate}
  % \end{proof}

  
\section{Gruppenstruktur hinreichend}\label{chap:gruppenstrukturhinreichend}
% ÄQ GRUPPENSTRUKTURERW.
% strikt henselsch nicht nötig!!
Die Hauptaussage dieses Kapitels wird sein, dass $E$ genau dann
Néron-Modell von $E_\K$ ist, wenn sich die Gruppenverknüpfung von
$E_\K$ auf ganz $E$ erweitert und $E$ damit zum $R$-Gruppenschema
macht.
\begin{Lemma}\label{thm:schnittearithmflglatt}
  Sei $R$ diskreter Bewertungsring mit perfektem Restklassenkörper
  $k$, $\K$, $E_\K$, $C$, $E$ wie in \ref{thm:exneronmodelle}.
  Sei $(R',\m')$ der Halm eines glatten $R$-Schemas und diskreter
  Bewertungsring
  und sei ein $R$-Morphismus $\phi\colon \Spec(R')\to C$ gegeben.
  Ist $\phi(\Spec(\K))\in E$, gilt $\im(\phi)\subset E$.
  \begin{proof}
    Für den generischen Punkt $\phi((0))$ ist dies klar, da die
    generische Faser von $C$ glatt ist und nach
    \cite[8.5, Proposition 17]{bosch} wegen Flachheit von $C$ in $E$
    liegt. 
    Betrachte also $\phi(\m')$ für den geschlossenen Punkt
    $\m'\in\Spec(R')$.
    $C_{R'}\coloneqq C\times_R R'$ ist nach
    \ref{thm:eigminmodell} wieder eine arithmetische Fläche über
    $R'$ mit Unterschema der glatten Punkte
    $E_{R'}\coloneqq E\times_R R'$.
    Nun ist der kanonische Morphismus
    $(\phi\times_{R}\Id_{R'})\colon\Spec(R')\to C_{R'}$ ein 
    $R'$-Schnitt, weshalb die spezielle Faser $(C_{R'})_{\m'}$ nach
    \ref{thm:arithflschnittbilderglatt} regulär in
    $(\phi\times\Id)(\m')$ ist.
    Als Schema endlichen Typs über einem
    perfekten Körper ist $(C_{R'})_{\m'}$ sogar glatt in diesem
    regulären Punkt.
    Da $C_{R'}$ als flaches, endlich präsentiertes $R'$-Schema wieder
    genau dann glatt in einem Punkt ist, wenn es die zugehörige Faser
    ist, ist $C_{R'}$ glatt in $(\phi\times\Id)(\m')$ bzw.
    $(\phi\times\Id)(\m')\in E\times_R R'$. Folglich muss bereits
    $\phi(\m')$ in $E$ liegen und $\phi\colon \Spec(R')\to E\subset C$
    faktorisiert über $E$.
  \end{proof}
\end{Lemma}  

\begin{Satz}\label{thm:fallstriktehenselisierung}
  Sei $R$ hier %strikt henselscher,
  diskreter Bewertungsring mit
  perfektem Restklassenkörper $k$
  und $\K$, $E_\K$, $C$, $E$ wie in \ref{thm:exneronmodelle}.
  Erweitert sich die $\K$-Gruppenschemastruktur von $E_\K$ auf
  $E=C^0$, dann ist $E$ Néron-Modell von $E_\K$.
  \begin{proof}
    Da $E$ bereits nach Voraussetzung glattes $R$-Gruppenschema ist,
    ist noch die \NAbbEig zu prüfen.
    Sei also $Y$ ein glattes $R$-Schema mit einem $\K$-Morphismus
    $\phi_\K\colon Y_\K\to E_\K$ der generischen Fasern.
    % , bzw. nach \ref{thm:dichtefaser}
    % einer rationalen Abbildung $\phi\colon Y\dra E$, die von $\phi_\K$
    % präsentiert wird.
    Nach \ref{thm:erwrrational} erweitert sich $\phi_\K$ auf eine
    $R$-rationale Abbildung $\phi\colon Y\dra C$, deren Bild nach dem
    vorhergehenden Lemma~\ref{thm:schnittearithmflglatt} in $E$
    liegt.
    
    Wir werden den Beweis durch Widerspruch führen.
    Angenommen, $\phi\colon Y\dra E$ sei kein Morphismus.
    Da $E$ ein $R$-Gruppenschema, $Y$ glatt und $\phi$ eine
    $R$-rationale Abbildung ist, hat $Y\setminus\Dom(\phi)$
    Kodimension~1 nach \ref{thm:rationalzumorphismus}.
    Das heißt es gibt ein irreduzibles, geschlossenes Unterschema
    $Z\subset Y$ der Kodimension~1, auf dessen generischem Punkt
    $\eta_Z$ die $R$-rationale Abbildung $\phi$ nicht definiert ist.
    % Dies ist äquivalent dazu, dass $\phi$ nicht im generischen Punkt
    % $\eta_Z$ von $Z$ definiert ist, da jede offene Teilmenge $U$ von
    % $Z$ Umgebung von $\eta_Z$ ist.
    $\O_{Y,\eta_Z}$ ist diskreter Bewertungsring, da er von
    Dimension~1, regulär ($Y$ ist glatt, also
    regulär) und lokal ist.
    Sein Quotientenkörper ist der Quotienten-~bzw. Restklassenkörper
    eines generischen Punkts $\eta$ der irreduziblen Komponente von
    $X$, in der $\eta_Z$ liegt. Wir erhalten somit ein kommutatives
    Diagramm
    \begin{center}
      \begin{tikzcd}
        \eta\arrow[r, equal]
        &\Spec(\k(\eta))
        \arrow[r, "{\phi_\K|_{\eta}}"] \arrow[d]
        & E_\K \arrow[r, hook]
        & C\arrow[d]\\
        \eta\cup\eta_Z\arrow[r, equal]
        &\Spec(\O_{Y,\eta_Z}) \arrow[rr]
        \arrow[urr, dashed, "{\exists!}", out=20, in=215]
        && \Spec(R)
      \end{tikzcd}
    \end{center}
    $C$ ist nach Voraussetzung eigentlich, daher erweitert sich
    $\phi_\K|_{\eta}$ nach dem Bewertungskriterium für
    Eigentlichkeit auf einen Morphismus
    $\phi|_{\O_{Y,\eta_Z}}\colon\Spec(\O_{Y,\eta_Z})\to C$.
    \ref{thm:morphismuserweiterung} besagt, dass sich
    dieser auf eine offene Umgebung $U_Z$ von $\eta_Z$ erweitert,
    was bedeutet, dass $\phi|_{U_Z}\colon U_Z\dra E\hookrightarrow C$
    Morphismus ist.
    Nun ist aber $\eta\in E_\K$, also
    $\phi|_{\O_{Y,\eta_Z}}(\eta)\in E$, weshalb nach
    \ref{thm:schnittearithmflglatt}
    $\im(\phi|_{\O_{Y,\eta_Z}})\subset E$.
    Wir können demnach eine glatte, offene Umgebung von
    $\im(\phi|_{\O_{Y,\eta_Z}})$ wählen, deren Urbild $U$ unter
    $\phi|_{U_Z}$ entsprechend wieder offen in $U_Z$, also $Y$, ist.
    Insgesamt ist $\phi|_U\colon U\to E$ ein Morphismus, was
    bedeutet, dass $\phi$ entgegen der Annahme in $\eta_Z$ definiert
    ist.
    Damit muss $\phi\colon Y\to E$ Morphismus sein, sprich die
    \NAbbEig für $E$ ist erfüllt.
    
    % \item Also kann $\phi(\eta_Z)$ nicht in $E$ liegen. Ansonsten
    %   wäre $\phi\colon Y\dra E$ auf dem Urbild einer genügend kleinen
    %   Umgebung von $\eta_Z$ in $E$ Morphismus und entgegen der Annahme in
    %   $\eta_Z$ definiert.
    %   %   C integer, urbilder dichter Mengen dicht
    % \item[?] $\phi(\eta_Z)\not\in E$ ist äquivalent dazu, dass für
    %   jeden geschlossenen (also $k$-wertigen) Punkt $x\in Z\setminus
    %   X_\K$, in dem $\phi$ definiert ist, $\phi(x)\not\in E$ gilt.
    % \item Da $R$ strikt henselsch ist, hat die Einschränkung $Y(R)\to
    %   Y_k(k)$ der $R$-wertigen Punkte auf die spezielle Faser ein
    %   dichtes Bild in $Y_k$ nach \ref{thm:eigstrikthenselsch}.
    % \item[?] Es gibt daher einen $R$-wertigen Punkt $x\in Y(R)$,
    %   dessen Einschränkung $x_0\in Y_k(k)$ einem geschlossen Punkt in
    %   $Z$ entspricht und in dem $\phi\colon Y\to C$ definiert ist.
    % \item Wir erhalten das kommutative Diagramm
    %   \begin{center}
    %     \begin{tikzcd}
    %       \Spec(\K) \arrow[r, "{x|_{\Spec(\K)}}"] \arrow[d]
    %       & Y_\K \arrow[r, "{\phi_\K}"]
    %       & E_\K \arrow[r, hook]
    %       & C \arrow[d]\\
    %       \Spec(R) \arrow[rrr, "\Id"{swap}]
    %       \arrow[urrr, dashed, "{\exists!}", in=210, out=15]
    %       &&&\Spec(R)
    %     \end{tikzcd}
    %   \end{center}
    %   und nach dem Bewertungskriterium für die Eigentlichkeit von $C$
    %   erweitert sich $\phi\circ x|_{\Spec(\K)}$ zu einem Morphismus
    %   $\Spec(R)\to C$, d.\,h.
    %   \begin{gather*}
    %     C(R)\ni\phi\circ x\colon
    %     \Spec(R)
    %     \overset{x}{\longto}
    %     \left\{ x|_{\Spec(\K)}, x_0 \right\}
    %     \overset{\phi}{\longto} C
    %   \end{gather*}
    %   ist Morphismus.
    % \item Es muss aber gelten, dass $\phi\circ x\not\in E(R)$, da
    %   ansonsten entgegen unserer Annahme $\phi$ in $x_0$ definiert
    %   wäre.
    % \item Dies ist ein Widerspruch, da nach
    %   \ref{thm:ratpkteregulaeremodelle} die Gleichheit
    %   $x\not\in E(R)=C(R)\ni x$ gilt.
    % \end{enumerate}
  \end{proof}
\end{Satz}

\section[Gruppenstruktur für strikt henselsche Ringe]%
{Nachweis der Gruppenstruktur für strikt henselsche Bewertungsringe}
\label{chap:gruppestrikthenselsch}
Für den Abschluss des Beweises zu \ref{thm:exneronmodelle} müssen wir
nach \ref{thm:fallstriktehenselisierung} nur noch nachweisen, dass
sich die Gruppenstruktur auf $E$ erweitert.
Dies werden wir in diesem Abschnitt vorerst für strikt henselsche
Bewertungsringe zeigen. Vorweg werden die Definition und Eigenschaften
strikt henselscher Ringe aufgeführt.
% Nur Erweiterung der Gruppenstruktur!
\subsection{Strikt henselsche Ringe}
% HENSELSCHE RINGE
\begin{Definition}[henselscher Ring]\label{def:henselscheringe}
  \optcite[2.3, Proposition 4]{neron}
  \optcite[Chapter IV.6]{silverman2}
  Ein lokaler Ring $(R,\m)$ heißt henselsch, wenn
  er eine der folgenden Bedingungen erfüllt
  \begin{enumerate}[label=(\roman*)]
  \item henselsches Lemma:
    Für jedes normierte Polynom $f\in R[X]$ und jede einfache Nullstelle
    $\bar a\in R/\m$ seiner Restklasse $\bar f\in (R/\m)[X]$
    gibt es eine eindeutige Nullstelle $a\in R$ von $f$, die ein Lift
    von $\bar a$ ist.
  \item Jede endliche $R$-Algebra ist ein Produkt lokaler Ringe.
  \end{enumerate}
  Ein henselscher Ring heißt strikt henselsch, falls sein
  Restklassenkörper $R/\m$ separabel abgeschlossen ist.
  Im Fall eines perfekten Restklassenkörpers ist dies äquivalent zu
  algebraisch abgeschlossen.
\end{Definition}

% \begin{Definition}[Henselisierung]\label{def:henselisierung}
%   \optcite[2.3, Definition 6]{neron}
%   Die Henselisierung eines lokalen Rings $R$ ist ein henselscher
%   lokaler Ring $\Rh$ mit einem lokalen Homomorphismus
%   $i^\text{h}\colon R\to\Rh$, s.\,d. folgende universelle Eigenschaft
%   erfüllt ist:
  
%   Für jeden weiteren lokalen Homomorphismus $u\colon R\to A$ in einen
%   henselschen lokalen Ring $A$ gibt es einen eindeutigen lokalen
%   Homomorphismus $u^\text{h}\colon\Rh\to A$, so dass das folgende Diagramm
%   kommutiert
%   \begin{center}
%     \begin{tikzcd}
%       R \arrow[r, "{u}"]\arrow[d, "{i^\text{h}}"]
%       & A \\
%       \Rh\arrow[ur, dashed, "{\exists!u^\text{h}}"{swap}]
%     \end{tikzcd}
%   \end{center}

%   Die Henselisierung ist nach Definition eindeutig bis auf
%   Isomorphismus und hat denselben Restklassenkörper wie $R$.
%   \optcite[Chapter 2.3, S. 47]{neron}
% \end{Definition}

\begin{Definition}[strikte Henselisierung]\label{def:striktehenselisierung}
  Die strikte Henselisierung eines lokalen Rings $(R,\m)$ mit
  Restklassenkörper $k=R/\m$ ist ein strikt henselscher lokaler Ring
  $\Rsh$, dessen Restklassenkörper der separable Abschluss $k_s$ von
  $k$ ist, zusammen mit einem lokalen Ringhomomorphismus
  $i^\text{sh}\colon R\to\Rsh$, so dass folgende universelle
  Eigenschaft erfüllt ist
  \begin{quote}
    Für jeden weiteren lokalen Ringhomomorphismus $u\colon R\to A$ in
    einen strikt henselschen lokalen Ring $A$ mit Restklassenkörper
    $k_A$ und für jeden $k$"=Körperhomomorphismus $\alpha\colon k_s\to k_A$
    gibt es einen eindeutigen lokalen Ringhomomorphismus
    $u^\text{sh}\colon\Rsh\to A$, so dass das folgende Diagramm
    kommutiert und so dass $\alpha$ die Projektion von $u^\text{sh}$
    auf die Restklassenkörper ist
    \begin{center}
      \begin{tikzcd}
        R \arrow[r, "{u}"]\arrow[d, "{i^\text{sh}}"]
        & A \\
        \Rsh\arrow[ur, dashed, "{\exists!u^\text{sh}}"{swap}]
      \end{tikzcd}
    \end{center}
  \end{quote}
  Die strikte Henselisierung ist nach Definition eindeutig bis auf
  Isomorphismus und hat $k_s$ als Restklassenkörper.
\end{Definition}

\begin{Satz}[{\cite[Proposition IV.6.5]{silverman2}}]
  \label{thm:exhenselisierung}
  Jeder diskrete Bewertungsring $(R,\m)$ hat eine
  % Henselisierung $\Rh$ und eine
  strikte Henselisierung $\Rsh$.
  \begin{proof}
    Siehe
    \cite[Proposition IV.6.5]{silverman2},
    \cite[Remark IV.6.6.2]{silverman2}
    oder \cite[Chapter 2.3, S. 48]{neron}.
  \end{proof}
\end{Satz}

Die für uns besonders wichtige Eigenschaft eines strikt henselschen
Rings ist folgende Aussage über die Einschränkung von Schnitten:
\begin{Satz}[{\cite[Proposition IV.6.4]{silverman2}}]
  \label{thm:eigstrikthenselsch}
  Sei $(R,\m)$ ein diskreter Bewertungsring mit perfektem
  Restklassenkörper $k=R/\m$ und $X$ glattes $R$-Schema mit spezieller
  Faser $X_k$.
  Die Einschränkung
  \begin{gather*}
    X(R)\longto X_k(k)\hookrightarrow X_k
  \end{gather*}
  ist
  \begin{enumerate}[label=(\roman*)]
  \item surjektiv auf $X_k(k)$, falls $R$ henselsch ist.
  \item surjektiv mit dichtem Bild in der speziellen Faser $X_k$,
    falls $R$ strikt henselsch ist.
  \end{enumerate}
\end{Satz}

\subsection{Nachweis der Gruppenstruktur für strikt henselsche Ringe}
Um die Erweiterung der Gruppenverknüpfung auf ganz $E$ für strikt
henselsche Ringe nachzuweisen, werden wir erst zeigen, dass wir eine
normale Komposition auf $E$ erhalten.
Diese ist dann nach \ref{thm:weil} $R$-birational zu dem gesuchten
Gruppenschema und wir werden zeigen, dass diese Birationalität ein
Isomorphismus ist.
% ERWEITERUNG DER GRUPPENSTRUKTUR
\begin{Satz}\label{thm:egruppenschema}
  Sei $R$ strikt henselscher diskreter Bewertungsring mit perfektem
  Restklassenkörper $k$ und $\K$, $E_\K$, $C$, $E$ wieder wie in
  \ref{thm:exneronmodelle}.
  Dann erweitert sich die Gruppenstruktur von $E_\K$ auf $E$ und macht
  $E$ zum Gruppenschema.
  \begin{proof}[Beweis zu \ref{thm:egruppenschema}]
    Die Aussage folgt direkt aus den folgenden beiden Lemmata.
  \end{proof}
\end{Satz}
\begin{Lemma}
  Die Gruppenstruktur $\mu_\K$ auf $E_\K$ erweitert sich zu einer
  normalen Komposition $\mu$ auf $E$.
  % (Hier muss $R$ nicht notwendigerweise strikt henselsch sein.)
  \begin{proof}
    $E\times_R E$ ist glattes $R$-Schema und $C$ ist eigentlich,
    weshalb sich die Gruppenverknüpfung $\mu_\K\colon E_\K\times_\K
    E_\K\to E_\K$ auf $E_\K$ nach \ref{thm:erwrrational} eindeutig zu
    einer $R$-rationalen Abbildung $\tilde\mu\colon E\times_R E\dra C$
    erweitert.
    Die generischen Punkte $\eta_k$ der speziellen Fasern von
    $E\times_R E$ haben wegen \autoref{eq:stalkdimension}
    Dimension~1, sind also diskrete Bewertungsringe, da
    $E\times_R E$ glatt und somit regulär ist. Nach
    \ref{thm:schnittearithmflglatt} liegt dann
    $\tilde\mu(\eta_k)$ und entsprechend das Bild einer offenen
    Umgebung von $\eta_k$ unter $\tilde\mu$ in $E$. Daher erhalten
    wir (unter Umständen durch Einschränkung von $\Dom(\tilde\mu)$)
    eine $R$-rationale Abbildung $\mu\colon E\times_R E\dra E$.

    Für eine normale Komposition bleiben Assoziativität und Inverse zu
    zeigen. Seien hierfür $\phi,\psi\colon E\times_R E\to E\times_R E$
    wie in der Definition \ref{def:normalekomposition}.

    \paragraph{Assoziativität}
    Für Assoziativität müssen für alle $x, z\in E$, für die es
    wohldefiniert ist, die folgenden $R$-rationalen 
    Abbildungen gleich sein
    \begin{gather*}
      \mu(\mu(x,\bullet),z), \mu(x,\mu(\bullet,z))
      \colon E \dra E
    \end{gather*}
    Da $E\times_R E$ reduziert ist, sind aber beide die nach
    \ref{thm:erweindeutig} eindeutige Erweiterung von
    \begin{gather*}
      \mu_\K(\mu_\K(x,\bullet),z)=\mu_\K(x,\mu_\K(\bullet,z))
      \colon E_\K \dra E_\K
    \end{gather*}
    Somit sind sie gleich als $R$-rationale Abbildungen.
    
    \paragraph{Inverses $R$-rational}
    Für $R$-Rationalität von $\phi$ und $\psi$ muss noch gezeigt
    werden, dass $\Dom(\phi)$ bzw. $\Dom(\psi)$ dicht in der
    speziellen Faser von $E$ ist, also in jedem generischen Punkt
    der speziellen Faser definiert ist.
    Wir werden den Fall für $\phi$ behandeln, $\psi$ folgt analog.
    
    Sei also $\eta\in E\times_R (R/\p)=E_k$ ein generischer Punkt der
    speziellen Faser von $E$ und $R'=\O_{E,\eta}$ der Halm in
    $\eta$ mit $\K'\coloneqq\Quot(R)$. % und $k'$ Restklassenkörper.
    Wegen \ref{eq:stalkdimension} hat $R'$
    Dimension 1, ist demnach ein diskreter Bewertungsring und
    $C_{R'}\coloneqq C\times_R R'$ nach \ref{thm:eigminmodell}
    minimales $R'$-Modell der elliptischen Kurve
    $E_\K\times_{\K'}\K'\eqqcolon E_{\K'}$ mit offenem Unterschema
    $E_{R'}=E\times_R R'$ der glatten Punkte.
    Wir wissen, dass per Konstruktion $\Spec(R')\subset E$, also dass
    es einen $R'$-wertigen Punkt $s\colon\Spec(R')\to E_{R'}$ gibt.
    Wir erhalten einen zugehörigen Translationsmorphismus auf der
    generischen Faser, der aufgrund der Gruppenstruktur Automorphismus
    ist:
    \begin{gather*}
      \tau_{\K'}\colon
      E_{\K'}\cong E_{\K'}\times_{\K'}\K'
      \xrightarrow{\Id\times s}
      E_{\K'}\times_{\K'}E_{\K'}
      \xrightarrow{\mu_{\K'}}
      E_{\K'}
    \end{gather*}
    Nach \ref{thm:eigminmodell} erweitert sich $\tau_{\K'}$ zu
    einem $R'$-Automorphismus $\tau_{R'}$ auf dem minimalen Modell
    $C_{R'}$ und auf dem glatten Unterschema $E_{R'}$.
    Für ein $e\in\Spec(R')$, $x=s(e)$ und einen Punkt in der Faser
    $y\in(E_{R'})_e$ sieht $\tau_{R'}$ wie folgt aus
    \begin{center}
      \begin{tikzcd}[row sep=0pt]
        R'\times_R E \arrow[r, equal]
        & R'\times_{R'}(R'\times_R E) \arrow[r, "{s\times\Id}"]
        & (R'\times_R E)\times_{R'}(R'\times_R E)
        \arrow[r, "{\mu_{R'}}", dashed]
        & R'\times_R E \\
        (e,y)\arrow[r, mapsto]
        & (e,(e,y))\arrow[r, mapsto]
        & ((e,x),(e,y))\arrow[r, mapsto]
        & (e,\mu(x,y))
      \end{tikzcd}
    \end{center}
    bzw. $\tilde\tau_{R'}\colon(e,y)\mapsto(e,\mu(y,x))$, wenn man
    $R'\times_{R'}(R'\times_R E)$ durch $(R'\times_R E)\times_{R'}R'$ ersetzt.
    Wir erhalten das kommutative Diagramm
    \begin{center}
      \begin{tikzcd}
        (e,y) \arrow[d, mapsto]\arrow[rrr, bend left=20, mapsto]
        & R'\times_R E \arrow[r, "\tau_{R'}", "\sim"{swap}]
        \arrow[d, "s\times\Id", hook]
        & R'\times_R E\arrow[d, "s\times\Id", hook]
        & (e,\mu(x,y))\arrow[d, mapsto]\\
        (x,y)\arrow[rrr, bend right=20, mapsto]
        & E\times_R E\arrow[r, "\phi", dashed]
        & E\times_R E
        & (x,\mu(x,y))
      \end{tikzcd}
    \end{center}
    und sehen, dass $\phi$ auf dem generischen Punkt
    $(\eta,\eta)\in\Spec(R')\times_R\Spec(R')\subset E\times_R E$
    von $(E\times_R E)_k=E_k\times_k E_k$ definiert ist.
    Mit $\tilde\tau$ statt $\tau$ erhält man dasselbe für $\psi$.
    \paragraph{Inverses $R$-birational}
    Für die Birationalität von $\phi$ bzw. $\psi$ siehe
    \cite[Propsition IV.6.10]{silverman2}.
    % TODO!! Warum birational?
  \end{proof}
\end{Lemma}

\begin{Lemma}\label{thm:äqgruppenschema}
  Sei nach \ref{thm:weil} $G$ das zu $E$ $R$-birationale
  Gruppenschema. Dann ist $G$ isomorph zu $E$.
  \begin{proof}
    Es seien $f\colon G\dra E$ und $g\colon G\dra E$ die
    zueinander inversen $R$-birationalen Abbildungen.
    Es ist zu zeigen, dass beide Morphismen sind.
    \begin{description}
    \item[$g$ Morphismus] Angenommen, $g$ ist kein
      Morphismus. Dann besagt \ref{thm:weil}, dass es eine
      irreduzible, geschlossene Teilmenge $Z\subset E$ gibt, so
      dass $g$ nicht in deren generischen Punkt $\eta_Z$ definiert
      ist. Die spezielle Faser $E_k$ von $E$ ist eindimensional und
      geschlossen, weshalb $\eta_Z\in E_k$ bedeuten würde, dass
      $\eta_Z$ ein generischer Punkt von $E_k$ ist.
      Nachdem $g$ als $R$-rationale Abbildung auf auf einer
      dichten Teilmenge von $E_k$, insbesondere auf dessen
      generischen Punkten, definiert ist, kann $\eta_Z$ also nicht in
      $E_k$ liegen.
      Allerdings ist $g$ auf der gesamten generischen Faser und daher
      dann auch in $\eta_Z$ definiert, was ein Widerspruch ist.
    \item[$f$ Morphismus] Eine Translation $\tau_P^G\colon G\to G$
      auf $G$ um einen Schnitt $P$ liefert eine Translation
      $\tau_P\coloneqq\tau_P^G|_{E_\K}$, also einen
      $\K$-Automorphismus auf $E_\K$ mit Inversem
      $\tau_{(-P)}=\tau_{(-P)}^G$. Diese 
      erweitern sich nach \ref{thm:eigminmodell} zu zueinander
      inversen $R$-Automorphismen
      $\tau_P^E,\tau_{(-P)}^E\colon E\to E$ auf $E$.
      Überall, wo es definiert ist, gilt
      \begin{gather*}
        f = \tau_{(-P)}^E\circ f\circ \tau_P^G
      \end{gather*}
      da beides die eindeutigen Erweiterungen von
      $\Id_{E_\K}=f|_{E_\K}$ sind, wo
      $\tau_P=\tau_P^E|_{E_\K}=\tau_P^G|_{E_\K}$ $\K$-Automorphismus
      ist und $\Id_{E_\K}=\tau_{(-P)}\circ\tau_P$ gilt.
      
      Für jedes $\beta\in G$, so dass es einen Schnitt $P$ gibt mit
      $\tau_P^G(\beta)\in\Dom(f)$, ist $\tau_{(-P)}^E\circ f\circ
      \tau_P^G$ also eine eindeutige Erweiterung von $f|_{E_\K}$
      bzw. von $f$ auf eine Umgebung von $\beta$.
      Wir können also durch Translation den Definitionsbereich von
      $f$ für passende $\beta$ erweitern.

      Jetzt können wir nutzen, dass $R$ nach Voraussetzung strikt
      henselsch ist. Das bedeutet nach
      \ref{thm:eigstrikthenselsch}, dass die Menge der Bildpunkte
      aller Schnitte $P\colon\Spec(R)\to G$ dicht in der
      speziellen Faser liegt.
      Insbesondere findet man in jeder offenen Umgebung eines
      Punkts $\beta\in G$ einen Punkt, um den man mithilfe einer
      Translation den Defintionsbereich von $f$ auf eine Umgebung von
      $\beta$ ausweiten kann, wie oben beschrieben.
      Damit ist $g$ in allen Punkten von $G$ definiert und somit
      Morphismus.
    \end{description}
  \end{proof}
\end{Lemma}


\section{Allgemeine Existenz}\label{chap:exneronmodellebeweis}
Zu allerletzt bleibt noch zu sehen, dass wir uns von Dedekindringen
auf Bewertungsringe und von dort auf strikt henselsche Bewertungsringe
beschränken können.
\begin{proof}[Beweis von \ref{thm:exneronmodelle}]
  Wegen \ref{thm:neronmodelllokal} können wir uns auf
  diskrete Bewertungsringe $R$ beschränken.
  Betrachte also einen diskreten Bewertungsring $R$ mit perfektem
  Restklassenkörper, eine elliptische Kurve $E_\K$ über $\K=\Quot(R)$,
  ein minimales Modell $X$ von $E_\K$ über $R$ und dessen Unterschema
  der glatten Punkte $E$.
  Sei $\Rsh$ die strikte Henselisierung von $R$.
  Nach \cite[Chapter 2.4, Corollary 9]{neron} ist $\Rsh$ treuflach
  über $R$ und unverzweigt per Konstruktion, vergleiche auch
  \cite[Chapter 2.3, Proposition 11]{neron}.
  Daher ist $\Rsh$ insbesondere étale über $R$ und nach
  \ref{thm:eigminmodell} ist $C\times_R\Rsh$ ein minimales Modell von
  $E_\Ksh=E_\K\times_\K\Ksh$ über $\Rsh$ mit Unterschema
  $E\times_R\Rsh$ der glatten Punkte.
  Wir wissen also aus \ref{thm:egruppenschema}, dass $E\times_R\Rsh$
  das Néron-Modell von $E_\Ksh$ über $\Rsh$ ist.

  Für die \NAbbEig von $E$ sei $Y$ glattes $R$-Schema und
  $\phi_\K\colon Y_\K\to E_\K$ Morphismus auf den generischen Fasern.
  $Y\times_R\Rsh$ ist glattes $\Rsh$-Schema, da Glattheit stabil unter
  Basiswechsel ist. Daher erweitert sich der Basiswechsel
  $\phi_\Ksh\colon Y_\Ksh\to E_\Ksh$ von $\phi_\K$ wegen der \NAbbEig
  von $E\times_R\Rsh$ auf einen Morphismus $\phi_\Rsh\colon
  Y\times_R\Rsh\to E\times_R\Rsh$.
  Nach \cite[Chapter 6.1, Theorem 6]{neron} ist wegen Treuflachheit
  von $\pi\colon\Spec(\Rsh)\to\Spec(R)$ der Basiswechselfunktor
  \begin{gather*}
    \text{(Sch/$R$)}\longto\text{(Sch/$\Rsh$)}\;,
    \qquad
    X\longmapsto \pi^*X=X\times_R\Rsh
  \end{gather*}
  volltreu. Deshalb erhalten wir einen
  treuflachen Abstieg $\phi\colon Y\to E$ von $\phi_\Rsh$, der
  $\phi_\K$ erweitert.

  Somit ist $E$ Néron-Modell von $E_\K$.
\end{proof}

%%% Local Variables:
%%% mode: latex
%%% TeX-master: "neron_models"
%%% End:


% Beispiele und Ausblick
\chapter{Beispiele und Ausblick}\label{chap:ausblick}
Inzwischen ist aus \ref{bsp:regweierstraßmodell} und
\ref{bem:konstruktionminmodell} bekannt, dass ein reguläres
Weierstraßmodell über einem Dedekindring bereits ein minimales Modell
seiner generischen Faser ist.
In einigen Fällen kann die Regularität an den
Koeffizienten der Weierstraßgleichung abgelesen werden.
Sei im Folgenden $(R,(\pi))$ ein diskreter Bewertungsring mit
perfektem Restklassenkörper, ${\K=\Quot(R)}$, $E_\K$ ein elliptische
Kurve über $\K$ und $W$ ein minimales Weierstraßmodell von $E_\K$ zur
Gleichung ($\Char(\K)\neq2,3$):
\begin{gather}\label{eq:allgweierstraßgleichung}
  f\colon y^2=x^3 + \beta x + \gamma\;.
\end{gather}
% \begin{Beispiel}[$b_n$"~Reduktion: $\ord_\K(\Delta)=1$]
%   \optcite[Lemma IV.9.5]{silverman2}
%   Ist $\ord_\K(\Delta)=1$, sind nach
%   \cite[Chapter 1.5, Lemma 3 and Lemma 4]{neron} die Koeffizienten
%   $\beta$ und $\gamma$ der Gleichung Einheiten in $R$. Eine Rechnung
%   wie in \cite[Lemma IV.9.5]{silverman2} zeigt, dass $W$ dann regulär
%   im singulären Punkt $(x,y,\pi)\subset R[x,y]/(f)$ ist.
% \end{Beispiel}

\begin{Korollar}\label{thm:c1red}
  % $\ord_\K(j)\geq0$]
  \label{bem:bedregweierstraßmodell}
  Eine hinreichende Bedingung dafür, dass $W$ regulär ist, ist
  \begin{gather*}
    \ord_\K(\beta)\geq 1
    \qquad\text{und}\qquad
    \ord_\K(\gamma)=1\;.
  \end{gather*}
  Das heißt $\gamma=u\pi$ und $\beta=v\pi^n$ für Einheiten
  $u,{v\in R^\times}$ und $n\in\N_{\geq1}$.

  \begin{proof}
    Wie aus \ref{thm:weierstrassglatt} bekannt, ist nur die
    Regularität des singulären Punkts der speziellen Faser zu
    überprüfen.
    Sei $\m_\pi\coloneqq(\pi)$ und $k=(R_{\m_\pi}/\m_\pi)$ der
    Restklassenkörper von $R$.
    Die spezielle Faser $W_k$ ist auf dem betrachteten Teil seiner
    affinen Überdeckung das Spektrum von 
    \begin{gather*}
      \dfrac{k[\bar x,\bar y]}{(\bar y^2-\bar x^3+\beta \bar x+\gamma)}
      =\dfrac{k[\bar x,\bar y]}{(\bar y^2-\bar x^3+v\pi^n \bar x+u\pi x)}
      =\dfrac{k[\bar x,\bar y]}{(\bar y^2-\bar x^3)}\;,
    \end{gather*}
    welches im abgeschlossenen Punkt $(0,0)\in\A_k^2$ beziehungsweise
    $(\bar x,\bar y)\in W_k$ singulär ist. Da der singuläre Punkt von
    $W_k$ eindeutig ist, muss nur dieser überprüft werden.
    Er entspricht in $W$ dem Maximalideal
    $\p\coloneqq(\pi,x,y)\subset R[x,y]/(f)$.
    Für Regularität gilt zu zeigen, dass $\p$ in seinem Halm
    $(R[x,y]/(f))_\p=\O_{W,\p}$ nur von zwei Elementen erzeugt wird.
    Durch Umstellen von \eqref{eq:allgweierstraßgleichung} ist aber
    ersichtlich, dass
    \begin{gather*}
      \pi = u^{-1}\left( y^2-x^3-\beta x \right)\in\p\;,
    \end{gather*}
    weshalb $\p=(x,y)$ tatsächlich nur zwei Erzeuger hat.
  \end{proof}
\end{Korollar}

\begin{Bemerkung}
Nach \cite[Lemma~1.5/3 und~1.5/4]{neron} ist die
Bedingung $\ord_\K(\beta)\geq 1$ und $\ord_\K(\gamma)=1$ für
minimale Weierstraßgleichungen äquivalent dazu,
dass $\ord_\K(\Delta)=2$ und $\ord_\K(j)\geq 1$,
wobei $j=2^8\cdot 3^3\cdot\beta^3$ die $j$"~Invariante der
Weierstraßgleichung ist.
\end{Bemerkung}

Wie in \cite[Chapter~1.5]{neron} beschrieben, ist die
Gruppenstruktur der speziellen Faser im Fall von \ref{thm:c1red} die
des additiven Gruppenschemas $\G_{a,k}$ über $k$. Dieses ist definiert
als $\G_{a}\times_\Z k$, wobei $\G_{a}\coloneqq\Spec(\Z[T])$ ist mit
der Gruppenstruktur
\begin{center}
  \begin{tikzcd}[row sep=5pt]
    \G_{a}\times_\Z \G_{a} \arrow[r,"\mu"] \arrow[d, equal]
    & \G_{a}\arrow[ddd, equal]\\
    \left(\Spec\Z[T_1]\right) \times_\Z
    \left(\Spec\Z[T_2]\right) \arrow[d, equal] \\
    \Spec\left(\Z[T_1]\otimes_\Z\Z[T_2]\right) \arrow[d,equal]\\
    \Spec\Z[T_1,T_2] \arrow[r, "\mu"]
    &\Spec(\Z[T])\\
    T_1+T_2\arrow[from=r, mapsto]
    & T
  \end{tikzcd}
\end{center}
    
\begin{Beispiel}
  Sei $R=\Z_{(5)}$ und betrachte die Weierstraßgleichungen
  \begin{align*}
    y^2&=x^3+5x\pm 10
    & \Delta&=-16\cdot3200=-16\cdot2^7\cdot5^2\\
    y^2&=x^3\pm 5
    & \Delta&=-16\cdot675=-16\cdot3^3\cdot5^2\\
    y^2&=x^3\pm 10
    & \Delta&=-16\cdot2700=-16\cdot2^2\cdot3^3\cdot5^2
  \end{align*}
  Für alle Gleichungen ist offensichtlich die Bedingung aus
  \ref{bem:bedregweierstraßmodell} erfüllt.
  Somit sind die projektiven Abschlüsse der elliptischen Kurven über
  $\Z_{(5)}$ alle bereits minimale Modelle und ihr glatter Teil,
  das heißt der Punkt $(x,y,5)$ ausgenommen, ist das jeweilige
  Néron-Modell.
\end{Beispiel}

\begin{Beispiel}
  Ein Anwendungsbeispiel für die Aussage aus
  \ref{thm:teilwabelscheneronmodelle} ist jetzt besonders ersichtlich:
  Bei einem regulären Weierstraßmodell $W$ ist die induzierte
  Gruppenstruktur auf dem glatten Teil in einer Umgebung der glatten
  Fasern von $W$ abelsch.
  Dies liegt daran, dass die glatten Fasern von $W$ bereits
  elliptische Kurven sind.
\end{Beispiel}

% TODO: Tate's alg
Wie zuvor bereits angedeutet, ist es möglich, die minimalen Modelle
elliptischer Kurven anhand ihrer Weierstraßgleichung nach der Form
ihrer speziellen Faser in zehn verschiedenen Kategorien zu
klassifizieren \cite[vgl.][Theorem~IV.8.2]{silverman2}.
Zur Klassifizierung kann ein Algorithmus von Tate verwendet werden,
der von einer Weierstraßgleichung der elliptischen Kurve ausgeht, diese
minimalisiert und dann einer der Kategorien in
\autoref{tab:faserklassifizierung} zuordnet.
\begin{table}[tbhp]
  \begin{center}
    \begin{tabularx}{\linewidth}{@{}l>{\hsize=.7\hsize}X>{\hsize=1.3\hsize}X@{}}
      \toprule
       & Bedingungen an \newline die Koeffizienten
      & Beschreibung der speziellen Faser\\
      \midrule[\heavyrulewidth]
      $a$& $\ord_\K(j)\geq0$, $\ord_\K(\Delta)=0$
      & elliptische Kurve\\\midrule
      $b_n$& ${\ord_\K(j)<0}$, \newline
      ${\ord_\K(\beta)=0}$, ${\ord_\K(\gamma)=0}$
      & $n=1$: Kurve mit Knoten\newline
        $n>1$: $n$ reguläre Kurven, die sich schneiden \\\midrule
      $c1$& $\ord_\K(j)\geq0$, $\ord_\K(\Delta)=2$
      & singuläre Kurve mit Scheitelpunkt\\\midrule
      $c2$& $\ord_\K(j)\geq0$, $\ord_\K(\Delta)=3$
      & 2 reguläre Kurven mit 1 tangentialem Schnittpunkt \\\midrule
      $c3$& $\ord_\K(j)\geq0$, $\ord_\K(\Delta)=4$
      & 3 reguläre Kurven mit 1 gem. Schnittpunkt \\\midrule
      $c4$& $\ord_\K(j)\geq0$, $\ord_\K(\Delta)=6$
      & 1 reguläre Kurve mit Multiplizität 2, die 4 weitere reguläre
      Kurven schneidet\\\midrule
      $c5_n$& ${\ord_\K(j)<0}$, \newline
      ${\ord_\K(\beta)=2}$, ${\ord_\K(\gamma)=3}$
      & $n$ reguläre Kurven mit Multiplizität 2 bilden eine Kette, die
        an den 2 Enden jeweils eine Kurve der Multiplizität 1 schneidet\\\midrule
      $c6$& $\ord_\K(j)\geq0$, $\ord_\K(\Delta)=8$
      & 1 reguläre Kurve der Multiplizität 3 schneidet 3 Kurven mit
        Multiplizität 2, die jeweils 1 mit Multiplizität 1 schneiden\\\midrule
      $c7$& $\ord_\K(j)\geq0$, $\ord_\K(\Delta)=9$
      & ähnlich $c6$ mit 8 Kurven und bis Multiplizität 4 \\\midrule
      $c8$& $\ord_\K(j)\geq0$, $\ord_\K(\Delta)=10$
      & ähnlich $c6$ mit 9 Kurven und bis Multiplizität 6\\\bottomrule
    \end{tabularx}
  \end{center}
  \caption{Kategorien für die
    Klassifizierung der speziellen Faser des minimalen Modells einer
    elliptischen Kurve über einem diskreten Bewertungsring mit
    perfektem Restklassenkörper,
    nach Néron,
    vergleiche \cite[Chapter~1.5]{neron} und
    \cite[Figure~IV.4.4]{silverman2}
    \label{tab:faserklassifizierung}} 
\end{table}
Für genaueres zu Tates Algorithmus siehe
\cite[Chapter~IV.9]{silverman2},
\cite[Chapter~1.5]{neron} und \cite{tate}.
Ein analoger Algorithmus, welcher auch den Fall von Dedekindringen mit
nicht perfekten Restklassenkörpern abdeckt, ist in \cite{nonperfect}
zu finden.


%%% Local Variables:
%%% mode: latex
%%% TeX-master: "neron_models"
%%% End:


% Appendix (Notizen)
% \appendix
\chapter{Verwendete Definitionen}
\section{Dedekindringe}
% DEDEKINDRINGE
\begin{Definition}[Dedekindring]\label{def:dedekindring}
  \optcite[S.\,40]{hartshorne}\optcite[Definition B.84]{wedhorn}
  Ein Integritätsring $R$ heißt Dedekindring, wenn er eine der
  äquivalenten Bedingungen erfüllt
  \begin{enumerate}[label=(\roman*)]
  \item regulär und $\dim(A)\leq 1$
  \item noethersch, normal und $\dim(A)\leq 1$
  \item $A$ Körper oder
    jede Lokalisierung nach einem Maximalideal ist ein diskreter
    Bewertungsring
    (d.\,h. lokal, noethersch, normal von Dimension 1)
  \item jedes echte Ideal ist endliches Produkt von Primidealen
    (diese ist dann eindeutig bis auf Reihenfolge) 
  \end{enumerate}
\end{Definition}
Dedekindringe sind genau die Hauptidealringe.
  \optcite[Proposition B.85]{wedhorn}
  
  \begin{Definition}[Dedekindschema]
    Ein Dedekindschema ist ein noethersches, integeres Schema $X$, das
    eine der äquivalenten Eigenschaften erfüllt
    \begin{enumerate}[label=(\roman*)]
    \item die offenen affinen Unterschemata sind Spektren von Dedekindringen
    \item $\dim(X)\leq 1$ und $X$ regulär
      (in Dimension 1 äquivalent zu normal nach
      \cite[Corollary 6.39, Proposition 6.40]{wedhorn})
    \end{enumerate}
    Ein integeres Schema mit endlicher Überdeckung durch Dedekindschemata
    ist selbst Dedekindschema.
    \optcite[Chapter (7.13)]{wedhorn}
  \end{Definition}

\begin{Bemerkung}
  Die Bedingung, dass $S$ integer sein muss, liefert einen eindeutigen
  generischen Punkt $\eta$ und einen Funktionenkörper
  $\K=\k(\eta)=\O_{S,\eta}$, so dass jeder Schnitt und jeder Halm ein
  Unterring von $\K$ ist, d.\,h. es gibt für jeden Punkt $x\in X$ einen
  Morphismus $\Spec(\K)\to\Spec(\k(x))$ und für jede offene Teilmenge
  $U\subset X$ einen Morphismus $\Spec(\K)\to\Spec(\Gamma(U,\O_X))$.
  
  Verzichtet man auf diese Bedingung, erhält man eine Familie
  generischer Punkte $\eta_i$ der endlich vielen irreduziblen
  Komponenten (beachte: $S$ noethersch). Dieser Fall kann analog
  behandelt werden, indem man $\K=\bigoplus_i \k(\eta_i)$ setzt.
\end{Bemerkung}

\begin{Bemerkung}[generische und spezielle Fasern]
  Ist $S=\Spec(R)$ mit $R$ Dedekindring, so hat ein $S$-Schema die
  \emph{generische Faser} über dem generische Punkt $\eta=(0)$ und die
  \emph{speziellen Fasern} über den verbleibenden Maximalidealen
  (geschlossenen Punkten).
\end{Bemerkung}


\section{Eigenschaften von Schemata}
% REDUZIERT
\begin{Definition}[reduziert, geometrisch reduziert]
  Ein Schema heißt reduziert, falls alle Halme reduzierte Ringe sind
  (d.\,h. $\Nil(0)=(0)$).

  Ein $\K$-Schema heißt geometrisch reduziert, wenn für jeden
  Basiswechsel zu einer Körpererweiterung von $\K$ die generische
  Faser reduziert ist. 
\end{Definition}

% LOKAL VON ENDLICHEM TYP
\begin{Definition}[lokal von endlichem Typ]
  Ein Morphismus $f\colon X\to Y$ von Schemata heißt lokal von endlichem
  Typ, wenn lokal auf affinen offenen Teilmengen $V\subset Y$ die
  induzierte Abbildung $f_U^\#: f_\ast\O_X(V)\to \O_Y(V)$
  (wobei $U=f^{-1}(V)$) $f_\ast\O_X(V)=\O_X(U)$ zu einer
  \emph{endlich erzeugten} $\O_Y(V)$-Algebra macht.

  Er heißt von endlichem Typ, wenn er lokal von endlichem Typ und
  quasikompakt (Urbilder quasikompakter Mengen sind quasikompakt) ist.
\end{Definition}

% REGULÄR
\begin{Definition}[Regulär]
  Ein Schema $X$ heißt regulär im Punk $x\in X$, wenn sein Halm
  $\O_{X,x}$ ein regulärer (lokaler) Ring ist
  (d.\,h. $\dim_{\O_{X,x}/\m_x}(\m_x/\m_x^2)=\dim(\O_{X,x})$).

  Regularität kann auf lokal noetherschen Schemata auf den geschlossenen
  Punkten geprüft werden:
  Ist $X$ lokal noethersch und regulär in allen geschlossenen Punkten,
  so ist $X$ regulär nach \cite[Remark 6.25 (3)]{wedhorn}.
\end{Definition}

\section{Eigenschaften von Morphismen von Schemata}
% FLACH
\begin{Definition}[flacher Morphismus]
  \optcite{bosch, goodreduction}
  Ein Morphismus $f\colon X\to Y$ von Schemata ist flach, wenn alle
  induzierten Ringhomomorphismen $f_x^\#\colon \O_{Y,f(x)}\to\O_{X,x}$  auf
  den Halmen flach sind.
\end{Definition}
\begin{Bemerkung}\label{thm:flachgenerischefaser}
  Für einen Dedekindring $R$ und ein reduziertes $R$-Schema $f\colon
  X\to R$ ist flach äquivalent dazu, dass alle generischen Punkte von 
  $X$ auf den generischen Punkt von $R$ geschickt werden
  nach \cite[Proposition III.9.7]{hartshorne}.
  Insbesondere ist dann die generische Faser nichtleer und dicht in $X$.
\end{Bemerkung}

% UNVERZWEIGT
% ??? Unverzweigt <=> Glatt für diskr. BWR?
\begin{Definition}[unverzweigter Morphismus]
  Ein Morphismus $f\colon X\to S$ von Schemata heißt unverzweigt im
  Punkt $x\in X$, falls es eine offene Umgebung
  $U\subset X$ von $x$ gibt und eine geschlossene $S$-Immersion
  $j\colon U\cong\O_W/\I\to W\subset\A_S^n$ in ein offenes Unterschema
  des affinen $n$-Raums mit zugehöriger Idealgarbe $\I\subset\O_W$,
  für die gilt
  \begin{enumerate}[label=(\alph*)]
  \item $\I$ ist in einer Umgebung von $z\coloneqq j(x)$ endlich erzeugt.
  \item $\Om{\A_S^n/S}$ ist erzeugt von den $\d g\in\Om{\A_S^n/S}$
    mit $g\in\I$.
  \end{enumerate}
  $f$ heißt unverzweigt, falls es in jedem Punkt $x\in X$
  unverzweigt ist.
  \begin{Bemerkung}
    \optcite[8.4, Theorem 3]{bosch}
    Ist $X$ lokal endlich präsentiert, sind äquivalent
    \begin{enumerate}[label=(\roman*)]
    \item $X$ unverzweigt in $x$,
    \item $\Om{X/S,x}=0$,
    \item $X_{f(x)}\to \Spec(\k(f(x)))$ unverzweigt.
    \end{enumerate}
  \end{Bemerkung}
\end{Definition}

% GLATT
\begin{Definition}[glatter Morphismus]
  \optcite[8.5, Definition 1]{bosch}
  Ein Morphismus $f\colon X\to S$ von Schemata heißt glatt im Punkt
  $x\in X$ von relativer Dimension $r$, wenn es eine offene Umgebung
  $U\subset X$ von $x$ gibt und eine geschlossene Immersion $j\colon
  U\cong\O_W/\I\to W\subset\A_S^n$ in ein offenes Unterschema des affinen
  $n$-Raums mit zugehöriger Idealgarbe $\I\subset\O_W$, für die
  gilt
  \begin{enumerate}[label=(\alph*)]
  \item $\I$ wird in einer Umgebung von $z\coloneqq j(x)$ von $n-r$ Schnitten
    $g_{r+1},\dotsc,g_{n}$ erzeugt.
  \item Die $\d g_{r+1}(z),\dotsc,\d g_{n}(z)\in 
    \Om{\A_S^n/S,z} \otimes_{\O_{\A_S^n,z}} \k(z)
    \cong \Om{\A_S^n/S,z}/\m_z\Om{\A_S^n/S,z}$ 
    sind linear unabhängig über $\k(z)$.
    Dies ist äquivalent zu
    \begin{gather*}\tag{Jacobi Bedingung}
      \rk\left(
        \frac{\partial g_j}{\partial t_i}\d t_i 
      \right)_{\substack{j=r+1,\dotsc,n\\i=1,\dotsc,n\phantom{+1}}}
      = n-r
    \end{gather*}
  \end{enumerate}
  Ist $f$ glatt in jedem Punkt heißt $f$ glatt bzw. $X$ ist glattes
  $S$-Schema.  
\end{Definition}
\begin{Bemerkung}\label{thm:eigglatt}
  Für $X$ lokal endlich präsentiert ist $f$ glatt in $x$ äquivalent zu
  $f$ flach in $x$ und $X_{f(x)}$ glatt über $\k(f(x))$ nach
  \cite[8.5, Proposition 17]{bosch}.
  
  Ist $X$ ein $\K$-Schema lokal endlichen Typs sind glatt und
  geometrisch regulär äquivalent nach \cite[Corollary 6.32]{wedhorn}.
  Ist $\K$ perfekter Körper (z.\,B. $\Char(\K)=0$), so ist dies sogar
  äquivalent zu regulär nach \cite[Remark 6.33]{wedhorn}.
    
  Glatt impliziert reduziert, lokal von endlichem Typ und flach
  \cite[für flach s.][Theorem 14.22]{wedhorn}.
\end{Bemerkung}


% ÉTALE
\begin{Definition}[étaler Morphismus]% Äquivalent zu lokal umkehrbar
  \optcite[8.5, Definition 1 und Corollary 12]{bosch}  
  Ein Morphismus $f\colon X\to Y$ von $S$-Schemata heißt étale in
  einem Punkt $x\in X$, falls er eine der äquivalenten Bedingungen
  erfüllt 
  \begin{enumerate}[label=(\roman*)]
  \item $f$ glatt von relativer Dimension 0 in $x$
  \item $X$ glatt in $x$, $Y$ glatt in $f(x)$ und 
    $(f^*\Om{Y/S})_x \overset{\sim}{\longto} \Om{X/S,x}$
  \end{enumerate}
  Er heißt étale, falls er étale in jedem Punkt von $X$ ist.
  
  $f$ ist genau dann étale, wenn $f$ glatt und unverzweigt ist.
  \optcite[8.5, Proposition 6]{bosch}

  Offene Immersionen sind genau die étalen Immersionen nach
  \cite[8.5, Lemma 7]{bosch}.
\end{Definition}

SEPARIERT
\begin{Definition}[separierter Morphismus]
  \cite[Theorem 15.8]{wedhorn}
  
  Unterschemata separierter Schemata sind wieder separiert:
  Sei $S$ Schema, $X$ separiertes $S$-Schema,
  $W\overset{\iota}\hookrightarrow X$ Unterschema.
  Dann ist
  \begin{center}
    \begin{tikzcd}
      W\arrow[r, hook]\arrow[d, "{\Delta}"]
      & X\arrow[d, "{\Delta}"]\\
      W\times_R W \arrow[r, hook] X\times_R X
    \end{tikzcd}
  \end{center}
  kartesisch, denn es kommutiert und für jedes weitere $V$, so dass
  \begin{center}
    \begin{tikzcd}
      V\arrow[r, "\phi"]\arrow[d, "{(\phi_1,\phi_2)}"]
      & X\arrow[d, "{\Delta}"]\\
      W\times_R W \arrow[r, hook] X\times_R X
    \end{tikzcd}
  \end{center}
  ebenfalls kommutiert bzw. $\iota\circ\phi_1=\phi=\iota\circ\phi_2$,
  muss wegen der Injektivität von $\iota$ gelten $\phi_1=\phi_2$ und
  $\phi_1$ lässt 
  
  http://math.stackexchange.com/questions/1088879/closed-immersion-of-separated-scheme-is-separated
  \end{Definition}

% EIGENTLICH
\begin{Definition}[eigentlicher Morphismus]
  % Verallgemeinerung von projektiv; Äquivalent zu kompakt
  Ein Morphismus $f\colon X\to Y$ heißt eigentlich, wenn er eine der
  folgenden äquivalenten Bedingungen erfüllt
  \begin{enumerate}[label=(\roman*)]
  \item $f$ ist separiert, von endlichem Typ und universell
    geschlossen (d.\,h. jeder Basiswechsel ist geschlossen)
  \item $f$ ist quasi-separiert, von endlichem Typ und für jeden
    Bewertungsring $R$ gilt
    \begin{gather*}
      X(R)=\Hom_Y(R, X) 
      \overset{\sim}{\underset{\circ\Spec(\operatorname{incl})}
        {\longrightarrow}}
      \Hom_Y(\Quot(R), X)=X(\Quot(R))\;.
    \end{gather*}
    In anderen Worten, es findet sich für jedes kommutative Diagramm
    wie folgt ein entsprechender, eindeutiger Lift
    \begin{center}
      \begin{tikzcd}
        \Spec(\Quot(R)) \arrow[r]\arrow[d]
        & X \arrow[d, "f"] \\
        \Spec(R) \arrow[r]\arrow[to=ur, dashed, "\exists!"]
        & Y
      \end{tikzcd}
    \end{center}
  \end{enumerate}
\end{Definition}
\begin{Bemerkung}
  Endliche sowie projektive Morphismen sind eigentlich
  \cite[9.5, Remark 5 und Theorem 9]{bosch}.
\end{Bemerkung}
\begin{Bemerkung}[Stein Faktorisierung]
  \cite[9.5, Theorem 12]{bosch}
  Jeder eigentliche Morphismus $f\colon X\to Y$ faktorisiert über
  einen eigentlichen, surjektiven Morphismus $g\colon X\to
  \Spec(f_*\O_X)$ mit zusammenhängenden Fasern und einen endlichen
  Morphismus $h\colon \Spec(f_*\O_X)\to Y$.
\end{Bemerkung}


\section{$\K$-Gruppenschemata}
% GUTE REDUKTION
\begin{Definition}[Gute Reduktion]\label{def:gutereduktion}
  % $X/S$ hat Gute Reduktion nach dem Primideal $P\in S$ heißt, dass die
  % Faser $X_P$ Glattheit erhält.
  Sei $S$ ein Dedekindring, $\K=\Quot(R)$ und $X/\K$
  ein eigentliches glattes $\K$-Schema.
  $X$ hat \emph{gute Reduktion im geschlossenen Punkt $s\in\Spec(R)$},
  falls es ein eigentliches glattes $R_s$-Schema $Y_s$ gibt, so dass
  $X=Y_s\times_{R_s}\Spec(\K)$ ist. D.\,h. folgendes Diagramm ist
  kartesisch
  \begin{center}
    \begin{tikzcd}
      X \arrow[r]\arrow[d]
      \arrow[dr, phantom, "\ulcorner", very near start]
      & Y_s \arrow[d] \\
      \Spec(\K) \arrow[r]
      & \Spec(R_s)
    \end{tikzcd}
  \end{center}
  \cite[vgl.][Chapter 1.4]{neron}

  Wir sagen $X$ hat \emph{gute Reduktion}, falls es ein eigentliches
  glattes $R$-schema $Y$ gibt, so dass $X$ isomorph zur generischen
  Faser $Y_{\K}$ von $Y$ ist. D.\,h. das folgende Diagramm ist
  kartesisch
  \begin{center}
    \begin{tikzcd}
      X \arrow[r]\arrow[d]
      \arrow[dr, phantom, "\ulcorner", very near start]
      & Y \arrow[d] \\
      \Spec(\K) \arrow[r]
      & \Spec(R)
    \end{tikzcd}
  \end{center}
  \cite[Definition 3.1]{goodreduction}, \cite[vgl.][Chapter 1.4]{neron}

  Alternative Definition für elliptische Kurven: \cite[Chapter VII.5]{silverman}
\end{Definition}

VARIETÄT, ELLIPTISCHE KURVE
\begin{Definition}[Varietät]
 \cite[Theorem 3.37]{wedhorn}
  Eine affine Varietät über einem Körper $\K$ ist eine Unterprägarbe der
  Prägarbe $U\mapsto \Hom(U,\K)$ über den geschlossenen Punkten eines
  geschlossenen, irreduziblen Unterschemas des $\A_\K^n$.
  Eine affine Prävarietät ist eine quasikompakte, zusammenhängende
  Prägarbe, die lokal affine Varietät ist.

  Für integere $\K$-Schemata endlichen Typs sind die geschlossenen
  Punkte isomorph zu den $\K$-Schnitten 
  $\Hom_{\Spec(\K)}(\Spec(\K),X)$.
  Dadurch gibt es eine Kategorienäquivalenz
  \begin{align*}
    \left\{\text{Integere $\K$-Schemata endlichen Typs}\right\}
    &\longto
      \left\{\text{Prävarietäten über $\K$}\right\}\\
    (X,\O_X) 
    &\longmapsto
      \left( X(\K), \O_{X(\K)}\subset \Hom(\bullet, \K) \right)
  \end{align*}
  wobei $f\in\O_{X(\K)}(U)$ der Abbildung 
  \begin{align*}
    U\cap X(\K) &\longto \K\\
    x  &\longmapsto f(x)
         \coloneqq \bar{f_x}\in\O_{X,x}/\m_x\cong \K
  \end{align*}
  entspricht (Analogon zum Einsetzhomomorphismus für Polynomringe).
  Der inverse Funktor ist gegeben durch
  \begin{align*}
    (P,\O_P) \longmapsto (t(P), t_\ast\O_P) 
  \end{align*}
  wobei $t\colon P\to t(P)$, $x\mapsto \{x\}$. Auf affinen Varietäten
  $V$ wird dies zu
  \begin{gather*}
    (V, \O_V) \mapsto (\Spec(\O_V(V)), \O_{\Spec(\O_V(V))})
  \end{gather*}

  Wir werden im folgenden mit Prävarietät ein integeres
  $\K$-Schema endlichen Typs meinen und mit affiner Varietät ein
  affines, integeres $\K$-Schema endlichen Typs.

  Eine algebraische Varietät sei im folgenden ein separiertes, integeres
  Schema über einem algebraisch abgeschlossen Körper $\K$
  (alternativ: über einem Körper $\K$, wobei die Basiserweiterung auf
  den algebraischen Abschluss weiterhin integer ist).
\end{Definition}
\begin{Definition}[Kurve, elliptische Kurve]
  Eine Kurve über einem Körper $\K$ ist eine projektive Varietät der
  Dimension 1,
  also ein projektives, integeres $\K$-Schema endlichen Typs.

  Eine elliptische Kurve über $\K$ ist eine reguläre Kurve $E$ vom
  Genus 1 zusammen mit einem $\K$-Schnitt $O\in E(\K)$.
  Elliptische Kurven sind glatte \cite[Proposition III.3.1]{silverman}
  $\K$"=Gruppenschemata mit Nullelement $O$ \cite{silverman}.
\end{Definition}

%%% Local Variables:
%%% mode: latex
%%% TeX-master: "neron_models"
%%% End:


% Literatur
%\nocite{*}
\printbibliography

% Erklärung
\chapter*{Eigenständigkeitserklärung}
\thispagestyle{empty}
Hiermit erkläre ich, Gesina Schwalbe, geboren am 29.01.1996 in Aachen,
dass ich die vorliegende Arbeit selbstständig verfasst und keine
anderen als die angegebenen Hilfsmittel und Quellen verwendet
habe. Weiter versichere ich, dass ich die Arbeit noch keiner anderen
Prüfungsbehörde vorgelegt habe.

\vspace*{5em}
Regensburg, den \abgabedatum
\hfill
Gesina Schwalbe


%%% Local Variables:
%%% mode: latex
%%% TeX-master: "neron_models"
%%% End:


\end{document}
